% =============================================================================
% Poisson inference and Monte Carlo methods for inference --- book exercises (verbatim statements),
% original math and conceptual questions, every one with a programming part; full solutions.
% Built from the MPE exam template (lista-exam.tex). One source, two outputs:
%   \solucoesfalse  -> blank list (space to work)
%   \solucoestrue   -> answer key (same file)
% Compile: pdflatex (2x); check: pdffonts <file>.pdf (Type 1 only).
% Every number printed in the solutions is reproduced and asserted by
%   python Listas/poisson-monte-carlo-check.py
% Scaffold notebook for the programming parts: Listas/lista-poisson-monte-carlo.ipynb
% =============================================================================
\documentclass[11pt,a4paper]{article}
\usepackage[utf8]{inputenc}
\usepackage[T1]{fontenc}
\usepackage{lmodern}   % Type 1 vector fonts (replaces PK bitmaps)
\usepackage{cmap}      % ToUnicode CMap -> copyable text and accents
\usepackage[brazil,english]{babel}
\usepackage{amsmath,amssymb,amsfonts}
\usepackage[a4paper,margin=2.2cm]{geometry}
\usepackage{enumitem}
\usepackage{fancyhdr}
\usepackage{comment}
% One cut file per job: parallel builds in one folder would otherwise share comment.cut and
% swap solutions. Redefine DefaultCutFileName; overriding CommentCutFile alone is reset by comment.sty.
\def\DefaultCutFileName{\def\CommentCutFile{\jobname.cut}}
\usepackage{needspace}
\usepackage{xcolor}
\usepackage{microtype}
\usepackage{booktabs}
\usepackage{pgfplots}
\pgfplotsset{compat=1.16}

% ====================== CONFIGURATION ======================
\newif\ifsolucoes
\solucoesfalse                 % <<<< \solucoestrue for the ANSWER KEY; \solucoesfalse for the LIST

\newcommand{\disciplina}{Econometria Avançada --- MPE (Insper)}
\newcommand{\listatema}{\textit{Poisson inference and Monte Carlo methods for inference}}
\newcommand{\espacoresol}{4.5cm}   % height of the working space in the blank list

% ====================== MACHINERY ======================
\ifsolucoes\includecomment{sol}\else\excludecomment{sol}\fi

\pagestyle{fancy}\fancyhf{}
\lhead{\footnotesize\disciplina}
\rhead{\footnotesize Poisson \& Monte Carlo\ifsolucoes{} --- Solutions\fi}
\cfoot{\footnotesize\thepage}
\renewcommand{\headrulewidth}{0.3pt}
\setlength{\headheight}{14pt}

\setlist{topsep=2pt,itemsep=2pt}
\newcounter{questao}
\numberwithin{equation}{questao}   % equation (q.m): question q, equation m

% Block header (one block per topic)
\newcommand{\bloco}[3]{\bigskip\noindent
  {\large\bfseries\color{blue!55!black}Block #1 --- #2}\hfill\textit{(#3)}\par
  \noindent\rule{\linewidth}{0.3pt}\smallskip}

% Question (automatic numbering)
\newcommand{\questao}[1]{\refstepcounter{questao}\medskip\needspace{5\baselineskip}\noindent
  \textbf{Question \thequestao.}~#1\par}
% Source line that opens every question
\newcommand{\fonte}[1]{\par\noindent{\small\textbf{Source:} #1}\par\smallskip}
% Programming part added to every question
\newcommand{\programacao}{\par\smallskip\noindent\textbf{Programming (added, not in the book).}~}
% Multiple-choice options and key
\newenvironment{alternativas}{\begin{enumerate}[label=\Alph*.,leftmargin=2.2em,topsep=1pt]}{\end{enumerate}}
\newcommand{\gabarito}[1]{\ifsolucoes\par\noindent\textbf{Key:}~#1\fi}
% Notation list that opens every solution (estilo-avaliacao §4, "Equations, notation and layout").
\newenvironment{notacao}{\par\smallskip\noindent\textbf{Notation.}\begin{itemize}[leftmargin=1.4em,topsep=1pt,itemsep=0pt]\small}{\end{itemize}}
\allowdisplaybreaks
\setlength{\emergencystretch}{3em}   % long \texttt seeds and function names

% Working space in the blank list
\newcommand{\resolespaco}{\ifsolucoes\else\par\textit{Solution:}\vspace{\espacoresol}\par\fi}

% Solution step (use inside sol)
\newcommand{\passo}[2]{\par\smallskip\noindent\textbf{Step --- #1.}~#2}
\newcommand{\intuicao}[1]{\par\smallskip\noindent\textit{Intuition.}~#1}
\newcommand{\refbib}[1]{\par\smallskip\noindent{\footnotesize\color{black!60}\textbf{Ref:} #1}}
\newcommand{\rubrica}[1]{\par\needspace{4\baselineskip}\smallskip\noindent\textbf{Grading rubric.}\begin{itemize}[leftmargin=1.4em,topsep=1pt]\small#1\end{itemize}}
\newcommand{\solhead}{\par\medskip\noindent{\color{blue!55!black}\rule{\linewidth}{0.3pt}}\par\noindent\textbf{Solution.}}
\newcommand{\codehead}{\par\smallskip\noindent\textbf{Code and output (programming part).}\par}

\newcommand{\E}{\mathbb{E}}
\newcommand{\Var}{\operatorname{Var}}
\newcommand{\Cov}{\operatorname{Cov}}
\newcommand{\1}{\mathbf{1}}

% =============================================================================
\begin{document}
\begin{center}
  {\Large\bfseries \disciplina}\\[2pt]
  {\large Exercise list: \listatema}\ifsolucoes\\[2pt]\textit{Solutions}\fi
\end{center}
\vspace{2pt}

{\footnotesize\noindent\textbf{Instructions.} Seven blocks, from Poisson likelihood theory to MCMC.
Each block has (i) textbook exercises reproduced verbatim from the book PDF (statements from
Portuguese-language books stay in Portuguese), (ii) an original math question (derivations and
proofs, ending with a parameter analysis) and (iii) an original conceptual question in exam form.
Every question carries a \emph{Programming} part added to the book's text: solve it in Python
(NumPy/SciPy/statsmodels) with the seeds given; the scaffold notebook
\texttt{Listas/lista-poisson-monte-carlo.ipynb} lists the steps. The script
\texttt{Listas/poisson-monte-carlo-check.py} reproduces and asserts every number printed in the
solutions. Multiple choice (3 points): one correct option, no penalty. True/False (2 points): score
$=\max\{0,\text{correct}-\text{wrong}/2\}\times$ item value; blank items do not count. Book
abbreviations: B = Bolfarine \& Sandoval (2010), G = Greenberg (2008, 1st ed.), I = ISLP (2023),
MS = Morettin \& Singer (2021). The coal-mining series used throughout (British mine disasters per
year, 1851--1962, Jarrett 1979) is
\par\smallskip
\noindent\texttt{1851--1878: 4 5 4 1 0 4 3 4 0 6 3 3 4 0 2 6 3 3 5 4 5 3 1 4 4 1 5 5}\\
\texttt{1879--1906: 3 4 2 5 2 2 3 4 2 1 3 2 2 1 1 1 1 3 0 0 1 0 1 1 0 0 3 1}\\
\texttt{1907--1934: 0 3 2 2 0 1 1 1 0 1 0 1 0 0 0 2 1 0 0 0 1 1 0 2 3 3 1 1}\\
\texttt{1935--1962: 2 1 1 1 1 2 4 2 0 0 0 1 4 0 0 0 1 0 0 0 0 0 1 0 0 1 0 1}
\par\smallskip\noindent
($n=112$ years, 191 disasters; 127 in the first 41 years, 1851--1891).\par}

% =============================================================================
\bloco{1}{Poisson likelihood: sufficiency, information, Cramér--Rao, UMVU}{B 2.4, 2.6, 2.5; original}

\questao{B 2.4 --- $X_1+2X_2$ is not sufficient for a Poisson mean.}
\fonte{Bolfarine, H. and Sandoval, M.~C. (2010), \textit{Introdução à Inferência Estatística},
Chapter~2, §2.6 Exercícios, Exercício~2.4, p.~32 [PDF p.~38]; the book gives no answer ---
reproduced verbatim.}

\begin{otherlanguage}{brazil}
\textbf{2.4.} Sejam $X_1,X_2$ uma amostra aleatória da variável aleatória $X\sim\mathit{Poisson}(\theta)$.
Mostre que $T=X_1+2X_2$ não é suficiente para $\theta$.
\end{otherlanguage}
\programacao Write \texttt{cond\_pmf(t, theta, w2)} that enumerates every pair $(x_1,x_2)$ with
$x_1+w_2x_2=t$ and returns $P(X_1=x_1,X_2=x_2\mid X_1+w_2X_2=t)$. Print it for $T=X_1+2X_2$
($w_2=2$) at $t=2$ and $t=4$, for $\theta=1$ and $\theta=2$. Repeat for $S=X_1+X_2$ ($w_2=1$) at
$s=4$ and check that the result is the Binomial$(4,1/2)$ p.m.f.\ for both values of $\theta$.
\resolespaco
\begin{sol}
\solhead
\begin{notacao}
\item $X_1,X_2$ i.i.d.\ Poisson$(\theta)$, $P(X=x)=\theta^xe^{-\theta}/x!$, $\theta>0$; $T=X_1+2X_2$; $S=X_1+X_2$.
\item Sufficiency (Definição 2.2.1): $P(\mathbf X=\mathbf x\mid T=t)$ does not depend on $\theta$.
\end{notacao}

\passo{Pick a value of $T$}{The pairs with $T=2$ are}
\begin{equation}\label{q1:set}
\{T=2\}=\{(X_1,X_2)=(2,0)\}\cup\{(X_1,X_2)=(0,1)\}.
\end{equation}

\passo{Probabilities}{By independence,}
\begin{align}
P(X_1=2,X_2=0)&=\frac{\theta^2e^{-\theta}}{2!}\cdot e^{-\theta}=\frac{\theta^2e^{-2\theta}}{2}, \label{q1:p20}\\
P(X_1=0,X_2=1)&=e^{-\theta}\cdot\theta e^{-\theta}=\theta e^{-2\theta}. \label{q1:p01}
\end{align}
Adding \eqref{q1:p20} and \eqref{q1:p01} as \eqref{q1:set} requires:
\begin{equation}\label{q1:pT}
P(T=2)=e^{-2\theta}\Big(\frac{\theta^2}{2}+\theta\Big).
\end{equation}

\passo{Conditional probability}{Dividing \eqref{q1:p20} by \eqref{q1:pT}:}
\[
P(X_1=2,X_2=0\mid T=2)=\frac{\theta^2e^{-2\theta}/2}{e^{-2\theta}(\theta^2/2+\theta)}=\frac{\theta^2/2}{\theta^2/2+\theta}.
\]
Dividing numerator and denominator by $\theta/2$:
\begin{equation}\label{q1:cond}
P(X_1=2,X_2=0\mid T=2)=\frac{\theta}{\theta+2},
\end{equation}
which depends on $\theta$ ($1/3=0.3333$ at $\theta=1$, $2/4=0.5$ at $\theta=2$). So $T$ is not
sufficient.

\passo{Contrast with $S$}{For $S=X_1+X_2$, $S\sim$ Poisson$(2\theta)$, and for $0\le x_1\le s$}
\begin{equation}\label{q1:S}
P(X_1=x_1\mid S=s)=\frac{\frac{\theta^{x_1}e^{-\theta}}{x_1!}\frac{\theta^{s-x_1}e^{-\theta}}{(s-x_1)!}}{\frac{(2\theta)^se^{-2\theta}}{s!}}
=\binom{s}{x_1}\Big(\frac12\Big)^{s},
\end{equation}
free of $\theta$: $S$ is sufficient (as the factorization criterion also shows).

\codehead
{\small
\begin{verbatim}
def cond_pmf(t, theta, w2):
    pts = [(t - w2 * x2, x2) for x2 in range(t // w2 + 1)]
    pr = np.array([poisson.pmf(a, theta) * poisson.pmf(b, theta) for a, b in pts])
    return pts, pr / pr.sum()
# T=2:  theta=1 -> P((2,0)|T=2) = 0.3333 ; theta=2 -> 0.5000
# T=4:  pairs (4,0),(2,1),(0,2): theta=1 -> 0.04, 0.48, 0.48
#                                 theta=2 -> 0.10, 0.60, 0.30
# S=4:  theta=1 and theta=2 -> 0.0625, 0.25, 0.375, 0.25, 0.0625 = Binomial(4, 1/2)
\end{verbatim}
}
The $t=4$ row confirms \eqref{q1:cond} is not an accident of $t=2$: the conditional law moves with
$\theta$, whereas \eqref{q1:S} does not.

\intuicao{The sufficient statistic here is $X_1+X_2$. Weighting $X_2$ twice merges sample points with
different totals: $T=2$ mixes ``two events'' with ``one event'', and the split between these cases
still carries information about $\theta$.}
\refbib{Bolfarine \& Sandoval (2010), §2.2 (estatísticas suficientes, pp.~20--23; Exemplo 2.2.4).}
\end{sol}

% -----------------------------------------------------------------------------
\questao{B 2.6 --- the information identity.}
\fonte{Bolfarine, H. and Sandoval, M.~C. (2010), \textit{Introdução à Inferência Estatística},
Chapter~2, §2.6 Exercícios, Exercício~2.6, p.~33 [PDF p.~39]; the book gives no answer ---
reproduced verbatim.}

\begin{otherlanguage}{brazil}
\textbf{2.6.} Seja $f(x|\theta)$ uma função densidade para a qual as condições de regularidade
estão satisfeitas. Mostre que
\[
E\left[\left(\frac{\partial\log f(X|\theta)}{\partial\theta}\right)^2\right]
=-E\left[\frac{\partial^2\log f(X|\theta)}{\partial\theta^2}\right].
\]
\end{otherlanguage}
\programacao (i) For $X\sim$ Poisson$(\theta)$ with $\theta=3$, compute $E[U]$, $E[U^2]$ and
$-E[\partial^2\log f/\partial\theta^2]$, with $U=\partial\log f(X|\theta)/\partial\theta$, by summing over
$x=0,\dots,399$, and compare with $1/\theta$. (ii) Do the same for $X\sim U(0,\theta)$ with
$\theta=2$, by numerical integration (\texttt{scipy.integrate.quad}), and explain the result.
\resolespaco
\begin{sol}
\solhead
\begin{notacao}
\item $f(x|\theta)$: density (sums replace integrals for a p.m.f.); $\theta$ in an open interval $\Theta$.
\item $U(\theta)=\partial\log f(X|\theta)/\partial\theta$: score of one observation;
  $I_F(\theta)=E[U^2]$: Fisher information.
\item Regularity conditions (B §2.3): the support $A=\{x:f(x|\theta)>0\}$ does not depend on
  $\theta$, and $\int f(x|\theta)dx$ can be differentiated twice under the integral sign.
\end{notacao}

\passo{The density integrates to one}{For every $\theta\in\Theta$,}
\begin{equation}\label{q2:norm}
\int_A f(x|\theta)\,dx=1 .
\end{equation}
Differentiating \eqref{q2:norm} in $\theta$ under the integral sign (regularity):
\begin{equation}\label{q2:d1}
\int_A\frac{\partial f(x|\theta)}{\partial\theta}\,dx=0 .
\end{equation}

\passo{Score has mean zero}{On $A$, the chain rule for the logarithm gives}
\begin{equation}\label{q2:dlog}
\frac{\partial f(x|\theta)}{\partial\theta}=\frac{\partial\log f(x|\theta)}{\partial\theta}\,f(x|\theta).
\end{equation}
Substituting \eqref{q2:dlog} into \eqref{q2:d1}:
\begin{equation}\label{q2:EU}
\int_A\Big(\frac{\partial\log f(x|\theta)}{\partial\theta}\,f(x|\theta)\Big)dx=0
\quad\Longleftrightarrow\quad E[U]=0 .
\end{equation}

\passo{Differentiate once more}{Differentiating the left side of \eqref{q2:EU} under the integral
sign and using the product rule:}
\begin{equation}\label{q2:d2}
\int_A\Big[\frac{\partial^2\log f(x|\theta)}{\partial\theta^2}f(x|\theta)
+\frac{\partial\log f(x|\theta)}{\partial\theta}\,\frac{\partial f(x|\theta)}{\partial\theta}\Big]dx=0 .
\end{equation}
Substituting \eqref{q2:dlog} for $\partial f/\partial\theta$ in the second term of \eqref{q2:d2}:
\[
\int_A\Big[\frac{\partial^2\log f(x|\theta)}{\partial\theta^2}f(x|\theta)
+\frac{\partial\log f(x|\theta)}{\partial\theta}\Big(\frac{\partial\log f(x|\theta)}{\partial\theta}f(x|\theta)\Big)\Big]dx=0 .
\]
Writing both integrals as expectations:
\begin{equation}\label{q2:res}
E\Big[\frac{\partial^2\log f(X|\theta)}{\partial\theta^2}\Big]+E\Big[\Big(\frac{\partial\log f(X|\theta)}{\partial\theta}\Big)^2\Big]=0,
\end{equation}
which rearranges to the identity. $\blacksquare$

\passo{Poisson check}{$\log f=x\log\theta-\theta-\log x!$, so}
\begin{equation}\label{q2:pois}
U=\frac{X}{\theta}-1,\qquad \frac{\partial^2\log f}{\partial\theta^2}=-\frac{X}{\theta^2}.
\end{equation}
From \eqref{q2:pois} and $E(X)=\Var(X)=\theta$: $E[U^2]=\Var(X)/\theta^2=1/\theta$ and
$-E[-X/\theta^2]=\theta/\theta^2=1/\theta$; both sides of \eqref{q2:res} equal $1/\theta$.

\passo{Where regularity fails}{For $U(0,\theta)$, $\log f=-\log\theta$ on $(0,\theta)$, so the
score is the constant $-1/\theta$ and $\partial^2\log f/\partial\theta^2=1/\theta^2$. Then
$E[U]=-1/\theta\neq0$, $E[U^2]=1/\theta^2$ but $-E[\partial^2\log f/\partial\theta^2]=-1/\theta^2$.
The step \eqref{q2:d1} is what breaks: by Leibniz's rule
$0=\frac{d}{d\theta}\int_0^\theta\frac1\theta dx=\frac1\theta+\int_0^\theta\big(-\frac{1}{\theta^2}\big)dx$,
the boundary term $1/\theta$ appears because the support moves with $\theta$.}

\codehead
{\small
\begin{verbatim}
th = 3.0; x = np.arange(400); f = poisson.pmf(x, th)
U, H = x / th - 1, -x / th**2
print(np.sum(f * U), np.sum(f * U**2), -np.sum(f * H))   # 0.000000  0.333333  0.333333
th = 2.0                                  # U(0, 2): score = -1/th on (0, th)
EU  = quad(lambda u: (-1 / th) / th, 0, th)[0]           # -0.5
EU2 = quad(lambda u: (1 / th**2) / th, 0, th)[0]         #  0.25
mEH = -quad(lambda u: (1 / th**2) / th, 0, th)[0]        # -0.25
\end{verbatim}
}
The Poisson values equal $1/\theta=0.333333$; for the uniform $E[U^2]=0.25$ but
$-E[\partial^2\log f/\partial\theta^2]=-0.25$ and $E[U]=-0.5$, as derived.

\intuicao{The score has mean zero because the model's total probability cannot change with
$\theta$; information measures how much the score fluctuates, and the curvature of the expected
log-likelihood measures the same thing. When the support depends on $\theta$, the data can rule
values out sharply, the likelihood is not smooth at its peak, and both identities fail.}
\refbib{Bolfarine \& Sandoval (2010), §2.3 (desigualdade da informação, pp.~23--28; Exemplo 2.3.1,
Poisson); §3.1 (Exemplo 3.1.1, the uniform).}
\end{sol}

% -----------------------------------------------------------------------------
\questao{B 2.5 --- the Cramér--Rao bound for $g(\theta)$.}
\fonte{Bolfarine, H. and Sandoval, M.~C. (2010), \textit{Introdução à Inferência Estatística},
Chapter~2, §2.6 Exercícios, Exercício~2.5, p.~33 [PDF p.~39]; the book gives no answer ---
reproduced verbatim.}

\begin{otherlanguage}{brazil}
\textbf{2.5.} Sejam $X_1,\dots,X_n$ uma amostra aleatória da variável aleatória $X$ com função de
densidade (ou de probabilidade) $f(x|\theta)$ para a qual as condições de regularidade estão
satisfeitas. Seja $\hat\gamma$ um estimador não viciado para $g(\theta)$. Mostre que
\[
Var(\hat\gamma)\ge\frac{(g'(\theta))^2}{nE\left[\left(\frac{\partial\log f(X|\theta)}{\partial\theta}\right)^2\right]}.
\]
\end{otherlanguage}
\programacao Take $X\sim$ Poisson$(\theta)$, $g(\theta)=\theta$, $\theta=2$, $n=10$. With
\texttt{np.random.default\_rng(3)} draw $200{,}000$ samples of size $n$; for each compute
$\overline X$ and $S^2=\sum_i(X_i-\overline X)^2/(n-1)$ (both unbiased for $\theta$). Print the
simulated means and variances of both estimators next to the bound $\theta/n$ and to the exact
$\Var(S^2)=\theta/n+2\theta^2/(n-1)$.
\resolespaco
\begin{sol}
\solhead
\begin{notacao}
\item $\mathbf X=(X_1,\dots,X_n)$ i.i.d.\ $f(x|\theta)$; $L(\theta;\mathbf x)=\prod_if(x_i|\theta)$.
\item $U_n=\sum_{i=1}^n\partial\log f(X_i|\theta)/\partial\theta$: score of the sample;
  $I_F(\theta)=E[(\partial\log f(X|\theta)/\partial\theta)^2]$.
\item $\hat\gamma=\hat\gamma(\mathbf X)$ with $E_\theta(\hat\gamma)=g(\theta)$ for all $\theta$; $g$ differentiable.
\end{notacao}

\passo{Unbiasedness, differentiated}{Unbiasedness reads}
\begin{equation}\label{q3:unb}
\int\hat\gamma(\mathbf x)\,L(\theta;\mathbf x)\,d\mathbf x=g(\theta).
\end{equation}
Differentiating \eqref{q3:unb} under the integral (regularity) and substituting
$\partial L/\partial\theta=U_n(\mathbf x)L$, the sample analogue of \eqref{q2:dlog}:
\begin{equation}\label{q3:cov0}
\int\hat\gamma(\mathbf x)\,\big(U_n(\mathbf x)L(\theta;\mathbf x)\big)\,d\mathbf x=g'(\theta)
\quad\Longleftrightarrow\quad E[\hat\gamma\,U_n]=g'(\theta).
\end{equation}

\passo{It is a covariance}{Each term of $U_n$ has mean zero by \eqref{q2:EU} (Question~2), so
$E[U_n]=0$ and \eqref{q3:cov0} becomes}
\begin{equation}\label{q3:cov}
\Cov(\hat\gamma,U_n)=E[\hat\gamma U_n]-E[\hat\gamma]\,(0)=g'(\theta).
\end{equation}

\passo{Variance of the score}{The terms of $U_n$ are i.i.d.\ with mean zero, so}
\begin{equation}\label{q3:varU}
\Var(U_n)=\sum_{i=1}^n\Var\Big(\frac{\partial\log f(X_i|\theta)}{\partial\theta}\Big)=nE\Big[\Big(\frac{\partial\log f(X|\theta)}{\partial\theta}\Big)^2\Big]=nI_F(\theta).
\end{equation}

\passo{Cauchy--Schwarz}{For any two variables with finite variance,
$\Cov(A,B)^2\le\Var(A)\Var(B)$. With $A=\hat\gamma$, $B=U_n$, substituting \eqref{q3:cov} and
\eqref{q3:varU}:}
\[
(g'(\theta))^2\le\Var(\hat\gamma)\,\big(nI_F(\theta)\big).
\]
Dividing by $nI_F(\theta)>0$:
\begin{equation}\label{q3:CR}
\Var(\hat\gamma)\ge\frac{(g'(\theta))^2}{nI_F(\theta)} .\qquad\blacksquare
\end{equation}
Equality in Cauchy--Schwarz holds iff $U_n=a(\theta)\,(\hat\gamma-g(\theta))$ with probability one:
the bound is attained only by estimators that are affine in the score.

\passo{Poisson, $g(\theta)=\theta$}{From \eqref{q2:pois}, $I_F(\theta)=1/\theta$, so
\eqref{q3:CR} gives the bound $\theta/n$. $U_n=\sum_i(X_i/\theta-1)=\frac{n}{\theta}(\overline X-\theta)$
is affine in $\overline X$, so $\overline X$ attains it: $\Var(\overline X)=\theta/n$. The sample
variance $S^2$ is also unbiased ($\Var X=\theta$), but by the standard formula
$\Var(S^2)=\mu_4/n-\sigma^4(n-3)/(n(n-1))$ with the Poisson moments $\sigma^2=\theta$,
$\mu_4=\theta+3\theta^2$:}
\begin{equation}\label{q3:varS2}
\Var(S^2)=\frac{\theta+3\theta^2}{n}-\frac{\theta^2(n-3)}{n(n-1)}=\frac{\theta}{n}+\frac{2\theta^2}{n-1}>\frac{\theta}{n}.
\end{equation}
At $\theta=2$, $n=10$: bound $=2/10=0.2$, and \eqref{q3:varS2} gives
$0.2+2(4)/9=0.2+0.888889=1.088889$.

\codehead
{\small
\begin{verbatim}
rng = np.random.default_rng(3)
X = rng.poisson(2.0, size=(200_000, 10))
xb, s2 = X.mean(1), X.var(1, ddof=1)
print(xb.mean(), xb.var(ddof=1))   # 2.0014  0.1997   (bound 0.2)
print(s2.mean(), s2.var(ddof=1))   # 1.9995  1.0871   (exact 1.088889)
\end{verbatim}
}
Both estimators are unbiased; only $\overline X$ sits at the bound, and $S^2$ is more than five
times as variable, as \eqref{q3:varS2} predicts.

\intuicao{An unbiased estimator must move one-for-one with $g(\theta)$; the score is the direction
in which the data move most when $\theta$ moves. The covariance between the two is pinned at
$g'(\theta)$, so an estimator can only be precise if it loads heavily on the score. $\overline X$ is
the score up to scale; $S^2$ also reacts to the spread of the data, which adds noise.}
\refbib{Bolfarine \& Sandoval (2010), §2.3, Teorema 2.3.1 and Corolário (pp.~24--27).}
\end{sol}

% -----------------------------------------------------------------------------
\questao{Original math question --- estimating the probability of an accident-free day.}
\fonte{Original question (not from a book).}

The Federal Highway Police (PRF) counts serious accidents per day on a stretch of the BR-116. Model
the daily counts $X_1,\dots,X_n$ as i.i.d.\ Poisson$(\theta)$, $\theta>0$, and let
$T=\sum_{i=1}^nX_i$. The indicator the PRF reports is $\psi=P(X=0)=e^{-\theta}$.
\begin{enumerate}[label=\textbf{(\alph*)},leftmargin=2.4em]
\item Find the maximum likelihood estimator $\hat\psi$ of $\psi$, its exact expectation, and the sign
of its bias for every $n\ge1$ and $\theta>0$.
\item Find the uniformly minimum-variance unbiased estimator (UMVU) $\tilde\psi$ of $\psi$. Prove
every step (sufficiency and completeness of $T$ may be quoted from B, Exemplo~2.5.4).
\item Find $\Var(\tilde\psi)$ in closed form and compare it with the Cramér--Rao bound for
$g(\theta)=e^{-\theta}$. Is there an efficient unbiased estimator of $\psi$ for finite $n$? Prove your
answer.
\item Parameter analysis: let $\rho(\theta,n)=\Var(\tilde\psi)/\text{CR bound}$. Prove the signs of
$\partial\rho/\partial\theta$ and of the change in $\rho$ with $n$, find the limits $n\to\infty$,
$\theta\to0$ and the value at $n=1$, and interpret them for the PRF.
\end{enumerate}
\programacao With $\theta=1.5$, $n=10$ and \texttt{np.random.default\_rng(4)}, simulate
$200{,}000$ samples; print the simulated mean of $\hat\psi$ against (a), the mean and variance of
$\tilde\psi$ against (b)--(c), the bound, and the mean squared error of $\hat\psi$. Then print
$\rho(1.5,n)$ for $n=1,2,5,10,50,1000$.
\resolespaco
\begin{sol}
\solhead
\begin{notacao}
\item $X_i\sim$ Poisson$(\theta)$ i.i.d., $\theta>0$; $T=\sum_iX_i\sim$ Poisson$(n\theta)$; $\overline X=T/n$.
\item $\psi=e^{-\theta}\in(0,1)$; $\hat\psi$: MLE; $\tilde\psi$: UMVU; $W=\1\{X_1=0\}$.
\item $G_T(s)=E[s^T]$, $s\ge0$: probability generating function of $T$.
\item $\rho(\theta,n)=\Var(\tilde\psi)/\mathrm{CR}(\theta,n)$; $x=\theta/n>0$.
\end{notacao}

\passo{(a) The generating function of $T$}{For $\mu=n\theta$,}
\[
G_T(s)=\sum_{t=0}^\infty s^t\frac{e^{-\mu}\mu^t}{t!}=e^{-\mu}\sum_{t=0}^\infty\frac{(s\mu)^t}{t!}=e^{-\mu}e^{s\mu},
\]
that is,
\begin{equation}\label{q4:pgf}
G_T(s)=\exp\{n\theta(s-1)\}.
\end{equation}

\passo{(a) MLE and its mean}{The MLE of $\theta$ is $\overline X$ (Question~14, \eqref{q14:score});
by invariance of the MLE (B, §3.1), $\hat\psi=e^{-\overline X}=(e^{-1/n})^T$. Substituting
$s=e^{-1/n}$ into \eqref{q4:pgf}:}
\begin{equation}\label{q4:Emle}
E[\hat\psi]=\exp\{n\theta(e^{-1/n}-1)\}.
\end{equation}

\passo{(a) Sign of the bias}{$e^{-u}$ is strictly convex, so it lies above its tangent at $0$:
$e^{-u}>1-u$ for $u\ne0$. With $u=1/n$:}
\begin{equation}\label{q4:tan}
e^{-1/n}-1>-\frac1n .
\end{equation}
Multiplying \eqref{q4:tan} by $n\theta>0$ and exponentiating (an increasing map), \eqref{q4:Emle}
gives
\[
E[\hat\psi]>\exp\{n\theta(-1/n)\}=e^{-\theta}=\psi .
\]
The MLE is biased upwards for every $n$ and $\theta$.

\passo{(b) An unbiased start}{$E[W]=P(X_1=0)=e^{-\theta}=\psi$, so $W$ is unbiased. By
Rao--Blackwell, $E[W\mid T]$ is unbiased with variance no larger; by Lehmann--Scheffé (B,
Teorema~2.5.3), with $T$ complete and sufficient, it is the UMVU.}

\passo{(b) Conditional law of $X_1$ given $T$}{$X_1$ and $\sum_{i\ge2}X_i\sim$ Poisson$((n-1)\theta)$
are independent, so for $0\le x_1\le t$:}
\[
P(X_1=x_1\mid T=t)=\frac{\dfrac{\theta^{x_1}e^{-\theta}}{x_1!}\cdot\dfrac{((n-1)\theta)^{t-x_1}e^{-(n-1)\theta}}{(t-x_1)!}}{\dfrac{(n\theta)^te^{-n\theta}}{t!}} .
\]
Cancelling $\theta^te^{-n\theta}$:
\begin{equation}\label{q4:bin}
P(X_1=x_1\mid T=t)=\binom{t}{x_1}\Big(\frac1n\Big)^{x_1}\Big(1-\frac1n\Big)^{t-x_1},
\end{equation}
a Binomial$(t,1/n)$, free of $\theta$.

\passo{(b) The UMVU}{Setting $x_1=0$ in \eqref{q4:bin}:}
\begin{equation}\label{q4:umvu}
\tilde\psi=E[W\mid T]=P(X_1=0\mid T)=\Big(1-\frac1n\Big)^{T}.
\end{equation}
Check: substituting $s=1-1/n$ into \eqref{q4:pgf}, $E[\tilde\psi]=\exp\{n\theta(-1/n)\}=e^{-\theta}$.

\passo{(c) Variance of the UMVU}{$\tilde\psi^2=((1-1/n)^2)^T$; substituting $s=(1-1/n)^2$ into
\eqref{q4:pgf}:}
\[
E[\tilde\psi^2]=\exp\Big\{n\theta\Big(\Big(1-\frac1n\Big)^2-1\Big)\Big\}=\exp\Big\{n\theta\Big(-\frac2n+\frac1{n^2}\Big)\Big\}=e^{-2\theta+\theta/n}.
\]
Subtracting $(E\tilde\psi)^2=e^{-2\theta}$:
\begin{equation}\label{q4:varU}
\Var(\tilde\psi)=e^{-2\theta}\big(e^{\theta/n}-1\big).
\end{equation}

\passo{(c) Cramér--Rao bound}{With $g(\theta)=e^{-\theta}$, $(g'(\theta))^2=e^{-2\theta}$, and
$I_F=1/\theta$ (Question~2), \eqref{q3:CR} gives}
\begin{equation}\label{q4:CR}
\mathrm{CR}(\theta,n)=\frac{e^{-2\theta}}{n(1/\theta)}=\frac{\theta e^{-2\theta}}{n}.
\end{equation}
Dividing \eqref{q4:varU} by \eqref{q4:CR}, with $x=\theta/n$:
\begin{equation}\label{q4:rho}
\rho(\theta,n)=\frac{e^{x}-1}{x}=\frac{1}{x}\sum_{k\ge1}\frac{x^k}{k!}=1+\frac x2+\frac{x^2}{6}+\dots>1 .
\end{equation}

\passo{(c) No efficient estimator}{By \eqref{q4:rho} the UMVU, which has the smallest variance among
unbiased estimators, stays strictly above the bound; hence no unbiased estimator of $e^{-\theta}$
attains it for finite $n$. Equivalently, by the equality condition after \eqref{q3:CR}, an efficient
estimator would be affine in $U_n=\frac n\theta(\overline X-\theta)$, i.e.\ $a+b\overline X$, whose
mean $a+b\theta$ is linear in $\theta$ and cannot equal $e^{-\theta}$. So a UMVU exists while an
efficient estimator does not.}

\passo{(d) Signs}{Let $r(x)=(e^x-1)/x$. Its derivative is}
\begin{equation}\label{q4:dr}
r'(x)=\frac{xe^x-(e^x-1)}{x^2}=\frac{h(x)}{x^2},\qquad h(x)=xe^x-e^x+1 .
\end{equation}
$h(0)=0$ and $h'(x)=e^x+xe^x-e^x=xe^x>0$ for $x>0$, so $h(x)>0$ and, by \eqref{q4:dr}, $r'(x)>0$.
Since $x=\theta/n$:
\[
\frac{\partial\rho}{\partial\theta}=\frac1n\,r'(x)>0,\qquad \rho\text{ decreases as }n\text{ grows}.
\]

\passo{(d) Limits}{From \eqref{q4:rho}: $\rho\to1$ as $x\to0$, i.e.\ as $n\to\infty$ (fixed $\theta$)
or $\theta\to0$ (fixed $n$); at $n=1$, $\rho=(e^\theta-1)/\theta$ and $\tilde\psi=(1-1)^{X_1}=\1\{X_1=0\}=W$
(with $0^0=1$): with one day of data the only unbiased estimator based on $T$ is the indicator itself.}

\passo{(d) Reading}{The bound is nearly attained when accidents are rare relative to the sample
(small $\theta/n$): the PRF loses little by using any reasonable estimator. On busy stretches
(large $\theta$) with few days of data, the bound is far from attainable, and it understates the
real uncertainty of the indicator. Asymptotically both $\hat\psi$ and $\tilde\psi$ are efficient.}

\passo{Numbers}{$\theta=1.5$, $n=10$: $\psi=e^{-1.5}=0.223130$; \eqref{q4:Emle}:
$\exp\{15(e^{-0.1}-1)\}=0.239923$, bias $0.016792$; \eqref{q4:varU}:
$e^{-3}(e^{0.15}-1)=0.008057$; \eqref{q4:CR}: $1.5e^{-3}/10=0.007468$; \eqref{q4:rho}:
$\rho=0.008057/0.007468=1.078895$.}

\codehead
{\small
\begin{verbatim}
rng = np.random.default_rng(4)
T = rng.poisson(1.5, size=(200_000, 10)).sum(1)
mle, umvu = np.exp(-T / 10), (1 - 1 / 10) ** T
print(mle.mean(), umvu.mean(), umvu.var(ddof=1))   # 0.239825  0.223033  0.008042
print(np.mean((mle - np.exp(-1.5))**2))            # 0.008638 (exact MSE 0.008657)
rho = lambda n: (np.exp(1.5 / n) - 1) / (1.5 / n)
# n = 1, 2, 5, 10, 50, 1000 -> 2.321126 1.489333 1.166196 1.078895 1.015151 1.000750
\end{verbatim}
}
The simulated values match (a)--(c) within Monte Carlo error; the $\rho$ column falls
monotonically to 1, as (d) proves.

\intuicao{The MLE plugs $\overline X$ into a convex function, so Jensen pushes it up. The UMVU
replaces $e^{-1/n}$ by its first-order approximation $1-1/n$, exactly the correction that removes
the bias. The Cramér--Rao bound is a local, linear benchmark: it is attainable only for parameters
that are linear in the sufficient statistic's mean, which is why $e^{-\theta}$ misses it.}
\refbib{Bolfarine \& Sandoval (2010), §2.3 (Cramér--Rao), §2.5 (Teorema 2.5.3 and Exemplo 2.5.4),
§3.1 (invariance); Exercício 2.10 for the same Rao--Blackwell step with Bernoulli data.}
\rubrica{
\item (a) MLE by invariance and $E[\hat\psi]$ via the p.g.f.: 2 pts; bias sign proved (convexity): 1 pt.
\item (b) $W$ unbiased and Rao--Blackwell/Lehmann--Scheffé named: 1 pt; \eqref{q4:bin} derived: 2 pts; \eqref{q4:umvu}: 1 pt.
\item (c) \eqref{q4:varU}: 2 pts; \eqref{q4:CR}: 1 pt; $\rho>1$ and no efficient estimator, proved: 2 pts.
\item (d) sign of $r'$ proved: 2 pts; limits: 1 pt; reading: 1 pt. Signs stated without proof: max.\ half of (d).
\item Programming: correct simulation and comparison with (a)--(c): 2 pts.
}
\end{sol}

% -----------------------------------------------------------------------------
\questao{Original conceptual question --- sufficiency, information and optimal estimators (10 points).}
\fonte{Original question (not from a book).}

The PRF models the daily number of accidents on a stretch of road as $X_1,\dots,X_n$ i.i.d.\
Poisson$(\theta)$, $T=\sum_iX_i$.

\textbf{a)} (3 points) Which statement about $T$ is correct?
\begin{alternativas}
\item $T$ is sufficient for $\theta$ because its expectation, $n\theta$, depends on $\theta$.
\item $T$ is sufficient for $\theta$ because the law of the sample given $T=t$ is multinomial $(t;1/n,\dots,1/n)$, free of $\theta$.
\item $T$ is sufficient for $\theta$ only for large $n$, when $\overline X$ is approximately normal by the CLT.
\item $\overline X$ is sufficient for $\theta$ but $T$ is not, because only $\overline X$ is unbiased for $\theta$.
\end{alternativas}
\gabarito{B}
\begin{sol}
\passo{Why}{B: by the computation of \eqref{q4:bin} for all coordinates, $P(\mathbf X=\mathbf x\mid T=t)=\frac{t!}{\prod x_i!}n^{-t}$, free of $\theta$ (Definição 2.2.1). A: a statistic whose mean depends on $\theta$ need not be sufficient ($X_1$ alone is a counterexample). C: sufficiency is an exact finite-sample property, unrelated to the CLT. D: $T$ and $\overline X=T/n$ are one-to-one, so they carry the same information; sufficiency does not involve unbiasedness.}
\refbib{B §2.2 (Definição 2.2.1, Exemplo 2.2.4).}
\end{sol}

\textbf{b)} (3 points) The Cramér--Rao lower bound for the variance of unbiased estimators of $\theta^2$ is:
\begin{alternativas}
\item $\theta/n$ (the bound for $\theta$ itself).
\item $4\theta^2/n$ (squared derivative, unit information).
\item $2\theta^3/n$ (derivative not squared).
\item $4\theta^3/n$ (squared derivative over $n/\theta$).
\item $4\theta^3/n^2$ (information counted twice).
\end{alternativas}
\gabarito{D}
\begin{sol}
\passo{Why}{By \eqref{q3:CR} with $g(\theta)=\theta^2$, $(g')^2=4\theta^2$ and $nI_F=n/\theta$: bound $=4\theta^2\theta/n=4\theta^3/n$ (D). A ignores $g$; B drops $I_F=1/\theta$; C forgets to square $g'$; E divides by $n$ twice.}
\end{sol}

\textbf{c)} (2 points, T or F) If an unbiased estimator of $g(\theta)$ attains the Cramér--Rao bound at every $\theta$, then it is the UMVU of $g(\theta)$.
\gabarito{T}
\begin{sol}
\passo{Why}{Every unbiased estimator has variance at least the bound \eqref{q3:CR}; one that attains it has the minimum variance at every $\theta$, which is the definition of UMVU.}
\end{sol}

\textbf{d)} (2 points, T or F) If a UMVU estimator of $g(\theta)$ exists, then it attains the Cramér--Rao bound.
\gabarito{F}
\begin{sol}
\passo{Why}{The converse of (c) fails. Question~4: $\tilde\psi=(1-1/n)^T$ is the UMVU of $e^{-\theta}$, yet $\Var(\tilde\psi)/\mathrm{CR}=\rho>1$ by \eqref{q4:rho}.}
\refbib{B §2.3--2.5.}
\end{sol}

\programacao For $\theta=2$, $n=10$: $T(T-1)/n^2=\overline X^2-\overline X/n$ is unbiased for $\theta^2$
and a function of the complete sufficient $T$, hence the UMVU. Compute its exact variance by summing
over $T\sim$ Poisson$(n\theta)$ ($t=0,\dots,399$) and compare it with the bound of (b).
\begin{sol}
\codehead
{\small
\begin{verbatim}
t = np.arange(400); pt = poisson.pmf(t, 20); est = t * (t - 1) / 100
np.sum(pt * est), np.sum(pt * est**2) - 16      # 4.0 (unbiased)  3.28
4 * 2**3 / 10                                   # bound 3.2 ; ratio 1.025
\end{verbatim}
}
The exact variance is $4\theta^3/n+2\theta^2/n^2=3.2+0.08=3.28>3.2$: a second UMVU that misses the
bound, confirming that (d) is false.
\end{sol}

% =============================================================================
\bloco{2}{Poisson tests and intervals}{B 6.5, 6.11; original}

\questao{B 6.5 --- UMP test for a Poisson mean and its power function.}
\fonte{Bolfarine, H. and Sandoval, M.~C. (2010), \textit{Introdução à Inferência Estatística},
Chapter~6, §6.7 Exercícios, Exercício~6.5, p.~115 [PDF p.~121]; the book gives no answer ---
reproduced verbatim.}

\begin{otherlanguage}{brazil}
\textbf{6.5.} Sejam $X_1,\dots,X_n$ uma amostra aleatória da variável aleatória $X\sim\mathit{Poisson}(\theta)$.\\
i) Encontre o teste UMP para testar $H_0:\theta=\theta_0$ versus $H_1:\theta>\theta_0$.\\
ii) Seja $\alpha=0,05$, faça o gráfico da função poder para $\theta_0=1$ e $n=25$ (use o Teorema
do limite central).
\end{otherlanguage}
\programacao For $\theta_0=1$, $n=25$, $\alpha=0.05$: (1) find the smallest integer $c$ with
$P(T>c\mid\theta_0)\le\alpha$, $T=\sum_iX_i$, and the randomization probability $\gamma$ that makes
the size exactly $0.05$; (2) print, for $\theta\in\{1,1.2,1.5,2\}$, the CLT power from (ii), the
exact power of the test ``reject if $T>c$'' and that of the randomized test; (3) plot the three
power curves on $\theta\in[0.8,2.2]$.
\resolespaco
\begin{sol}
\solhead
\begin{notacao}
\item $X_i\sim$ Poisson$(\theta)$ i.i.d.; $T=\sum_iX_i\sim$ Poisson$(n\theta)$; $\overline X=T/n$.
\item $\phi(T)\in[0,1]$: test function (probability of rejecting $H_0$); $\pi(\theta)=E_\theta[\phi(T)]$: power.
\item $c\in\{0,1,\dots\}$, $\gamma\in[0,1)$: critical value and randomization; $z_{0.95}=1.644854$.
\end{notacao}

\passo{(i) Likelihood ratio}{For a fixed $\theta_1>\theta_0$,}
\begin{equation}\label{q6:lr}
\frac{L(\theta_1;\mathbf x)}{L(\theta_0;\mathbf x)}=\frac{\theta_1^te^{-n\theta_1}/\prod x_i!}{\theta_0^te^{-n\theta_0}/\prod x_i!}
=\Big(\frac{\theta_1}{\theta_0}\Big)^{t}e^{-n(\theta_1-\theta_0)},
\end{equation}
strictly increasing in $t$ because $\theta_1/\theta_0>1$.

\passo{(i) Neyman--Pearson}{The most powerful size-$\alpha$ test of $\theta_0$ against $\theta_1$
rejects when \eqref{q6:lr} exceeds a constant (B, Lema de Neyman--Pearson); by monotonicity, when
$T$ exceeds a constant. With the randomization needed for a discrete $T$:}
\begin{equation}\label{q6:phi}
\phi(T)=\begin{cases}1,&T>c,\\ \gamma,&T=c,\\ 0,&T<c,\end{cases}\qquad
P_{\theta_0}(T>c)+\gamma P_{\theta_0}(T=c)=\alpha .
\end{equation}
$c$ and $\gamma$ depend only on $\theta_0$, $n$, $\alpha$, not on $\theta_1$: the same test is most
powerful against every $\theta_1>\theta_0$, hence UMP for $H_1:\theta>\theta_0$ (B, §6.3: families
with monotone likelihood ratio in $T$).

\passo{(ii) CLT version}{Under $\theta$, $\overline X\approx N(\theta,\theta/n)$. The size-$\alpha$
rule rejects when $\overline X>k$ with}
\begin{equation}\label{q6:k}
k=\theta_0+z_{0.95}\sqrt{\theta_0/n}.
\end{equation}
Substituting $\theta_0=1$, $n=25$ into \eqref{q6:k}:
\[
k=1+(1.644854)\sqrt{1/25}=1+1.644854\,(0.2)=1.328971 .
\]
The power function is
\begin{equation}\label{q6:pow}
\pi(\theta)=P_\theta(\overline X>k)\approx1-\Phi\Big(\frac{k-\theta}{\sqrt{\theta/n}}\Big).
\end{equation}
Substituting $\theta=1.5$: $\pi(1.5)\approx1-\Phi\big((1.328971-1.5)/\sqrt{1.5/25}\big)=1-\Phi(-0.698223)=0.757481$;
likewise $\pi(1)=0.05$, $\pi(1.2)=0.278042$, $\pi(2)=0.991165$.
\begin{center}
\begin{tikzpicture}
\begin{axis}[width=0.7\linewidth,height=5cm,xlabel={$\theta$},ylabel={$\pi(\theta)$},
  xmin=0.8,xmax=2.2,ymin=0,ymax=1.02,grid=major,legend pos=south east,legend style={font=\scriptsize}]
\addplot[blue!60!black,thick,mark=*,mark size=1.2pt] coordinates {(0.8,0.0016) (0.9,0.0119) (1.0,0.0500)
  (1.1,0.1375) (1.2,0.2780) (1.3,0.4495) (1.4,0.6180) (1.5,0.7575) (1.6,0.8580) (1.7,0.9226)
  (1.8,0.9604) (1.9,0.9808) (2.0,0.9912) (2.1,0.9961) (2.2,0.9983)};
\addlegendentry{CLT approximation \eqref{q6:pow}}
\addplot[red!70!black,dashed] coordinates {(0.8,0.05) (2.2,0.05)};
\addlegendentry{$\alpha=0.05$}
\end{axis}
\end{tikzpicture}
\end{center}

\passo{Exact test (programming)}{$T\sim$ Poisson$(25)$ under $H_0$: the smallest $c$ with
$P(T>c)\le0.05$ is $c=33$, size $P(T>33)=0.049780$; \eqref{q6:phi} gives
$\gamma=(0.05-0.049780)/P(T=33)=0.010133$. Since $25k=33.22$, the CLT rule $\overline X>k$ is the
same as $T>33$: its true size is $0.049780$, and its true power is the exact column below.}

\codehead
{\small
\begin{verbatim}
c = 33; size = poisson.sf(c, 25)                       # 0.049780
gam = (0.05 - size) / poisson.pmf(c, 25)               # 0.010133
for th in (1, 1.2, 1.5, 2):
    clt = 1 - norm.cdf((1.328971 - th) / np.sqrt(th / 25))
    ex = poisson.sf(c, 25 * th); rnd = ex + gam * poisson.pmf(c, 25 * th)
# theta   CLT       exact     randomized
# 1.0     0.050000  0.049780  0.050000
# 1.2     0.278042  0.255551  0.256158
# 1.5     0.757481  0.738112  0.738642
# 2.0     0.991165  0.993021  0.993047
\end{verbatim}
}
The CLT curve overstates the power near $\theta_0$ by about two points (the Poisson is right-skewed)
and understates it slightly at $\theta=2$; randomization adds at most $\gamma P(T=33)$.

\intuicao{All the information about $\theta$ is in the total count, and larger totals always point
to larger rates, so ``reject when the count is large'' is optimal against every larger rate at
once. The discreteness of counts is why an exact 5\% test needs a coin flip at the boundary.}
\refbib{Bolfarine \& Sandoval (2010), §6.2 (Lema de Neyman--Pearson), §6.3 (testes UMP), Exemplo 6.3.2.}
\end{sol}

% -----------------------------------------------------------------------------
\questao{B 6.11 --- likelihood ratio test for two Poisson means.}
\fonte{Bolfarine, H. and Sandoval, M.~C. (2010), \textit{Introdução à Inferência Estatística},
Chapter~6, §6.7 Exercícios, Exercício~6.11, p.~116 [PDF p.~122]; the book gives no answer ---
reproduced verbatim.}

\begin{otherlanguage}{brazil}
\textbf{6.11.} Sejam $X_1,\dots,X_n$ uma amostra aleatória da variável aleatória $X\sim\mathit{Poisson}(\theta_1)$
e sejam $Y_1,\dots,Y_m$ uma amostra aleatória da variável aleatória $Y\sim\mathit{Poisson}(\theta_2)$
sendo as amostras independentes.\\
i) Encontre o teste da RVG(aproximado) para testar $H_0:\theta_1=\theta_2$ versus $H_1:\theta_1\neq\theta_2$.\\
ii) Sendo $n=5$, $\sum x_i=3,8$; $m=8$; $\sum y_i=4,8$, qual a sua conclusão a um nível de
significância de 5\%?
\end{otherlanguage}
\programacao Write \texttt{lr2(sx, sy, n, m)} returning $-2\log\lambda$ from the two totals. Apply it
to the printed data. Then estimate the true size of the approximate test (reject when
$-2\log\lambda>3.841$) for $n=5$, $m=8$ by simulating $100{,}000$ pairs of samples under $H_0$ with
$\theta_1=\theta_2=0.66$ and with $\theta_1=\theta_2=3$ (\texttt{np.random.default\_rng(7)}).
\resolespaco
\begin{sol}
\solhead
\begin{notacao}
\item $s_x=\sum_ix_i$, $s_y=\sum_jy_j$; $\overline x=s_x/n$, $\overline y=s_y/m$.
\item $\hat\theta_0=(s_x+s_y)/(n+m)$: MLE of the common mean under $H_0$.
\item $\lambda(\mathbf x,\mathbf y)$: generalized likelihood ratio (RVG); $\chi^2_{1,0.95}=3.841459$.
\end{notacao}

\passo{(i) Likelihood}{By independence,}
\begin{equation}\label{q7:L}
L(\theta_1,\theta_2)=\frac{\theta_1^{s_x}e^{-n\theta_1}\,\theta_2^{s_y}e^{-m\theta_2}}{\prod_ix_i!\prod_jy_j!}.
\end{equation}
Maximizing $\log L$ in each rate separately gives $\hat\theta_1=\overline x$ and
$\hat\theta_2=\overline y$ (as in \eqref{q14:score}). Under $H_0$, \eqref{q7:L} becomes
\[
L(\theta,\theta)=\frac{\theta^{s_x+s_y}e^{-(n+m)\theta}}{\prod_ix_i!\prod_jy_j!},
\]
maximized at $\hat\theta_0$.

\passo{(i) The statistic}{Dividing the restricted by the unrestricted maximum (factorials cancel):}
\[
\lambda=\frac{\hat\theta_0^{\,s_x+s_y}e^{-(n+m)\hat\theta_0}}{\overline x^{\,s_x}e^{-n\overline x}\,\overline y^{\,s_y}e^{-m\overline y}} .
\]
Taking $-2\log$:
\[
-2\log\lambda=2\Big[s_x\log\frac{\overline x}{\hat\theta_0}+s_y\log\frac{\overline y}{\hat\theta_0}\Big]
-2\big[n\overline x+m\overline y-(n+m)\hat\theta_0\big].
\]
The second bracket is $(s_x+s_y)-(s_x+s_y)=0$, so
\begin{equation}\label{q7:stat}
-2\log\lambda=2\Big[s_x\log\frac{\overline x}{\hat\theta_0}+s_y\log\frac{\overline y}{\hat\theta_0}\Big].
\end{equation}
$H_1$ has two free parameters and $H_0$ one, so under $H_0$, $-2\log\lambda\to\chi^2_1$ in
distribution (B, Teorema~6.5.1): reject at 5\% when \eqref{q7:stat} exceeds $3.841459$.

\passo{(ii) The printed data}{\emph{Book error:} a sum of Poisson counts is an integer, but the book
prints $\sum x_i=3{,}8$ and $\sum y_i=4{,}8$ (reading them as means does not help either:
$m\overline y=38.4$). Statistic \eqref{q7:stat} only needs the totals, so we evaluate it with the
printed values: $\overline x=3.8/5=0.76$, $\overline y=4.8/8=0.6$, $\hat\theta_0=8.6/13=0.661538$.}
Substituting into \eqref{q7:stat}:
\[
-2\log\lambda=2\Big[(3.8)\log\frac{0.76}{0.661538}+(4.8)\log\frac{0.6}{0.661538}\Big]
=2\,[0.527251+(-0.468665)]=0.117173 .
\]
Since $0.117173<3.841459$ (p-value $0.7321$), $H_0:\theta_1=\theta_2$ is not rejected at 5\%.

\codehead
{\small
\begin{verbatim}
def lr2(sx, sy, n, m):
    th0 = (sx + sy) / (n + m)
    t = lambda s, k: np.where(s > 0, s * np.log(np.where(s > 0, s, 1) / (k * th0)), 0.)
    return 2 * (t(sx, n) + t(sy, m))           # 0 log 0 = 0
lr2(3.8, 4.8, 5, 8)                            # 0.117173 ; p = chi2.sf(., 1) = 0.7321
rng = np.random.default_rng(7)
for th in (0.66, 3.0):
    sx, sy = rng.poisson(5 * th, 100_000), rng.poisson(8 * th, 100_000)
    print(np.mean(lr2(sx, sy, 5, 8) > 3.841459))   # 0.06949 (th=0.66), 0.05167 (th=3)
\end{verbatim}
}
With expected totals of only 3.3 and 5.3 events the $\chi^2_1$ approximation over-rejects (size
about 7\%); with means near 3 per observation it is close to 5\%.

\intuicao{The statistic compares how well one common rate explains both samples against two separate
rates; it is twice the gain in log-likelihood from letting the rates differ. Wilks' $\chi^2_1$ is a
large-count result: with a handful of events its tail is too light.}
\refbib{Bolfarine \& Sandoval (2010), §6.5 (razão de verossimilhanças generalizada, Teorema 6.5.1,
Exemplo 6.5.5 for one Poisson sample).}
\end{sol}

% -----------------------------------------------------------------------------
\questao{Original math question --- three confidence intervals for a Poisson rate.}
\fonte{Original question (not from a book).}

The National Civil Defence records the yearly number of dam incidents in a region, modelled as
$X_1,\dots,X_n$ i.i.d.\ Poisson$(\theta)$; $T=\sum_iX_i$, $z=z_{1-\alpha/2}$.
\begin{enumerate}[label=\textbf{(\alph*)},leftmargin=2.4em]
\item Derive the Wald interval from the CLT and Slutsky's theorem. Show that it degenerates when
$T=0$ and prove that, for fixed $n$, its coverage tends to $0$ as $\theta\to0$.
\item Derive the score interval by inverting $|\overline X-\theta|\le z\sqrt{\theta/n}$. Prove that
both endpoints are non-negative and find the interval when $T=0$.
\item Prove that $P(T\le t;\mu)=P(G_{t+1}>\mu)$ for $T\sim$ Poisson$(\mu)$ and $G_{t+1}\sim$
Gamma$(t+1,1)$. Use it to write the exact (Garwood) interval with gamma quantiles and prove that its
coverage is at least $1-\alpha$ for every $\theta$.
\item Parameter analysis: sign of the shift of the score interval's centre relative to
$\overline X$; ratio of the score and Wald half-widths and its limit as $n\to\infty$; limits of the
three coverages as $\theta\to0$.
\end{enumerate}
\programacao Write \texttt{intervals(t, n)} returning the three 95\% intervals and
\texttt{coverage(theta, n)} computing their \emph{exact} coverage by summing $P(T=t)$ over the
$t$ whose interval contains $\theta$ (no simulation). Print the coverages for $n=10$ and
$\theta\in\{0.1,1,5\}$, and the three intervals for the coal data 1851--1891 ($n=41$, $T=127$).
\resolespaco
\begin{sol}
\solhead
\begin{notacao}
\item $T\sim$ Poisson$(n\theta)$, $\overline X=T/n$; $\alpha\in(0,1)$, $z=z_{1-\alpha/2}$ ($1.959964$ for 95\%).
\item $F_t(\mu)=P(T\le t;\mu)$ for $T\sim$ Poisson$(\mu)$; $G_a\sim$ Gamma$(a,1)$ with quantile $q_p(a)$.
\item $C(\theta)$: coverage probability $P_\theta(\theta\in\text{interval})$.
\end{notacao}

\passo{(a) Wald}{By the CLT, $\sqrt n(\overline X-\theta)/\sqrt\theta\to N(0,1)$; $\overline X\to\theta$ in
probability, so by Slutsky}
\begin{equation}\label{q8:wpiv}
\frac{\overline X-\theta}{\sqrt{\overline X/n}}\xrightarrow{d}N(0,1),
\qquad\text{interval}\quad \overline X\pm z\sqrt{\overline X/n}.
\end{equation}
If $T=0$, \eqref{q8:wpiv} gives $[0,0]$, which excludes every $\theta>0$. Hence
\begin{equation}\label{q8:wcov}
C_{W}(\theta)\le1-P_\theta(T=0)=1-e^{-n\theta}\xrightarrow[\theta\to0]{}0 .
\end{equation}

\passo{(b) Score}{Squaring the inequality, $(\overline X-\theta)^2\le z^2\theta/n$; expanding:}
\begin{equation}\label{q8:quad}
\theta^2-\Big(2\overline X+\frac{z^2}{n}\Big)\theta+\overline X^2\le0 .
\end{equation}
The roots of \eqref{q8:quad} are
\begin{equation}\label{q8:score}
\theta_\pm=\overline X+\frac{z^2}{2n}\pm z\sqrt{\frac{\overline X}{n}+\frac{z^2}{4n^2}} .
\end{equation}
Their product is $\overline X^2\ge0$ and their sum $2\overline X+z^2/n>0$, so neither root is
negative. With $T=0$, \eqref{q8:score} gives $[0,\,z^2/n]$.

\passo{(c) The Poisson--gamma identity}{Differentiate $F_t(\mu)=\sum_{j=0}^te^{-\mu}\mu^j/j!$:}
\[
F_t'(\mu)=\sum_{j=0}^t\Big(-\frac{e^{-\mu}\mu^j}{j!}\Big)+\sum_{j=1}^t\frac{e^{-\mu}\mu^{j-1}}{(j-1)!}.
\]
The second sum is the first without its last term, so the sum telescopes:
\begin{equation}\label{q8:dF}
F_t'(\mu)=-\frac{e^{-\mu}\mu^t}{t!}<0 .
\end{equation}
Since $F_t(\mu)\to0$ as $\mu\to\infty$, integrating \eqref{q8:dF} from $\mu$ to $\infty$:
\begin{equation}\label{q8:id}
F_t(\mu)=\int_\mu^\infty\frac{u^te^{-u}}{t!}\,du=P(G_{t+1}>\mu). \qquad\blacksquare
\end{equation}

\passo{(c) Garwood limits}{The upper limit solves $F_T(n\theta_U)=\alpha/2$; by \eqref{q8:id},
$P(G_{T+1}>n\theta_U)=\alpha/2$. The lower limit solves $P(T'\ge T;n\theta_L)=\alpha/2$, i.e.\
$1-F_{T-1}(n\theta_L)=\alpha/2$, and by \eqref{q8:id} $P(G_T\le n\theta_L)=\alpha/2$. So}
\begin{equation}\label{q8:gar}
\theta_L=\frac{q_{\alpha/2}(T)}{n}\ (\theta_L=0\text{ if }T=0),\qquad \theta_U=\frac{q_{1-\alpha/2}(T+1)}{n}.
\end{equation}

\passo{(c) Coverage}{Fix $\theta$ and let $t^*$ be the largest $t$ with $F_t(n\theta)\le\alpha/2$.
By \eqref{q8:dF}, $F_t$ is decreasing in $\mu$, so $\theta>\theta_U(t)\iff F_t(n\theta)<F_t(n\theta_U(t))=\alpha/2$;
and $F_t(n\theta)$ increases in $t$, so this happens only for $t\le t^*$. Hence
$P_\theta(\theta>\theta_U)\le P_\theta(T\le t^*)=F_{t^*}(n\theta)\le\alpha/2$. The same argument on the
upper tail gives $P_\theta(\theta<\theta_L)\le\alpha/2$, so $C_{G}(\theta)\ge1-\alpha$ for every
$\theta$. $\blacksquare$}

\passo{(d) Parameter analysis}{By \eqref{q8:score} the score interval is centred at
$\overline X+z^2/(2n)$: shifted \emph{up} (sign $+$), towards the side where the Poisson is
skewed, and the shift vanishes at rate $1/n$. The half-width ratio is}
\begin{equation}\label{q8:ratio}
\frac{z\sqrt{\overline X/n+z^2/(4n^2)}}{z\sqrt{\overline X/n}}=\sqrt{1+\frac{z^2}{4n\overline X}}>1,
\qquad\to1\ (n\to\infty).
\end{equation}
As $\theta\to0$: $C_W\to0$ by \eqref{q8:wcov}; $C_S\ge P(T=0)=e^{-n\theta}\to1$, because the $T=0$
interval $[0,z^2/n]$ contains every $\theta\le z^2/n$; and $C_G\ge0.95$ always, tending to 1 since
the $T=0$ interval $[0,q_{1-\alpha/2}(1)/n]$ covers small $\theta$. Reading: for rare events the
Wald interval is unusable, the score interval is a good default, and the exact interval is the safe
(conservative) choice.

\codehead
{\small
\begin{verbatim}
def intervals(t, n, z=1.959964, al=0.05):
    xb = t / n
    wald = (xb - z*np.sqrt(xb/n), xb + z*np.sqrt(xb/n))
    c, h = xb + z**2/(2*n), z*np.sqrt(xb/n + z**2/(4*n**2))
    exact = (0.0 if t == 0 else gamma.ppf(al/2, t)/n, gamma.ppf(1-al/2, t+1)/n)
    return wald, (c - h, c + h), exact
# exact coverage, n = 10:   theta   Wald      score     exact
#                           0.1     0.631526  0.919699  0.981012
#                           1       0.925728  0.962622  0.975386
#                           5       0.948781  0.944398  0.952638
# coal 1851-1891 (n=41, T=127): Wald [2.558837, 3.636285]
#     score [2.603651, 3.685165]  exact [2.582297, 3.685510]
\end{verbatim}
}
At $\theta=0.1$ the Wald coverage $0.6315$ is close to the bound \eqref{q8:wcov},
$1-e^{-1}=0.6321$; the exact interval is above $0.95$ at every $\theta$, as proved.

\intuicao{Wald uses the estimated variance $\overline X/n$, which is zero exactly when the data are
least informative; the score interval uses the variance at the hypothesized $\theta$ and never
collapses. Inverting exact tail probabilities trades length for a guaranteed coverage.}
\refbib{Bolfarine \& Sandoval (2010), §5.2 (quantidades pivotais), §5.3 (intervalos aproximados);
Class~1 deck, \textit{Poisson} (asymptotic interval).}
\rubrica{
\item (a) pivot \eqref{q8:wpiv} with CLT and Slutsky named: 2 pts; \eqref{q8:wcov}: 2 pts.
\item (b) \eqref{q8:quad}--\eqref{q8:score}: 2 pts; non-negativity proved: 1 pt; $T=0$ case: 1 pt.
\item (c) telescoping \eqref{q8:dF} and identity \eqref{q8:id}: 3 pts; limits \eqref{q8:gar}: 2 pts; coverage proof: 2 pts.
\item (d) shift sign, \eqref{q8:ratio} and the three limits with reading: 3 pts (limits without proof: max.\ half).
\item Programming: exact coverage by enumeration (not simulation) and coal intervals: 2 pts.
}
\end{sol}

% -----------------------------------------------------------------------------
\questao{Original conceptual question --- tests and intervals for counts (10 points).}
\fonte{Original question (not from a book).}

\textbf{a)} (3 points) To test $H_0:\theta=5$ against $H_1:\theta\ne5$ with a Poisson sample (B,
Exemplo~6.5.5), one uses $-2\log\lambda=2n[\overline x\log(\overline x/5)-(\overline x-5)]$. For
large $n$, under $H_0$ this statistic is approximately:
\begin{alternativas}
\item $\chi^2_n$, one degree of freedom for each observation in the sample.
\item $N(0,1)$, because $\overline X$ is asymptotically normal under $H_0$.
\item $\chi^2_1$, the number of free parameters in $H_1$ minus that in $H_0$.
\item $\chi^2_2$, because the alternative $\theta\ne5$ has two sides.
\end{alternativas}
\gabarito{C}
\begin{sol}
\passo{Why}{C: Wilks (B, Teorema 6.5.1), $1-0=1$ degree of freedom. A confuses sample size with
dimension; B: the signed root is $N(0,1)$, the statistic itself is its square; D: two-sidedness is
already built into squaring, not into extra degrees of freedom.}
\end{sol}

\textbf{b)} (3 points) Over $n=12$ years a dam recorded no incident ($\overline x=0$). The 95\%
Wald interval for the yearly rate $\theta$ is:
\begin{alternativas}
\item $[0,0]$, a degenerate interval of zero length.
\item $[0,\,3.84/12]$, from inverting the score statistic.
\item $[0,\,1.96/12]$, from the normal quantile over $n$.
\item $[-0.16,\,0.16]$, from a symmetric normal band.
\end{alternativas}
\gabarito{A}
\begin{sol}
\passo{Why}{A: by \eqref{q8:wpiv}, $0\pm1.96\sqrt{0/12}$. B is the \emph{score} interval
$[0,z^2/n]=[0,0.320122]$ of \eqref{q8:score}, a correct answer to a different question; C and D
have no derivation behind them.}
\end{sol}

\textbf{c)} (2 points, T or F) For Poisson data, a UMP test of $H_0:\theta\le\theta_0$ against
$H_1:\theta>\theta_0$ exists because the likelihood ratio is monotone in $T=\sum X_i$.
\gabarito{T}
\begin{sol}
\passo{Why}{\eqref{q6:lr} is increasing in $t$ for $\theta_1>\theta_0$; the test \eqref{q6:phi} does
not depend on $\theta_1$ (B §6.3).}
\end{sol}

\textbf{d)} (2 points, T or F) Because a UMP test exists for the one-sided alternative, a UMP test
also exists for $H_0:\theta=\theta_0$ against $H_1:\theta\ne\theta_0$.
\gabarito{F}
\begin{sol}
\passo{Why}{The most powerful test against $\theta_1>\theta_0$ rejects for large $T$, the one
against $\theta_1<\theta_0$ for small $T$; no single test is most powerful on both sides. Two-sided
problems use the GLR test of (a) instead.}
\refbib{B §6.3--6.5.}
\end{sol}

\programacao Check (a) by simulation. Draw $100{,}000$ samples of size $n=50$ from
Poisson$(5)$, through their totals $T\sim$ Poisson$(250)$, using the generator
\texttt{np.random.default\_rng(9)}; compute the statistic of (a) for each and report the fraction
above $\chi^2_{1,0.95}=3.841$.
\begin{sol}
\codehead
{\small
\begin{verbatim}
rng = np.random.default_rng(9); xb = rng.poisson(250, 100_000) / 50
lr = 2 * 50 * (xb * np.log(xb / 5) - (xb - 5))
np.mean(lr > 3.841459)                          # 0.04866
\end{verbatim}
}
The rejection rate $0.0487$ is close to $0.05$: with $n\theta_0=250$ expected events the $\chi^2_1$
approximation is accurate (compare Question~7, where 3 to 5 events gave about 7\%).
\end{sol}

% =============================================================================
\bloco{3}{Two Poisson regimes: equal rates and the change point}{G 3.8, 7.3; original}

\questao{G 3.8 --- Bayes factor for two typists.}
\fonte{Greenberg, E. (2008), \textit{Introduction to Bayesian Econometrics}, 1st ed., Cambridge
University Press, Chapter~3, §3.5 Exercises, Exercise~3.8, p.~40 [PDF p.~55]; the book gives no
answer --- reproduced verbatim. Table~3.2, to which it refers, is reproduced from §3.2.4, p.~38
[PDF p.~53].}

3.8 Suppose the number of typographical errors per page has a Poisson distribution $P(\theta)$. Sam
types the first $m$ pages of a manuscript and makes $e_1$ errors in total, and Levi types the last
$m$ pages and makes $e_2$ errors. Let Sam's error rate be $\theta_1$ and Levi's error rate be
$\theta_2$. To compare model $M_1:\theta_1=\theta_2=\theta$ with model $M_2:\theta_1\neq\theta_2$,
produce a table like Table~3.2 for $m=100$ with 10 total errors and for $m=200$ with 20 total errors
and Sam's error proportions of 0.9, 0.8, 0.7, 0.6, and 0.5. Take $G(1,1)$ as a prior for $\theta_1$
and $\theta_2$.

\begin{center}\small
\emph{Referred to (Greenberg 2008, p.~38):}\quad \textbf{Table 3.2.} Bayes Factors for Selected Possible Outcomes\\[2pt]
\begin{tabular}{ccrr}
\toprule
Michaela & Lila & \multicolumn{2}{c}{log$_{10}$(Bayes factor)}\\
Proportion heads & Proportion heads & $m=10$ & $m=100$\\
\midrule
0.1 & 0.9 & $-2.506$ & $-30.775$\\
0.2 & 0.8 & $-1.200$ & $-15.793$\\
0.3 & 0.7 & $-0.348$ & $-6.316$\\
0.4 & 0.6 & $0.138$ & $-0.975$\\
0.5 & 0.5 & $0.297$ & $0.756$\\
\bottomrule
\end{tabular}
\end{center}
\programacao Compute the table from the closed form. Then check the cell $m=100$, $(e_1,e_2)=(9,1)$
by Monte Carlo integration over the prior: with \texttt{np.random.default\_rng(10)} draw
$10^6$ values from $G(1,1)$ for each of the three integrals ($M_1$; $M_2$'s two factors), average
the integrands, and report $\log_{10}\hat B_{12}$ with its delta-method standard error.
\resolespaco
\begin{sol}
\solhead
\begin{notacao}
\item $y_{1i}$, $y_{2i}$ ($i=1,\ldots,m$): errors on page $i$ typed by Sam and by Levi; $e_1=\sum_iy_{1i}$, $e_2=\sum_iy_{2i}$, $e=e_1+e_2$.
\item $\theta,\theta_1,\theta_2>0$: error rates per page; prior $G(1,1)$, i.e.\ $\pi(\theta)=e^{-\theta}$.
\item $P=\prod_iy_{1i}!\prod_iy_{2i}!$; $m_j(y)$: marginal likelihood of $M_j$; $B_{12}=m_1(y)/m_2(y)$.
\end{notacao}

\passo{The gamma integral}{For $a,b>0$ (Greenberg App.~A.1.7),}
\begin{equation}\label{q10:gi}
\int_0^\infty\theta^{a-1}e^{-b\theta}d\theta=\frac{\Gamma(a)}{b^{a}}.
\end{equation}

\passo{Marginal likelihood of $M_1$}{All $2m$ pages share $\theta$:}
\begin{align}
m_1(y)&=\int_0^\infty\prod_{i}\frac{\theta^{y_{1i}}e^{-\theta}}{y_{1i}!}\prod_{i}\frac{\theta^{y_{2i}}e^{-\theta}}{y_{2i}!}\,e^{-\theta}\,d\theta\label{q10:m1a}\\
&=\frac1P\int_0^\infty\theta^{e}e^{-(2m+1)\theta}d\theta.\label{q10:m1b}
\end{align}
Substituting \eqref{q10:gi} with $a=e+1$, $b=2m+1$:
\begin{equation}\label{q10:m1}
m_1(y)=\frac1P\,\frac{\Gamma(e+1)}{(2m+1)^{e+1}}.
\end{equation}

\passo{Marginal likelihood of $M_2$}{The two rates have independent priors, so the integral factorises:}
\begin{align}
m_2(y)&=\frac1P\int_0^\infty\theta_1^{e_1}e^{-(m+1)\theta_1}d\theta_1\int_0^\infty\theta_2^{e_2}e^{-(m+1)\theta_2}d\theta_2\label{q10:m2a}\\
&=\frac1P\,\frac{\Gamma(e_1+1)}{(m+1)^{e_1+1}}\,\frac{\Gamma(e_2+1)}{(m+1)^{e_2+1}}.\label{q10:m2}
\end{align}

\passo{Bayes factor}{Dividing \eqref{q10:m1} by \eqref{q10:m2}; $P$ cancels, so only the totals matter:}
\begin{align}
B_{12}&=\frac{\Gamma(e+1)\,(m+1)^{e+2}}{(2m+1)^{e+1}\,\Gamma(e_1+1)\,\Gamma(e_2+1)},\label{q10:B}\\
\log_{10}B_{12}&=\frac{\ln\Gamma(e+1)+(e+2)\ln(m+1)-(e+1)\ln(2m+1)-\ln\Gamma(e_1+1)-\ln\Gamma(e_2+1)}{\ln10}.\label{q10:logB}
\end{align}

\passo{One entry worked out}{$m=100$, $e=10$, Sam's proportion $0.9$, so $e_1=9$, $e_2=1$. The
pieces of \eqref{q10:logB}: $\ln\Gamma(11)=\ln10!=15.104413$; $12\ln101=12(4.615121)=55.381446$;
$11\ln201=11(5.303305)=58.336354$; $\ln\Gamma(10)=\ln9!=12.801827$; $\ln\Gamma(2)=0$. Substituting:}
\begin{align}
\log_{10}B_{12}&=\frac{(15.104413)+(55.381446)-(58.336354)-(12.801827)-(0)}{2.302585}\label{q10:w1}\\
&=\frac{-0.652323}{2.302585}=-0.283.\label{q10:w2}
\end{align}

\passo{The table}{The same computation for every cell:}
\begin{center}
\begin{tabular}{ccrcr}
\toprule
& \multicolumn{2}{c}{$m=100$, $e=10$} & \multicolumn{2}{c}{$m=200$, $e=20$}\\
\cmidrule(lr){2-3}\cmidrule(lr){4-5}
Sam's proportion & $(e_1,e_2)$ & $\log_{10}B_{12}$ & $(e_1,e_2)$ & $\log_{10}B_{12}$\\
\midrule
0.9 & (9, 1) & $-0.283$ & (18, 2) & $-1.717$\\
0.8 & (8, 2) & $0.370$ & (16, 4) & $-0.310$\\
0.7 & (7, 3) & $0.796$ & (14, 6) & $0.593$\\
0.6 & (6, 4) & $1.039$ & (12, 8) & $1.105$\\
0.5 & (5, 5) & $1.118$ & (10, 10) & $1.271$\\
\bottomrule
\end{tabular}
\end{center}

\passo{Reading it}{On Jeffreys' scale (Greenberg Table~3.1), with 10 errors even a 9--1 split
favours $M_2$ only weakly ($\log_{10}B_{21}=0.283<1/2$), while an even split gives strong support to
$M_1$ ($1<1.118<3/2$). Doubling the evidence (20 errors) sharpens everything: an 18--2 split now
gives very strong evidence for $M_2$ ($3/2<1.717<2$), and equal shares strengthen $M_1$ ($1.271$).}

\passo{Monte Carlo check (programming)}{Each marginal likelihood in \eqref{q10:m1b}--\eqref{q10:m2a}
is a prior expectation, e.g.\ $P\,m_1(y)=E_{\pi}[\theta^{e}e^{-2m\theta}]$ with $\theta\sim G(1,1)$, so it
can be estimated by averaging the integrand over prior draws (Greenberg §5.4 with $h=\pi$). If
$\hat m=\bar d$ averages $G$ draws $d_g$, the delta method gives
$\Var(\ln\hat m)\approx\Var(d)/(G\bar d^{\,2})$, and the three estimates are independent.}

\codehead
{\small
\begin{verbatim}
rng = np.random.default_rng(10); G, m, e1, e2 = 10**6, 100, 9, 1
h = lambda t, k, r: np.exp(k * np.log(t) - r * t)
d1 = h(rng.exponential(1, G), e1 + e2, 2 * m)          # M1 integrand
d2a, d2b = h(rng.exponential(1, G), e1, m), h(rng.exponential(1, G), e2, m)
lB = np.log10(d1.mean() / (d2a.mean() * d2b.mean()))   # -0.28295
cv2 = lambda d: d.var(ddof=1) / (G * d.mean()**2)
se = np.sqrt(cv2(d1) + cv2(d2a) + cv2(d2b)) / np.log(10)   # 0.00312
\end{verbatim}
}
$-0.28295\pm0.00312$ brackets the exact $-0.28330$. Prior sampling works here because the prior
$G(1,1)$ still puts mass near the likelihood's peak ($\theta\approx0.05$); with a more concentrated
likelihood most draws would contribute nothing.

\intuicao{The Bayes factor automatically penalises the extra parameter of $M_2$: it pays off only
when the split is lopsided enough to beat that Occam penalty, and the same proportions carry more
weight when they rest on more errors (compare the two $m$ columns, as in Table~3.2).}
\refbib{Greenberg (2008), §3.2.4 (Bayes factors, Tables 3.1--3.2, pp.~33--38) and §5.4
(Monte Carlo integration).}
\end{sol}

% -----------------------------------------------------------------------------
\questao{G 7.3 --- Gibbs sampler for the coal-mining switch point (computer exercise).}
\fonte{Greenberg, E. (2008), \textit{Introduction to Bayesian Econometrics}, 1st ed., Cambridge
University Press, Chapter~7, §7.5 Exercises, Exercise~7.3, p.~106 [PDF p.~121]; the book gives no
answer --- reproduced verbatim. Equations (7.3) and (7.4), to which the exercise refers, are
reproduced verbatim from §7.1.1, pp.~93--94 [PDF pp.~108--109].}

\emph{Equations referred to (Greenberg 2008, pp.~93--94):}
\[
p(y_i)=\begin{cases}\dfrac{e^{-\theta_1}\theta_1^{y_i}}{y_i!}, & \text{for } i=1,\dots,k,\\[8pt]
\dfrac{e^{-\theta_2}\theta_2^{y_i}}{y_i!}, & \text{for } i=k+1,\dots,n,\end{cases}\tag{7.3}
\]
``where $y_i=0,1,\dots,$ and the switch point $k$ is unknown. The specification is completed by
assigning the conditionally conjugate priors,
\[
\theta_1\sim G(\alpha_{10},\beta_{10}),\quad \theta_2\sim G(\alpha_{20},\beta_{20}),\quad
\pi(k=j)=1/n,\quad j=1,\dots,n.\text{''}\tag{7.4}
\]

7.3 Construct a Gibbs algorithm to analyze the Poisson model with unknown switch point. Given the
specification in Equations (7.3) and (7.4), show that
\[
\pi(\theta_1,\theta_2,k\mid y)\propto\theta_1^{\alpha_{10}-1}e^{-\beta_{10}\theta_1}\theta_2^{\alpha_{20}-1}e^{-\beta_{20}\theta_2}
\prod_{i=1}^{k}e^{-\theta_1}\theta_1^{y_i}\prod_{i=k+1}^{n}e^{-\theta_2}\theta_2^{y_i},
\]
and verify that
\begin{gather*}
\theta_1\mid y,k\sim G\Big(\alpha_{10}+\sum_{1}^{k}y_i,\ \beta_{10}+k\Big),\\
\theta_2\mid k,y\sim G\Big(\alpha_{20}+\sum_{k+1}^{n}y_i,\ \beta_{20}+n-k\Big),\\
\pi(k\mid y,\theta_1,\theta_2)=\frac{e^{k(\theta_2-\theta_1)}(\theta_1/\theta_2)^{\sum_1^ky_i}}
{\sum_{k=1}^{n}e^{k(\theta_2-\theta_1)}\,(\theta_1/\theta_2)^{\sum_1^ky_i}},\qquad k=1,\dots,n.
\end{gather*}
Apply your algorithm to the mining disaster data analyzed by Carlin et al.\ (1992a); the data may be
found in the article. Your analysis should generate marginal posterior distributions for
$\theta_1$, $\theta_2$, $p(k)$, and the marginal likelihood computed by the method of Chib (1995).
For simplicity, set $\alpha_{10}=\alpha_{20}=0.5$ and $\beta_{10}=\beta_{20}=1.0$. (Curiously,
although these data have been analyzed from the Bayesian viewpoint in several papers, none of them
present historical material that might shed light on when changes in parameters took place. For
example, changes in government regulation, technology, or the location of mines might be expected
to affect the probability of a disaster. A more general version of the model permits multiple
switch points if necessary.)
\programacao Beyond the book's tasks (use the coal series of the Instructions,
\texttt{np.random.default\_rng(2026)}, start $k^{(0)}=56$, 1{,}000 burn-in and $G=10{,}000$ kept
draws): (1) compute the \emph{exact} $p(k\mid y)$ by integrating $\theta_1,\theta_2$ out analytically
and report the total-variation distance between it and the Gibbs frequencies of $k$; (2) report the
numerical standard error of the posterior mean of $\theta_1$ by batch means (50 batches of 200),
the i.i.d.\ standard error $s/\sqrt G$, and their squared ratio (inefficiency factor).
\resolespaco
\begin{sol}
\solhead
\begin{notacao}
\item $y_i$: number of British coal-mining disasters in year $i$, $i=1$ (1851) to $n=112$ (1962);
  the series of Carlin, Gelfand \& Smith (1992) is Jarrett (1979), total $S=\sum_iy_i=191$.
\item $k\in\{1,\dots,n\}$: switch point, the last year of regime 1 (year $1850+k$); $k=n$ means no change.
\item $\theta_1,\theta_2>0$: disaster rates before and after the switch;
  $S_k=\sum_{i=1}^ky_i$ (so $\sum_{k+1}^ny_i=S-S_k$).
\item $G(a,b)$: gamma with shape $a$, rate $b$; $\alpha_{10}=\alpha_{20}=\alpha=0.5$,
  $\beta_{10}=\beta_{20}=\beta=1$.
\item $g=1,\dots,G$: Gibbs iteration; $\theta^*=(\theta_1^*,\theta_2^*)$: point where Chib's
  identity is evaluated; $m(y)$: marginal likelihood.
\end{notacao}

\passo{Joint posterior}{By (7.3) and independence, the likelihood is}
\begin{equation}\label{q11:lik}
p(y\mid\theta_1,\theta_2,k)=\prod_{i=1}^k\frac{e^{-\theta_1}\theta_1^{y_i}}{y_i!}
\prod_{i=k+1}^n\frac{e^{-\theta_2}\theta_2^{y_i}}{y_i!},
\end{equation}
and by (7.4) the prior is
\begin{equation}\label{q11:prior}
\pi(\theta_1,\theta_2,k)=\frac{\beta_{10}^{\alpha_{10}}}{\Gamma(\alpha_{10})}\theta_1^{\alpha_{10}-1}e^{-\beta_{10}\theta_1}
\cdot\frac{\beta_{20}^{\alpha_{20}}}{\Gamma(\alpha_{20})}\theta_2^{\alpha_{20}-1}e^{-\beta_{20}\theta_2}\cdot\frac1n .
\end{equation}
Multiplying \eqref{q11:lik} by \eqref{q11:prior} and dropping the factors free of
$(\theta_1,\theta_2,k)$ (the $y_i!$, the gamma constants, $1/n$) gives the expression in the statement:
\begin{equation}\label{q11:joint}
\pi(\theta_1,\theta_2,k\mid y)\propto\theta_1^{\alpha_{10}-1}e^{-\beta_{10}\theta_1}\theta_2^{\alpha_{20}-1}e^{-\beta_{20}\theta_2}
\prod_{i=1}^{k}e^{-\theta_1}\theta_1^{y_i}\prod_{i=k+1}^{n}e^{-\theta_2}\theta_2^{y_i}.
\end{equation}

\passo{Full conditional of $\theta_1$}{Keeping in \eqref{q11:joint} only the factors with $\theta_1$:}
\[
\pi(\theta_1\mid y,k,\theta_2)\propto\theta_1^{\alpha_{10}-1}e^{-\beta_{10}\theta_1}\big(e^{-k\theta_1}\theta_1^{S_k}\big).
\]
Adding exponents:
\begin{equation}\label{q11:t1}
\pi(\theta_1\mid y,k)\propto\theta_1^{\alpha_{10}+S_k-1}e^{-(\beta_{10}+k)\theta_1}
\;\Longrightarrow\;\theta_1\mid y,k\sim G(\alpha_{10}+S_k,\ \beta_{10}+k).
\end{equation}

\passo{Full conditional of $\theta_2$}{In the same way, the $\theta_2$ factors of \eqref{q11:joint}
are}
\[
\pi(\theta_2\mid y,k,\theta_1)\propto\theta_2^{\alpha_{20}-1}e^{-\beta_{20}\theta_2}\big(e^{-(n-k)\theta_2}\theta_2^{S-S_k}\big),
\]
so, adding exponents,
\begin{equation}\label{q11:t2}
\theta_2\mid y,k\sim G(\alpha_{20}+S-S_k,\ \beta_{20}+n-k).
\end{equation}
Neither \eqref{q11:t1} nor \eqref{q11:t2} involves the other rate: given $k$, $\theta_1$ and $\theta_2$
are independent.

\passo{Full conditional of $k$}{The factors of \eqref{q11:joint} that change with $k$ are the two
products:}
\begin{equation}\label{q11:k0}
\pi(k\mid y,\theta_1,\theta_2)\propto e^{-k\theta_1}\theta_1^{S_k}\,e^{-(n-k)\theta_2}\theta_2^{S-S_k}.
\end{equation}
Dividing \eqref{q11:k0} by the $k$-free factor $e^{-n\theta_2}\theta_2^{S}$:
\begin{equation}\label{q11:k}
\pi(k\mid y,\theta_1,\theta_2)\propto e^{k(\theta_2-\theta_1)}\Big(\frac{\theta_1}{\theta_2}\Big)^{S_k},
\qquad k=1,\dots,n,
\end{equation}
and normalising over $k=1,\dots,n$ gives the formula in the statement.

\passo{Algorithm}{Start at $k^{(0)}=56$. For $g=1,\dots,11{,}000$: draw $\theta_1^{(g)}$ from
\eqref{q11:t1} and $\theta_2^{(g)}$ from \eqref{q11:t2} given $k^{(g-1)}$; compute the $n$ weights
\eqref{q11:k} on the log scale, normalise, and draw $k^{(g)}$. Discard 1{,}000 burn-in draws and keep
$G=10{,}000$ (NumPy seed 2026).}

\passo{Exact benchmark}{Because the priors are conjugate given $k$, $\theta_1,\theta_2$ integrate out
of \eqref{q11:lik}$\times$\eqref{q11:prior} in closed form, by \eqref{q10:gi}:}
\begin{equation}\label{q11:pyk}
p(y\mid k)=\frac{1}{\prod_iy_i!}\cdot\frac{\beta^\alpha\,\Gamma(\alpha+S_k)}{\Gamma(\alpha)\,(\beta+k)^{\alpha+S_k}}
\cdot\frac{\beta^\alpha\,\Gamma(\alpha+S-S_k)}{\Gamma(\alpha)\,(\beta+n-k)^{\alpha+S-S_k}} ,
\end{equation}
so the exact marginal likelihood and posterior of $k$ are
\begin{equation}\label{q11:exact}
m(y)=\sum_{k=1}^n\frac1n\,p(y\mid k),\qquad p(k\mid y)=\frac{p(y\mid k)/n}{m(y)},
\end{equation}
giving $\log m(y)=-179.4590$, $\E(\theta_1\mid y)=3.0500$, $\E(\theta_2\mid y)=0.9147$. These check the
sampler.

\passo{Marginal posteriors}{Gibbs output (exact values from \eqref{q11:exact} for $k$):}
\begin{center}\small
\begin{tabular}{lccc}
\toprule
 & Posterior mean & Posterior sd & 95\% interval \\
\midrule
$\theta_1$ (before) & 3.0506 & 0.2809 & $[2.530,\ 3.635]$\\
$\theta_2$ (after)  & 0.9166 & 0.1163 & $[0.701,\ 1.156]$\\
\bottomrule
\end{tabular}
\qquad
\begin{tabular}{cccc}
\toprule
$k$ & year & $p(k\mid y)$ Gibbs & exact \\
\midrule
41 & 1891 & 0.2471 & 0.2482\\
40 & 1890 & 0.1849 & 0.1848\\
39 & 1889 & 0.1433 & 0.1414\\
42 & 1892 & 0.0987 & 0.1010\\
37 & 1887 & 0.0881 & 0.0898\\
36 & 1886 & 0.0742 & 0.0745\\
\bottomrule
\end{tabular}
\end{center}
The posterior of $k$ has mean $40.14$ and mode $k=41$ (1891); $P(36\le k\le46\mid y)=0.98$, i.e.\ the
switch happened between 1886 and 1896 with high probability. Disasters fell from about 3.05 to
about 0.92 per year.

\passo{Marginal likelihood by Chib (1995)}{For any $\theta^*$, Bayes' theorem rearranged gives the
identity}
\begin{equation}\label{q11:chib}
\log m(y)=\log p(y\mid\theta^*)+\log\pi(\theta^*)-\log\pi(\theta^*\mid y).
\end{equation}
Here $k$ is a latent parameter, so the likelihood term sums it out with its prior:
\begin{equation}\label{q11:lstar}
p(y\mid\theta^*)=\sum_{k=1}^n\frac1n\,p(y\mid\theta_1^*,\theta_2^*,k).
\end{equation}
For the ordinate, conditional independence of $\theta_1,\theta_2$ given $k$ (from \eqref{q11:t1} and
\eqref{q11:t2}) gives the Rao--Blackwell estimate
\begin{equation}\label{q11:ord}
\hat\pi(\theta^*\mid y)=\frac1G\sum_{g=1}^G G\big(\theta_1^*\mid\alpha+S_{k^{(g)}},\beta+k^{(g)}\big)\,
G\big(\theta_2^*\mid\alpha+S-S_{k^{(g)}},\beta+n-k^{(g)}\big),
\end{equation}
with $G(\cdot\mid a,b)$ the gamma density. At the posterior means $\theta^*=(3.0506,\ 0.9166)$:
$\log p(y\mid\theta^*)=-172.2563$ from \eqref{q11:lstar}, $\log\pi(\theta^*)=-5.6261$ from (7.4), and
$\log\hat\pi(\theta^*\mid y)=1.5765$ from \eqref{q11:ord}. Substituting into \eqref{q11:chib}:
\[
\log\hat m(y)=(-172.2563)+(-5.6261)-(1.5765)=-179.4589 ,
\]
against the exact $-179.4590$ from \eqref{q11:exact}.

\passo{Added: exactness and numerical standard error}{The total-variation distance
$\frac12\sum_k|\hat p(k)-p(k\mid y)|$ between the Gibbs frequencies and \eqref{q11:exact} is
$0.0077$. For $\E(\theta_1\mid y)$, batch means give n.s.e.\ $=0.003035$ against the i.i.d.\ value
$0.2809/\sqrt{10{,}000}=0.002809$; the inefficiency factor $(0.003035/0.002809)^2=1.167$ says the
chain is almost as informative as independent draws (Greenberg §7.3).}

\codehead
{\small
\begin{verbatim}
cs, ks = np.cumsum(y), np.arange(1, n + 1)        # S_k and k = 1..n
rng, k = np.random.default_rng(2026), n // 2
for g in range(11000):
    t1 = rng.gamma(a0 + cs[k-1], 1 / (b0 + k))          # (11.4) theta1 | k
    t2 = rng.gamma(a0 + S - cs[k-1], 1 / (b0 + n - k))  # (11.5) theta2 | k
    lp = ks * (t2 - t1) + cs * np.log(t1 / t2)          # (11.7) log pi(k | ...)
    p = np.exp(lp - lp.max()); k = rng.choice(ks, p=p / p.sum())
# theta1 3.0506 (0.2809), theta2 0.9166 (0.1163), mode k = 41 (1891)
# Chib: -172.2563 + (-5.6261) - 1.5765 = -179.4589 ; exact -179.4590
bm = th1.reshape(50, 200).mean(1)
nse = bm.std(ddof=1) / np.sqrt(50)                  # 0.003035 (iid: 0.002809)
\end{verbatim}
}

\intuicao{The data choose a sharp break: about three disasters a year until around 1890, under one
afterwards, with the switch year pinned down within a few years. The Gibbs sampler works because
each block is easy given the others: rates are conjugate gammas given $k$, and $k$ is a discrete
draw over 112 years given the rates. On the historical question Greenberg raises, the posterior
mass near 1887--1891 sits close to the Coal Mines Regulation Act of 1887, a candidate explanation
that this model does not test. Chib's identity recovers the exact marginal likelihood to the fourth
decimal.}
\refbib{Greenberg (2008), §7.1.1 (pp.~91--95), §7.3 (n.s.e., batch means) and §7.4 (Chib 1995);
Class~1 deck \textit{Poisson with breakpoint}.}
\end{sol}

% -----------------------------------------------------------------------------
\questao{Original math question --- likelihood ratio for a change point.}
\fonte{Original question (not from a book).}

Counts $y_1,\dots,y_n$ are independent with $y_i\sim$ Poisson$(\theta_1)$ for $i\le\tau$ and
Poisson$(\theta_2)$ for $i>\tau$, $1\le\tau<n$. Let $S=\sum_{i=1}^ny_i>0$ and $S_\tau=\sum_{i\le\tau}y_i$.
\begin{enumerate}[label=\textbf{(\alph*)},leftmargin=2.4em]
\item For known $\tau$, derive the MLEs and the statistic $\Lambda_\tau=-2\log\lambda_\tau$ for
$H_0:\theta_1=\theta_2$. Show that $\Lambda_\tau=2S\,\mathrm{KL}(p\,\|\,\pi)$, the Kullback--Leibler
divergence between Bernoulli$(p)$ and Bernoulli$(\pi)$ with $p=S_\tau/S$ and $\pi=\tau/n$, and
prove that $\Lambda_\tau\ge0$ with equality iff $p=\pi$.
\item Prove that under $H_0$, conditionally on $S$, $S_\tau\sim$ Binomial$(S,\tau/n)$ whatever the
common rate, and use it to build an exact test for known $\tau$.
\item For unknown $\tau$, let $M=\max_{1\le\tau<n}\Lambda_\tau$ (the profile-likelihood statistic) and
$\hat\tau$ its argmax. Prove that, under $H_0$, $P(M>c)\ge P(\Lambda_{\tau_0}>c)$ for any fixed
$\tau_0$ and $c$, and conclude what happens to a test that compares $M$ with the $\chi^2_1$ critical
value.
\item Parameter analysis: for fixed $S$ and $\pi$, prove the sign of $\partial\Lambda_\tau/\partial p$
and the convexity of $\Lambda_\tau$ in $p$; find its limits as $p\to0$ and $p\to1$; say how
$\Lambda_\tau$ scales with $S$ at a fixed $p$. Interpret.
\end{enumerate}
\programacao For the coal series ($n=112$): compute $\Lambda_\tau$ for $\tau=1,\dots,111$, $\hat\tau$
and $M$; the conditional binomial p-value of (b) at $\tau=\hat\tau$ (and why it is not a valid
p-value once $\tau$ was chosen by the data); and the null distribution of $M$ by a parametric
bootstrap (2{,}000 series of 112 i.i.d.\ Poisson$(S/n)$ counts, \texttt{np.random.default\_rng(12)}):
its 95\% quantile, the rejection rate of the rule ``$M>3.841$'' and the bootstrap p-value of the
observed $M$.
\resolespaco
\begin{sol}
\solhead
\begin{notacao}
\item $n_1=\tau$, $n_2=n-\tau$; $\hat\theta_1=S_\tau/\tau$, $\hat\theta_2=(S-S_\tau)/(n-\tau)$, $\hat\theta=S/n$.
\item $p=S_\tau/S\in[0,1]$: share of events before $\tau$; $\pi=\tau/n\in(0,1)$: share of time before $\tau$.
\item $\mathrm{KL}(p\|\pi)=p\log\frac p\pi+(1-p)\log\frac{1-p}{1-\pi}$ (with $0\log0=0$).
\end{notacao}

\passo{(a) MLEs}{The log-likelihood is
$\ell=S_\tau\log\theta_1-\tau\theta_1+(S-S_\tau)\log\theta_2-(n-\tau)\theta_2+\text{const}$; its two
score equations separate and give $\hat\theta_1$, $\hat\theta_2$ of the Notation (as in
\eqref{q14:score}); under $H_0$ the pooled MLE is $\hat\theta=S/n$.}

\passo{(a) The statistic}{As in \eqref{q7:stat} of Question~7, with $(n,m)$ replaced by
$(\tau,n-\tau)$ (the linear terms cancel because $\tau\hat\theta_1+(n-\tau)\hat\theta_2=S=n\hat\theta$):}
\begin{equation}\label{q12:L}
\Lambda_\tau=2\Big[S_\tau\log\frac{\hat\theta_1}{\hat\theta}+(S-S_\tau)\log\frac{\hat\theta_2}{\hat\theta}\Big].
\end{equation}

\passo{(a) KL form}{The ratios in \eqref{q12:L} are}
\begin{equation}\label{q12:ratios}
\frac{\hat\theta_1}{\hat\theta}=\frac{S_\tau/\tau}{S/n}=\frac{S_\tau/S}{\tau/n}=\frac p\pi,\qquad
\frac{\hat\theta_2}{\hat\theta}=\frac{(S-S_\tau)/S}{(n-\tau)/n}=\frac{1-p}{1-\pi}.
\end{equation}
Substituting \eqref{q12:ratios} and $S_\tau=Sp$ into \eqref{q12:L}:
\begin{equation}\label{q12:KL}
\Lambda_\tau=2\Big[(Sp)\log\frac p\pi+(S(1-p))\log\frac{1-p}{1-\pi}\Big]=2S\,\mathrm{KL}(p\|\pi).
\end{equation}

\passo{(a) Non-negativity}{$-\log$ is strictly convex, so by Jensen's inequality over the two-point
law $(p,1-p)$:}
\[
\mathrm{KL}(p\|\pi)=-\Big[p\log\frac\pi p+(1-p)\log\frac{1-\pi}{1-p}\Big]\ge-\log\Big[p\frac\pi p+(1-p)\frac{1-\pi}{1-p}\Big]=-\log1=0,
\]
with equality iff $\pi/p=(1-\pi)/(1-p)$, i.e.\ $p=\pi$. By \eqref{q12:KL}, $\Lambda_\tau\ge0$, zero iff
the share of events equals the share of time.

\passo{(b) Conditional law}{Under $H_0$ with common rate $\theta$, $S_\tau\sim$ Poisson$(\tau\theta)$
and $S-S_\tau\sim$ Poisson$((n-\tau)\theta)$ independently, so for $0\le s\le S$:}
\[
P(S_\tau=s\mid S)=\frac{\frac{(\tau\theta)^se^{-\tau\theta}}{s!}\frac{((n-\tau)\theta)^{S-s}e^{-(n-\tau)\theta}}{(S-s)!}}{\frac{(n\theta)^Se^{-n\theta}}{S!}}
=\binom Ss\Big(\frac\tau n\Big)^s\Big(1-\frac\tau n\Big)^{S-s},
\]
the same cancellation as \eqref{q4:bin}. An exact level-$\alpha$ test for known $\tau$ rejects when
$S_\tau$ is in the tails of Binomial$(S,\pi)$; its size is at most $\alpha$ for every $\theta$
because the conditional law does not involve $\theta$.

\passo{(c) The maximum dominates}{For every sample, $M=\max_\tau\Lambda_\tau\ge\Lambda_{\tau_0}$.
Hence $\{\Lambda_{\tau_0}>c\}\subseteq\{M>c\}$ and}
\begin{equation}\label{q12:dom}
P_{H_0}(M>c)\ge P_{H_0}(\Lambda_{\tau_0}>c).
\end{equation}
For fixed $\tau_0$, $\Lambda_{\tau_0}\to\chi^2_1$ under $H_0$ (one restriction; B Teorema 6.5.1), so
with $c=3.841$ the right side of \eqref{q12:dom} tends to $0.05$: comparing $M$ with $\chi^2_1$
gives a test whose size is at least 5\% and, since the maximum runs over many correlated
statistics, in practice far larger. Searching for the break and then testing it as if it were
known is a multiple-testing error; the null law of $M$ must be simulated.

\passo{(d) Signs and convexity}{Differentiating \eqref{q12:KL} in $p$:}
\begin{equation}\label{q12:dp}
\frac{\partial\Lambda_\tau}{\partial p}=2S\Big[\log\frac p\pi+1-\log\frac{1-p}{1-\pi}-1\Big]
=2S\log\frac{p(1-\pi)}{(1-p)\pi},
\end{equation}
negative for $p<\pi$, zero at $p=\pi$, positive for $p>\pi$ (the argument of the log exceeds 1
iff $p>\pi$). Differentiating \eqref{q12:dp} again:
\[
\frac{\partial^2\Lambda_\tau}{\partial p^2}=2S\Big(\frac1p+\frac1{1-p}\Big)=\frac{2S}{p(1-p)}>0 .
\]
Limits from \eqref{q12:KL}: $p\to0$ gives $2S\log\frac1{1-\pi}$; $p\to1$ gives $2S\log\frac1\pi$.
At fixed $p$, $\Lambda_\tau$ is linear in $S$.

\passo{(d) Reading}{Evidence of a break grows with the gap between the share of disasters and
the share of years before $\tau$, symmetrically in the KL sense, and doubles when the number of
events doubles at the same shares: a 66\%--37\% split is weak with 10 events and overwhelming with
191.}

\passo{Coal data (programming)}{The maximum is at $\hat\tau=41$ (1891): $S_{41}=127$, $p=127/191$,
$\pi=41/112$; $\hat\theta_1=127/41=3.097561$, $\hat\theta_2=64/71=0.901408$, $\hat\theta=191/112=1.705357$,
and $M=\Lambda_{41}=69.99$. Under $H_0$, $E(S_{41}\mid S)=191\cdot41/112=69.92$, while 127 were
observed: the binomial tail $P(S_{41}\ge127)=5.4\times10^{-17}$. That number is not a valid p-value,
since $\tau=41$ was chosen because it maximizes the evidence. The bootstrap null law of $M$ has 95\%
quantile $10.08$ (not $3.84$); the rule $M>3.841$ rejects in $56.75\%$ of null series, which
quantifies (c); none of the 2{,}000 null values reaches $69.99$ (bootstrap p-value $<1/2000$).}

\codehead
{\small
\begin{verbatim}
def lr_tau(y):
    n, S = len(y), y.sum(); tau = np.arange(1, n); s1 = np.cumsum(y)[:-1]
    xlogy = lambda a, b: np.where(a > 0, a * np.log(np.where(a > 0, a, 1) / b), 0.)
    return 2 * (xlogy(s1, tau * S / n) + xlogy(S - s1, (n - tau) * S / n))
L = lr_tau(y); L.argmax() + 1, L.max()          # 41, 69.99
binom.sf(126, 191, 41 / 112)                    # 5.4e-17
rng = np.random.default_rng(12)
M = np.array([lr_tau(rng.poisson(191 / 112, 112)).max() for _ in range(2000)])
np.quantile(M, .95), np.mean(M > 3.841), np.mean(M >= 69.99)   # 10.08 0.5675 0.0
\end{verbatim}
}

\intuicao{A change point turns a count series into a split of events between two periods; the
likelihood ratio measures how far that split is from the split of time. Choosing the period that
looks most different and then testing it as if fixed inflates the evidence, the same trap as
reporting the best of many specifications.}
\refbib{Class~1 deck \textit{Poisson with breakpoint} (LR at 1891 $=69.99$); Bolfarine \& Sandoval
(2010), §6.5; Greenberg (2008), §7.1.1 for the Bayesian treatment.}
\rubrica{
\item (a) MLEs and \eqref{q12:L}: 3 pts; \eqref{q12:ratios}--\eqref{q12:KL}: 2 pts; Jensen proof: 2 pts.
\item (b) conditional binomial derived and the exact test: 3 pts.
\item (c) \eqref{q12:dom} and the conclusion about the $\chi^2_1$ rule: 3 pts (claim without the inclusion argument: max.\ 1).
\item (d) \eqref{q12:dp} with sign, convexity, limits, scaling and reading: 4 pts.
\item Programming: $\hat\tau$, $M$, binomial tail and bootstrap quantile: 3 pts.
}
\end{sol}

% -----------------------------------------------------------------------------
\questao{Original conceptual question --- switch points and Bayes factors (10 points).}
\fonte{Original question (not from a book).}

\textbf{a)} (3 points) In the Gibbs sampler of Question~11, the full conditional distribution of the
switch point $k$ is:
\begin{alternativas}
\item uniform on $\{1,\dots,n\}$, its prior, because $k$ does not appear in the gamma conditionals.
\item a gamma distribution, because the gamma prior is conjugate to the Poisson likelihood.
\item Binomial$(S,k/n)$, the conditional law of the first-segment count given the total.
\item discrete on $\{1,\dots,n\}$ with weights proportional to $e^{k(\theta_2-\theta_1)}(\theta_1/\theta_2)^{S_k}$.
\end{alternativas}
\gabarito{D}
\begin{sol}
\passo{Why}{D is \eqref{q11:k}. A: $k$ enters the likelihood through which years use which rate,
so the data update it. B: $k$ is discrete; conjugacy concerns the rates. C confuses the
conditional law of the \emph{data} $S_k$ under $H_0$ (Question~12(b)) with the posterior of the
\emph{parameter} $k$.}
\end{sol}

\textbf{b)} (3 points) The Bayes factor $B_{12}$ of Question~10 depends on the page-level data only
through:
\begin{alternativas}
\item the individual page counts, through the product of factorials $\prod_iy_{1i}!\prod_iy_{2i}!$.
\item the totals $e_1$ and $e_2$ (given $m$), because the factorials cancel in the ratio.
\item the grand total $e=e_1+e_2$, because model $M_1$ pools the two typists' pages.
\item the proportion $e_1/e$ alone, whatever the number of errors behind it.
\end{alternativas}
\gabarito{B}
\begin{sol}
\passo{Why}{B: \eqref{q10:B}. A: $P$ appears in both marginal likelihoods and cancels. C: $M_2$ uses
$e_1$ and $e_2$ separately. D: the table shows that the same proportion gives different factors
for $e=10$ and $e=20$.}
\end{sol}

\textbf{c)} (2 points, T or F) In the switch-point model with independent gamma priors, $\theta_1$
and $\theta_2$ are independent a posteriori \emph{given} $k$.
\gabarito{T}
\begin{sol}
\passo{Why}{\eqref{q11:t1} does not involve $\theta_2$ and \eqref{q11:t2} does not involve
$\theta_1$; the joint \eqref{q11:joint} factorizes for fixed $k$.}
\end{sol}

\textbf{d)} (2 points, T or F) Hence $\theta_1$ and $\theta_2$ are also independent in their
marginal posterior, after $k$ is integrated out.
\gabarito{F}
\begin{sol}
\passo{Why}{The marginal posterior is a mixture over $k$ of product densities, and mixing products
creates dependence. Here a later $k$ lowers both conditional means ($k=36$: $3.18$ and $0.97$;
$k=44$: $2.90$ and $0.89$, from \eqref{q11:t1}--\eqref{q11:t2}), so averaging over $p(k\mid y)$
gives an exact posterior correlation of $0.060$, not $0$. Conditional independence does not imply
marginal independence.}
\end{sol}

\programacao Using the Gibbs draws of Question~11, compute the sample correlation of
$(\theta_1^{(g)},\theta_2^{(g)})$; compute the exact posterior correlation from the mixture over
$p(k\mid y)$ of the moments of the two conditional gammas $\theta_j\mid y,k$ given in the statement
of Question~11.
\begin{sol}
\codehead
{\small
\begin{verbatim}
np.corrcoef(th1, th2)[0, 1]                     # 0.0708 (Gibbs)
m1, m2 = (a0 + cs) / (b0 + ks), (a0 + S - cs) / (b0 + n - ks)     # E(theta_j | k)
v1 = np.sum(pk * ((a0 + cs) / (b0 + ks)**2 + m1**2)) - np.sum(pk * m1)**2
v2 = np.sum(pk * ((a0 + S - cs) / (b0 + n - ks)**2 + m2**2)) - np.sum(pk * m2)**2
(np.sum(pk * m1 * m2) - np.sum(pk * m1) * np.sum(pk * m2)) / np.sqrt(v1 * v2)   # 0.0600
\end{verbatim}
}
Both are clearly positive (the Gibbs value carries Monte Carlo error of about $0.01$), so (d) is false.
\end{sol}

% =============================================================================
\bloco{4}{Bayesian gamma--Poisson: posterior, predictive, Bayes estimators}{G 2.2, 3.2, 3.5; B 4.9/5.12; original}

\questao{G 2.2 --- gamma--Poisson conjugacy.}
\fonte{Greenberg, E. (2008), \textit{Introduction to Bayesian Econometrics}, 1st ed., Cambridge
University Press, Chapter~2, §2.5 Exercises, Exercise~2.2, p.~19 [PDF p.~34]; the book gives no
answer --- reproduced verbatim.}

2.2 The Poisson distribution has probability mass function
\[ p(y_i\mid\theta)=\frac{\theta^{y_i}e^{-\theta}}{y_i!},\qquad \theta>0,\ y_i=0,1,\dots, \]
and let $y_1,\dots,y_n$ be a random sample from this distribution.
\begin{enumerate}[label=\textbf{(\alph*)},leftmargin=2.4em]
\item Show that the gamma distribution $G(\alpha,\beta)$ is a conjugate prior distribution for the
Poisson distribution.
\item Show that $\bar y$ is the MLE for $\theta$.
\item Write the mean of the posterior distribution as a weighted average of the mean of the prior
distribution and the MLE.
\item What happens to the weight on the prior mean as $n$ becomes large?
\end{enumerate}
\programacao Use the coal counts 1851--1891 ($n=41$) and the prior $G(0.5,1)$ of G~7.3. Evaluate
the unnormalized posterior $\theta^{\alpha+s-1}e^{-(\beta+n)\theta}$ on a grid of $200{,}001$ points in
$(0,10]$ (work on the log scale), normalize it with the trapezoidal rule, and compare its mean with
the closed form of (c). Print the prior weight $w$ for $n=1,10,41,112,1000$.
\resolespaco
\begin{sol}
\solhead
\begin{notacao}
\item $y_i\in\{0,1,\dots\}$: counts; $\theta>0$: Poisson rate; $s=\sum_{i=1}^ny_i$, $\bar y=s/n$.
\item $G(\alpha,\beta)$, $\alpha,\beta>0$: gamma with shape $\alpha$ and rate $\beta$ (Greenberg's
  convention, App.~A), $\pi(\theta)=\frac{\beta^\alpha}{\Gamma(\alpha)}\theta^{\alpha-1}e^{-\beta\theta}$,
  mean $\alpha/\beta$.
\item $\ell(\theta)$: log-likelihood; $\hat\theta$: MLE; $w$: weight on the prior mean.
\end{notacao}

\passo{(a) Likelihood}{Independence gives}
\begin{equation}\label{q14:lik}
p(y\mid\theta)=\prod_{i=1}^n\frac{\theta^{y_i}e^{-\theta}}{y_i!}=\frac{\theta^{s}e^{-n\theta}}{\prod_iy_i!}.
\end{equation}

\passo{(a) Posterior}{Bayes' theorem, $\pi(\theta\mid y)\propto p(y\mid\theta)\pi(\theta)$, with
\eqref{q14:lik} and the gamma prior, dropping factors free of $\theta$:}
\begin{equation}\label{q14:post0}
\pi(\theta\mid y)\propto\big(\theta^{s}e^{-n\theta}\big)\big(\theta^{\alpha-1}e^{-\beta\theta}\big).
\end{equation}
Adding exponents in \eqref{q14:post0}:
\begin{equation}\label{q14:post}
\pi(\theta\mid y)\propto\theta^{\alpha+s-1}e^{-(\beta+n)\theta},
\end{equation}
the kernel of $G(\alpha+s,\ \beta+n)$. Prior and posterior are both gamma, so the gamma family is
conjugate for the Poisson likelihood.

\passo{(b) MLE}{Taking logs of \eqref{q14:lik}:}
\begin{equation}\label{q14:ll}
\ell(\theta)=s\log\theta-n\theta-\sum_i\log y_i!.
\end{equation}
Differentiating \eqref{q14:ll} and setting the score to zero:
\begin{equation}\label{q14:score}
\ell'(\theta)=\frac{s}{\theta}-n=0\;\Longrightarrow\;\hat\theta=\frac{s}{n}=\bar y .
\end{equation}
The second derivative $\ell''(\theta)=-s/\theta^2<0$ when $s>0$, so \eqref{q14:score} is a maximum. (If
$s=0$, $\ell(\theta)=-n\theta$ decreases and the supremum is at the boundary $0=\bar y$.)

\passo{(c) Posterior mean as a weighted average}{The mean of the gamma \eqref{q14:post} is}
\begin{equation}\label{q14:pm}
\E(\theta\mid y)=\frac{\alpha+s}{\beta+n}.
\end{equation}
Substituting $\alpha=\beta(\alpha/\beta)$ and $s=n\bar y$ into \eqref{q14:pm}:
\[
\E(\theta\mid y)=\frac{\beta\,(\alpha/\beta)+n\,(\bar y)}{\beta+n}.
\]
Splitting the fraction:
\begin{equation}\label{q14:wavg}
\E(\theta\mid y)=\underbrace{\frac{\beta}{\beta+n}}_{w}\cdot\frac{\alpha}{\beta}
+\underbrace{\frac{n}{\beta+n}}_{1-w}\cdot\bar y .
\end{equation}

\passo{(d) Large $n$}{From \eqref{q14:wavg},}
\begin{equation}\label{q14:lim}
w=\frac{\beta}{\beta+n}\xrightarrow[n\to\infty]{}0,
\end{equation}
so the posterior mean converges to the MLE $\bar y$; the posterior variance $(\alpha+s)/(\beta+n)^2$
also goes to zero at rate $1/n$.

\passo{Numerical illustration}{Coal-mining disasters 1851--1891 ($n=41$, $s=127$,
$\bar y=3.097561$) with the prior of G~7.3, $\alpha=0.5$, $\beta=1$. By \eqref{q14:lim},
$w=1/(1+41)=0.023810$; by \eqref{q14:wavg},}
\[
\E(\theta\mid y)=0.023810\,(0.5)+0.976190\,(3.097561)=0.011905+3.023810=3.035714,
\]
which equals \eqref{q14:pm}: $(0.5+127)/(1+41)=127.5/42=3.035714$.

\codehead
{\small
\begin{verbatim}
grid = np.linspace(1e-6, 10, 200_001)
ker = np.exp((0.5 + 127 - 1) * np.log(grid) - (1 + 41) * grid - 300)  # -300: scale
pm = trapezoid(grid * ker, grid) / trapezoid(ker, grid)                 # 3.035714
w = {n: 1 / (1 + n) for n in (1, 10, 41, 112, 1000)}
# 0.500000  0.090909  0.023810  0.008850  0.000999
\end{verbatim}
}

\intuicao{The prior acts like $\beta$ extra periods of data containing $\alpha$ events in total.
After $n$ real periods it weighs $\beta/(\beta+n)$; with 41 years of data a prior worth one year
moves the estimate only from $3.098$ to $3.036$.}
\refbib{Greenberg (2008), §2.2 (pp.~12--18) and App.~A (gamma distribution); the Class~1 deck
\textit{first Bayesian example} (Poisson with an exponential prior, the case $\alpha=1$).}
\end{sol}

% -----------------------------------------------------------------------------
\questao{G 3.2 --- Poisson--gamma predictive distribution.}
\fonte{Greenberg, E. (2008), \textit{Introduction to Bayesian Econometrics}, 1st ed., Cambridge
University Press, Chapter~3, §3.5 Exercises, Exercise~3.2, p.~39 [PDF p.~54]; the book gives no
answer --- reproduced verbatim.}

3.2 Compute the predictive distribution for $y_{n+1}$ if the $y_i$ have independent Poisson
distributions with parameter $\theta$, given that the first $n$ experiments yielded
$y_1,y_2,\ldots,y_n$ events, respectively. Assume the prior distribution $G(\alpha,\beta)$.
\programacao For the coal data 1851--1891 with prior $G(0.5,1)$, simulate the predictive by
composition (Greenberg §5.2): draw $\theta^{(g)}$ from the posterior and then
$y^{(g)}\sim$ Poisson$(\theta^{(g)})$, $g=1,\dots,10^6$, with \texttt{np.random.default\_rng(15)}.
Compare the simulated frequencies of $y=0,1,2,3$, mean and variance with your closed form, and
check the closed form at $y=0,\dots,3$ against numerical integration of the predictive integral.
\resolespaco
\begin{sol}
\solhead
\begin{notacao}
\item $y_i$, $s=\sum_{i\le n}y_i$, $\theta$, $\alpha$, $\beta$ as in Question~14; $y\in\{0,1,\ldots\}$: value of $y_{n+1}$.
\item $\alpha_1=\alpha+s$, $\beta_1=\beta+n$: posterior parameters.
\end{notacao}
\passo{Posterior}{By \eqref{q14:post} of Question~14, $\theta\mid y_{1:n}\sim G(\alpha_1,\beta_1)$, and
$y_{n+1}$ is independent of $y_{1:n}$ given $\theta$.}
\passo{Predictive integral}{Integrating the Poisson p.m.f.\ against the posterior:}
\begin{align}
p(y_{n+1}=y\mid y_{1:n})&=\int_0^\infty\frac{\theta^{y}e^{-\theta}}{y!}\,\frac{\beta_1^{\alpha_1}}{\Gamma(\alpha_1)}\theta^{\alpha_1-1}e^{-\beta_1\theta}\,d\theta\label{q15:a}\\
&=\frac{\beta_1^{\alpha_1}}{\Gamma(\alpha_1)\,y!}\int_0^\infty\theta^{\alpha_1+y-1}e^{-(\beta_1+1)\theta}d\theta.\label{q15:b}
\end{align}
Substituting the gamma integral \eqref{q10:gi} with $a=\alpha_1+y$, $b=\beta_1+1$ into \eqref{q15:b}:
\begin{align}
p(y_{n+1}=y\mid y_{1:n})&=\frac{\beta_1^{\alpha_1}}{\Gamma(\alpha_1)\,y!}\cdot\frac{\Gamma(\alpha_1+y)}{(\beta_1+1)^{\alpha_1+y}}\label{q15:c}\\
&=\frac{\Gamma(\alpha_1+y)}{\Gamma(\alpha_1)\,y!}\left(\frac{\beta_1}{\beta_1+1}\right)^{\alpha_1}\left(\frac{1}{\beta_1+1}\right)^{y},\qquad y=0,1,\ldots\label{q15:nb}
\end{align}
This is the negative binomial with $r=\alpha_1$ and success probability $p=\beta_1/(\beta_1+1)$.
\passo{Moments}{For that negative binomial, $E=r(1-p)/p$ and $\operatorname{Var}=r(1-p)/p^2$; substituting $r=\alpha_1$, $p=\beta_1/(\beta_1+1)$:}
\begin{align}
E(y_{n+1}\mid y_{1:n})&=\frac{\alpha_1}{\beta_1},\label{q15:m}\\
\operatorname{Var}(y_{n+1}\mid y_{1:n})&=\frac{\alpha_1}{\beta_1}\Big(1+\frac{1}{\beta_1}\Big).\label{q15:v}
\end{align}
\passo{Numbers}{Coal data 1851--1891 with $G(0.5,1)$: $\alpha_1=127.5$, $\beta_1=42$. From
\eqref{q15:nb}, $P(y_{n+1}=0)=(42/43)^{127.5}=0.049780$ and
$P(y_{n+1}=3)=\frac{(129.5)(128.5)(127.5)}{3!}(42/43)^{127.5}(1/43)^3=0.221403$; from
\eqref{q15:m}--\eqref{q15:v}, $E=127.5/42=3.035714$ and $\operatorname{Var}=(3.035714)(1+1/42)=3.107993$.}

\codehead
{\small
\begin{verbatim}
a1, b1 = 127.5, 42.0
pmf = lambda y: nbinom.pmf(y, a1, b1 / (b1 + 1))  # (15.4)
[pmf(k) for k in range(4)]                         # 0.049780 0.147604 0.220548 0.221403
quad(lambda t: poisson.pmf(3, t) * gamma.pdf(t, a1, scale=1/b1), 0, 10)[0]   # 0.221403
rng = np.random.default_rng(15)
yp = rng.poisson(rng.gamma(a1, 1 / b1, 10**6))    # composition
yp.mean(), yp.var(ddof=1)                          # 3.035098  3.101461
[np.mean(yp == k) for k in range(4)]               # 0.049581 0.147294 0.221058 0.221849
\end{verbatim}
}
The simulated frequencies differ from \eqref{q15:nb} by at most $0.0006$, about one Monte Carlo
standard error ($\sqrt{p(1-p)/10^6}\le0.0005$).

\intuicao{The predictive variance exceeds the mean (a Poisson with known $\theta$ would have them
equal): the factor $1+1/\beta_1$ is the extra spread coming from not knowing $\theta$, and it
vanishes as $n$ (hence $\beta_1$) grows. Composition sampling is exactly how the predictive is
defined: draw a parameter, then data given that parameter.}
\refbib{Greenberg (2008), §3.2.3 (prediction), §5.2 (composition), App.~A.1.3 (negative binomial)
and A.1.7.}
\end{sol}

% -----------------------------------------------------------------------------
\questao{G 3.5 --- the posterior median minimizes expected absolute loss.}
\fonte{Greenberg, E. (2008), \textit{Introduction to Bayesian Econometrics}, 1st ed., Cambridge
University Press, Chapter~3, §3.5 Exercises, Exercise~3.5, p.~39 [PDF p.~54]; the book gives no
answer --- reproduced verbatim.}

3.5 Show that the median of the posterior distribution minimizes loss under the absolute value loss
function.
\programacao For the coal posterior $G(127.5,42)$ of Question~15, write the posterior expected
loss $\rho(d)=\int|\theta-d|\,\pi(\theta\mid y)\,d\theta$ with \texttt{scipy.integrate.quad}, minimize it
with \texttt{minimize\_scalar} on $[2.5,3.5]$, and compare the minimizer with the posterior median
(\texttt{gamma.median}). Print $\rho$ at the median, the mean and the mode, and check numerically
that $\rho'(d)=2F(d)-1$ at $d=3$.
\resolespaco
\begin{sol}
\solhead
\begin{notacao}
\item $\pi(\theta\mid y)$: posterior density, c.d.f.\ $F$, assumed continuous with finite mean.
\item $L(\hat\theta,\theta)=|\hat\theta-\theta|$: absolute loss; $d\in\mathbb R$: a candidate estimate.
\item $\rho(d)=\E[|\theta-d|\mid y]$: posterior expected loss; $m$: a posterior median, $F(m)=1/2$.
\end{notacao}

\passo{Split the loss}{On $\theta<d$ the loss is $d-\theta$, on $\theta>d$ it is $\theta-d$:}
\begin{equation}\label{q16:rho}
\rho(d)=\int_{-\infty}^{d}(d-\theta)\,\pi(\theta\mid y)\,d\theta+\int_{d}^{\infty}(\theta-d)\,\pi(\theta\mid y)\,d\theta .
\end{equation}

\passo{Differentiate}{By Leibniz's rule on each integral of \eqref{q16:rho} (the boundary terms are
$(d-d)\pi(d\mid y)=0$):}
\[
\rho'(d)=\int_{-\infty}^{d}\pi(\theta\mid y)\,d\theta-\int_{d}^{\infty}\pi(\theta\mid y)\,d\theta .
\]
Writing both integrals with $F$:
\begin{equation}\label{q16:d1}
\rho'(d)=F(d)-\big(1-F(d)\big)=2F(d)-1 .
\end{equation}

\passo{Minimum}{By \eqref{q16:d1}, $\rho'(d)<0$ for $F(d)<1/2$ and $\rho'(d)>0$ for $F(d)>1/2$:
$\rho$ decreases up to the median and increases after it. Hence $\rho$ is minimized at $d=m$ with
$F(m)=1/2$. Equivalently, $\rho''(d)=2\pi(d\mid y)\ge0$: $\rho$ is convex and \eqref{q16:d1} is
the first-order condition of a global minimum. $\blacksquare$ (If $F$ is flat at $1/2$ on an interval,
every point of it is a median and minimizes $\rho$.)}

\passo{Coal posterior}{$G(127.5,42)$ has median $3.027781$, mean $127.5/42=3.035714$ and mode
$126.5/42=3.011905$: the right skew of the gamma orders them mode $<$ median $<$ mean.}

\codehead
{\small
\begin{verbatim}
post = gamma(127.5, scale=1/42)
rho = lambda d: quad(lambda t: abs(t - d) * post.pdf(t), 0, 10, points=[d])[0]
res = minimize_scalar(rho, bounds=(2.5, 3.5), method="bounded")
res.x, post.median()                    # 3.027781  3.027781
rho(post.median()), rho(post.mean()), rho(126.5 / 42)   # 0.214276 0.214369 0.214650
(rho(3 + 1e-5) - rho(3 - 1e-5)) / 2e-5, 2 * post.cdf(3) - 1   # equal, as in (16.2)
\end{verbatim}
}
The numerical minimizer is the median; using the mean or the mode instead costs little here
because the posterior is nearly symmetric.

\intuicao{Moving the estimate up by a small step lowers the loss for every draw above it and raises
it for every draw below, each by the same amount. The estimate stops moving when half the posterior
mass lies on each side. Squared loss weights far draws more and pulls towards the mean.}
\refbib{Greenberg (2008), §3.2.1 (point estimates and loss functions, pp.~29--31).}
\end{sol}

% -----------------------------------------------------------------------------
\questao{B 5.12 (with B 4.9) --- Bayes estimator of $P(X=0)$ and a Bayesian interval.}
\fonte{Bolfarine, H. and Sandoval, M.~C. (2010), \textit{Introdução à Inferência Estatística},
Chapter~5, §5.6 Exercícios, Exercício~5.12, p.~89 [PDF p.~95]; it refers to Exercício~4.9,
Chapter~4, §4.5 Exercícios, p.~72 [PDF p.~78], reproduced first. The book gives no answer to
either --- both reproduced verbatim.}

\begin{otherlanguage}{brazil}
\emph{Exercício referido:} \textbf{4.9.} Em uma área de reflorestamento, o número de árvores de
determinada espécie, por hectare, com certa doença tem uma distribuição $\mathit{Poisson}(\theta)$. A
distribuição a priori de $\theta$ é exponencial com média igual a 1. Encontre o estimador de Bayes
de $P_\theta(X=0)$ com relação à perda quadrática..

\textbf{5.12.} Considere o Exercício 4.9. Obtenha um intervalo de confiança Bayesiano para $\theta$
com coeficiente de confiança $\gamma=0,95$. Como fica seu intervalo se $x=4$?
\end{otherlanguage}
\programacao For $x=4$: (1) compute the equal-tailed 95\% interval with \texttt{scipy.stats.gamma}
and check it against the $\chi^2_{10}$ quantiles; (2) compute the HPD interval numerically (minimize
the length of $[F^{-1}(p),F^{-1}(p+0.95)]$ over $p$) and verify that the density is equal at both
ends; (3) compute the exact frequentist coverage of the equal-tailed rule $x\mapsto I(x)$ at
$\theta\in\{0.5,1,4,8\}$ by summing Poisson probabilities.
\resolespaco
\begin{sol}
\solhead
\begin{notacao}
\item $X\mid\theta\sim$ Poisson$(\theta)$: one observation $x$; prior $\pi(\theta)=e^{-\theta}$, $\theta>0$ (exponential, mean 1 $=G(1,1)$).
\item $G(a,b)$: shape $a$, rate $b$; $q_p$: posterior $p$-quantile; $\chi^2_{\nu,p}$: chi-square quantile.
\end{notacao}

\passo{Posterior (4.9)}{By Bayes' theorem,}
\begin{equation}\label{q17:post0}
\pi(\theta\mid x)\propto\frac{\theta^xe^{-\theta}}{x!}\cdot e^{-\theta}\propto\theta^{x}e^{-2\theta},
\end{equation}
the kernel of $G(x+1,2)$ (the case $\alpha=\beta=1$, $n=1$ of \eqref{q14:post}):
\begin{equation}\label{q17:post}
\pi(\theta\mid x)=\frac{2^{x+1}}{x!}\theta^xe^{-2\theta}.
\end{equation}

\passo{Bayes estimator (4.9)}{Under quadratic loss the Bayes estimator of $g(\theta)=P_\theta(X=0)=e^{-\theta}$
is its posterior mean. Using \eqref{q17:post}:}
\[
\E[e^{-\theta}\mid x]=\int_0^\infty e^{-\theta}\,\frac{2^{x+1}}{x!}\theta^xe^{-2\theta}d\theta=\frac{2^{x+1}}{x!}\int_0^\infty\theta^xe^{-3\theta}d\theta .
\]
Substituting \eqref{q10:gi} with $a=x+1$, $b=3$:
\begin{equation}\label{q17:bayes}
\E[e^{-\theta}\mid x]=\frac{2^{x+1}}{x!}\cdot\frac{x!}{3^{x+1}}=\Big(\frac23\Big)^{x+1}.
\end{equation}
At $x=4$: $(2/3)^5=0.131687$. (It is not $e^{-\E(\theta|x)}=e^{-5/2}$: the Bayes estimator of a
function is the posterior mean of the function.)

\passo{Pivot for the interval (5.12)}{If $\theta\mid x\sim G(x+1,2)$ then $4\theta\mid x\sim G(x+1,1/2)$,
which is $\chi^2_{2(x+1)}$. The symmetric (equal-tailed) interval with $\gamma=0.95$ is}
\begin{equation}\label{q17:int}
\Big[\frac{\chi^2_{2x+2,\,0.025}}{4},\ \frac{\chi^2_{2x+2,\,0.975}}{4}\Big].
\end{equation}
For $x=4$, $\chi^2_{10,0.025}=3.246973$ and $\chi^2_{10,0.975}=20.483177$; substituting into
\eqref{q17:int}:
\[
\Big[\frac{3.246973}{4},\ \frac{20.483177}{4}\Big]=[0.811743,\ 5.120794].
\]

\passo{HPD (programming)}{The posterior is right-skewed, so the shortest 95\% interval is shifted
left: $[0.603480,\ 4.715108]$, length $4.111628$ against $4.309051$ for \eqref{q17:int}, with equal
density at both ends.}

\codehead
{\small
\begin{verbatim}
pg = gamma(5, scale=0.5)
pg.ppf([.025, .975])                      # 0.811743 5.120794 = chi2(10) quantiles / 4
L = lambda p: pg.ppf(p + .95) - pg.ppf(p)
p = minimize_scalar(L, bounds=(1e-9, .05 - 1e-9), method="bounded").x
pg.ppf([p, p + .95])                      # HPD 0.603480 4.715108
xs = np.arange(200); lo, hi = gamma.ppf([[.025], [.975]], xs + 1, scale=.5)
cov = lambda th: np.sum(poisson.pmf(xs, th) * ((lo <= th) & (th <= hi)))
# coverage at theta = 0.5, 1, 4, 8:  0.985612  0.996340  0.761820  0.407449
\end{verbatim}
}
A 95\% \emph{Bayesian} interval need not have 95\% \emph{frequentist} coverage: with a single
observation the prior (mean 1) dominates, so the rule over-covers near $\theta=1$ and badly
under-covers large $\theta$ (40.7\% at $\theta=8$).

\intuicao{With one count the posterior is a compromise between the observation and a prior that
expects about one diseased tree per hectare; the interval is a statement about $\theta$ given this
$x$, not a promise about repeated sampling.}
\refbib{Bolfarine \& Sandoval (2010), §4.3 (estimador de Bayes, perda quadrática), §5.5 (intervalos
Bayesianos, Exemplo 5.5.1).}
\end{sol}

% -----------------------------------------------------------------------------
\questao{Original math question --- Bayes estimators and their frequentist risk.}
\fonte{Original question (not from a book).}

An insurer prices a policy on the number of claims per client and year, $X_1,\dots,X_n$ i.i.d.\
Poisson$(\theta)$, with a prior $\theta\sim G(\alpha,\beta)$ (shape, rate) from past portfolios;
$S=\sum_iX_i$.
\begin{enumerate}[label=\textbf{(\alph*)},leftmargin=2.4em]
\item Derive the Bayes estimator under squared loss, $d_1$, and under the weighted loss
$L(\theta,d)=(\theta-d)^2/\theta$, $d_2$, stating when $d_2$ exists. Relate $d_2$ to a feature of the
posterior.
\item Derive the frequentist risk $R_1(\theta)=E_\theta[(d_1-\theta)^2]$ and prove that
$R_1(\theta)<\theta/n$ (the risk of the MLE $\overline X$) exactly on an interval of $\theta$ that
contains the prior mean $\alpha/\beta$; give the interval.
\item Compute the Bayes risks $r=\int R(\theta)\pi(\theta)d\theta$ of $d_1$ and of $\overline X$ and prove
$r(d_1)<r(\overline X)$.
\item Parameter analysis: sign of the bias of $d_1$ as a function of $\theta$; the variance ratio
$\Var(d_1)/\Var(\overline X)$; the limits of the interval of (b) as $n\to\infty$ and as $\beta\to0$
with $\alpha$ fixed. Interpret for the insurer.
\end{enumerate}
\programacao With $\alpha=2$, $\beta=1$, $n=10$: print the interval of (b), its $n\to\infty$ limit and
the two Bayes risks; with \texttt{np.random.default\_rng(18)} simulate $400{,}000$ values of $S$ at
$\theta\in\{1,2,5,6\}$ and compare the simulated mean squared errors of $d_1$ and $\overline X$ with
the formulas.
\resolespaco
\begin{sol}
\solhead
\begin{notacao}
\item $S\sim$ Poisson$(n\theta)$; posterior $\theta\mid\mathbf x\sim G(a,b)$, $a=\alpha+S$, $b=\beta+n$ (by \eqref{q14:post}).
\item $\rho_j(d)$: posterior expected loss; $R(\theta)$: frequentist risk (MSE); $r$: Bayes risk.
\item $\alpha,\beta>0$, $n\ge1$.
\end{notacao}

\passo{(a) Squared loss}{$\rho_1(d)=\E[(\theta-d)^2\mid\mathbf x]=\Var(\theta\mid\mathbf x)+(\E(\theta\mid\mathbf x)-d)^2$
is minimized at the posterior mean, so by \eqref{q14:pm}}
\begin{equation}\label{q18:d1}
d_1=\frac{\alpha+S}{\beta+n}.
\end{equation}

\passo{(a) Weighted loss}{Expanding the square,}
\begin{equation}\label{q18:rho2}
\rho_2(d)=\E\Big[\frac{(\theta-d)^2}{\theta}\Big|\,\mathbf x\Big]=\E(\theta\mid\mathbf x)-2d+d^2\,\E(\theta^{-1}\mid\mathbf x),
\end{equation}
a convex quadratic in $d$ with minimum where $-2+2d\,\E(\theta^{-1}\mid\mathbf x)=0$:
\begin{equation}\label{q18:d2a}
d_2=\frac{1}{\E(\theta^{-1}\mid\mathbf x)} .
\end{equation}
For $G(a,b)$, using \eqref{q10:gi} with exponent $a-1$:
$\E(\theta^{-1})=\frac{b^a}{\Gamma(a)}\int_0^\infty\theta^{a-2}e^{-b\theta}d\theta=\frac{b^a}{\Gamma(a)}\cdot\frac{\Gamma(a-1)}{b^{a-1}}=\frac{b}{a-1}$,
finite iff $a>1$. Substituting into \eqref{q18:d2a}:
\begin{equation}\label{q18:d2}
d_2=\frac{a-1}{b}=\frac{\alpha+S-1}{\beta+n},\qquad\text{if }\alpha+S>1,
\end{equation}
which is the posterior \emph{mode}: dividing the loss by $\theta$ makes errors at small $\theta$
costly and pulls the estimate down by $1/(\beta+n)$.

\passo{(b) Bias and variance of $d_1$}{From \eqref{q18:d1} with $E_\theta S=\Var_\theta S=n\theta$:}
\begin{equation}\label{q18:bv}
E_\theta d_1-\theta=\frac{\alpha+n\theta-(\beta+n)\theta}{\beta+n}=\frac{\alpha-\beta\theta}{\beta+n},\qquad
\Var_\theta(d_1)=\frac{n\theta}{(\beta+n)^2}.
\end{equation}
Adding variance and squared bias:
\begin{equation}\label{q18:R1}
R_1(\theta)=\frac{n\theta+(\alpha-\beta\theta)^2}{(\beta+n)^2}.
\end{equation}

\passo{(b) Where $d_1$ beats the MLE}{$R_1(\theta)<\theta/n$; multiplying by $n(\beta+n)^2>0$ and
substituting \eqref{q18:R1}:}
\[
n\big[n\theta+(\alpha-\beta\theta)^2\big]<\theta(\beta+n)^2 .
\]
Expanding both sides and collecting powers of $\theta$:
\begin{equation}\label{q18:quad}
Q(\theta)\equiv n\beta^2\theta^2-\big(\beta^2+2n\beta+2n\alpha\beta\big)\theta+n\alpha^2<0 .
\end{equation}
$Q$ is an upward parabola, so \eqref{q18:quad} holds exactly between its roots $\theta_-<\theta_+$.
At the prior mean, substituting $\theta=\alpha/\beta$ into \eqref{q18:quad}:
\[
Q(\alpha/\beta)=n\alpha^2-(\beta^2+2n\beta+2n\alpha\beta)\frac\alpha\beta+n\alpha^2=-\alpha\beta-2n\alpha<0,
\]
so the roots are real and distinct and $\alpha/\beta\in(\theta_-,\theta_+)$. $\blacksquare$

\passo{(c) Bayes risks}{Under the prior, $E\theta=\alpha/\beta$ and $E(\alpha-\beta\theta)^2=\beta^2\Var(\theta)=\beta^2(\alpha/\beta^2)=\alpha$.
Averaging \eqref{q18:R1}:}
\begin{equation}\label{q18:r1}
r(d_1)=\frac{n\alpha/\beta+\alpha}{(\beta+n)^2}=\frac{\alpha(n+\beta)}{\beta(\beta+n)^2}=\frac{\alpha}{\beta(\beta+n)},
\qquad r(\overline X)=E\Big[\frac\theta n\Big]=\frac{\alpha}{\beta n}.
\end{equation}
Since $\beta+n>n$, $r(d_1)<r(\overline X)$ (as it must: $d_1$ minimizes the Bayes risk).

\passo{(d) Signs and limits}{By \eqref{q18:bv} the bias is positive for $\theta<\alpha/\beta$, zero at
the prior mean and negative above it: $d_1$ shrinks towards $\alpha/\beta$. The variance ratio is
$\Var(d_1)/\Var(\overline X)=(n/(\beta+n))^2<1$. Dividing \eqref{q18:quad} by $n$:
$\beta^2\theta^2-(\beta^2/n+2\beta+2\alpha\beta)\theta+\alpha^2<0$; as $n\to\infty$ this tends to
$\beta^2\theta^2-2\beta(1+\alpha)\theta+\alpha^2<0$, with roots}
\begin{equation}\label{q18:lim}
\theta_\pm^\infty=\frac{(1+\alpha)\pm\sqrt{(1+\alpha)^2-\alpha^2}}{\beta}=\frac{(1+\alpha)\pm\sqrt{1+2\alpha}}{\beta}:
\end{equation}
the region does not vanish with $n$ (both risks shrink like $1/n$ and their ratio stays informative).
As $\beta\to0$ with $\alpha$ fixed, $\theta_-\approx n\alpha^2/(\beta^2+2n\beta(1+\alpha))\to\infty$
and the region disappears: \eqref{q18:R1} tends to $\theta/n+\alpha^2/n^2>\theta/n$, a pure bias with
no shrinkage.

\passo{(d) Reading}{The prior helps the insurer when the true claim rate is near the portfolio
history and hurts when the client is very different; the band \eqref{q18:lim} measures how far is
``very different''. A flat-but-shifted prior ($\beta\to0$) is the worst of both worlds.}

\passo{Numbers}{$\alpha=2$, $\beta=1$, $n=10$: \eqref{q18:quad} is
$10\theta^2-(1+20+40)\theta+40=10\theta^2-61\theta+40<0$, discriminant $61^2-4(10)(40)=2121$, roots
$(61\mp\sqrt{2121})/20=(0.747284,\ 5.352716)$; \eqref{q18:lim}: $3\mp\sqrt5=(0.763932,\ 5.236068)$;
\eqref{q18:r1}: $r(d_1)=2/11=0.181818<r(\overline X)=2/10=0.2$.}

\codehead
{\small
\begin{verbatim}
al, be, n = 2., 1., 10
rB = lambda t: (n * t + (al - be * t)**2) / (be + n)**2
rng = np.random.default_rng(18)
for th in (1, 2, 5, 6):
    S = rng.poisson(n * th, 400_000)
    print(th, rB(th), np.mean(((al + S)/(be + n) - th)**2), th/n, np.mean((S/n - th)**2))
# theta  R1 formula  R1 sim    MLE formula  MLE sim
# 1      0.090909    0.091203  0.1          0.100222
# 2      0.165289    0.164913  0.2          0.199545
# 5      0.487603    0.487363  0.5          0.499468
# 6      0.628099    0.629657  0.6          0.602826
\end{verbatim}
}
Inside $(0.747,5.353)$ (at $\theta=1,2,5$) the Bayes estimator wins; at $\theta=6$ it loses, as (b)
predicts.

\intuicao{Shrinkage trades a little bias for less variance. The trade pays off near the prior
mean and on average over the prior, never uniformly: by (b), $d_1$ loses to the MLE outside a
bounded interval of $\theta$, whatever $n$.}
\refbib{Bolfarine \& Sandoval (2010), §4.2--4.3 (risco, risco de Bayes, estimador de Bayes);
Greenberg (2008), §3.2.1 (loss functions).}
\rubrica{
\item (a) $d_1$: 2 pts; \eqref{q18:rho2}--\eqref{q18:d2} with the existence condition and the mode: 3 pts.
\item (b) \eqref{q18:bv}--\eqref{q18:R1}: 3 pts; \eqref{q18:quad} and $Q(\alpha/\beta)<0$: 3 pts.
\item (c) \eqref{q18:r1} and the inequality: 2 pts.
\item (d) bias sign, variance ratio, \eqref{q18:lim}, $\beta\to0$, reading: 4 pts (limits without derivation: max.\ half).
\item Programming: interval, Bayes risks and simulated risks: 3 pts.
}
\end{sol}

% -----------------------------------------------------------------------------
\questao{Original conceptual question --- gamma--Poisson Bayes (10 points).}
\fonte{Original question (not from a book).}

An analyst at ANS (the health-insurance regulator) models claims per beneficiary as Poisson$(\theta)$
with prior $G(\alpha,\beta)$ and observes $n$ beneficiaries with total $s$.

\textbf{a)} (3 points) In the posterior mean written as a weighted average of the prior mean
$\alpha/\beta$ and the MLE $\bar y$, the weight on $\bar y$ is:
\begin{alternativas}
\item $n/(\beta+n)$, the share of real observations in the total ``sample size''.
\item $\beta/(\beta+n)$, the share of the prior's pseudo-observations.
\item $n/(\alpha+n)$, with the prior's events counted as observations.
\item $\alpha/(\alpha+\beta)$, which does not depend on the sample.
\end{alternativas}
\gabarito{A}
\begin{sol}
\passo{Why}{A: \eqref{q14:wavg}. B is the weight on the \emph{prior} mean; C mixes up the prior's
events $\alpha$ with its periods $\beta$; D ignores the data entirely.}
\end{sol}

\textbf{b)} (3 points) The predictive distribution of a new beneficiary's count $y_{n+1}$ is:
\begin{alternativas}
\item Poisson with mean $(\alpha+s)/(\beta+n)$, plugging in the posterior mean.
\item gamma $G(\alpha+s,\beta+n)$, the posterior of $\theta$.
\item negative binomial, with variance larger than its mean.
\item negative binomial, with variance smaller than its mean.
\end{alternativas}
\gabarito{C}
\begin{sol}
\passo{Why}{C: \eqref{q15:nb}, with \eqref{q15:v} $=$ mean $\times(1+1/\beta_1)$. A ignores the
uncertainty about $\theta$ (plug-in); B is a distribution for the parameter, not for a count;
D has the inequality reversed.}
\end{sol}

\textbf{c)} (2 points, T or F) Under absolute-error loss the Bayes estimator of $\theta$ is the
posterior median.
\gabarito{T}
\begin{sol}
\passo{Why}{Question~16, \eqref{q16:d1}.}
\end{sol}

\textbf{d)} (2 points, T or F) Since the Bayes estimator under squared loss is the posterior mean, it
is unbiased for $\theta$ in repeated sampling.
\gabarito{F}
\begin{sol}
\passo{Why}{Its frequentist bias is $(\alpha-\beta\theta)/(\beta+n)$ by \eqref{q18:bv}, zero only at
$\theta=\alpha/\beta$. ``Mean'' refers to the posterior, not to the sampling distribution.}
\end{sol}

\programacao Check (d) with $\alpha=2$, $\beta=1$, $n=10$, $\theta=1$: with
\texttt{np.random.default\_rng(19)} simulate $400{,}000$ totals $S\sim$ Poisson$(10)$ and report the
average of $(\alpha+S)/(\beta+n)-\theta$.
\begin{sol}
\codehead
{\small
\begin{verbatim}
rng = np.random.default_rng(19); S = rng.poisson(10, 400_000)
np.mean((2 + S) / 11) - 1                       # 0.090347 ; formula 1/11 = 0.090909
\end{verbatim}
}
\end{sol}

% =============================================================================
\bloco{5}{Monte Carlo for inference: integration, importance sampling, accept--reject}{G 5.6, 5.7; MS B.15, B.18; original}

\questao{G 5.6 --- the GHK importance-sampling estimator (5.3).}
\fonte{Greenberg, E. (2008), \textit{Introduction to Bayesian Econometrics}, 1st ed., Cambridge
University Press, Chapter~5, §5.8 Exercises, Exercise~5.6, p.~75 [PDF p.~90]; the book gives no
answer --- reproduced verbatim. The text it refers to (§5.4, pp.~71--72 [PDF pp.~86--87], including
Equation (5.3)) and Equation (5.1) (§5.1, p.~65 [PDF p.~80]) are reproduced first.}

\emph{Referred to (Greenberg 2008, p.~65):} ``we generate $U\sim U(0,1)$ and set
$U=\frac{F(X)-F(c_1)}{F(c_2)-F(c_1)}$, which implies that
\[
X=F^{-1}\big(F(c_1)+U[F(c_2)-F(c_1)]\big)\tag{5.1}
\]
is a drawing from the truncated distribution.''

\emph{Referred to (Greenberg 2008, pp.~71--72):} ``As a second example, we approximate
$P(a_1<X_1\le b_1,a_2<X_2\le b_2)$, where $(X_1,X_2)\sim N_2(\mu,\Sigma)$, $\mu=(\mu_1,\mu_2)'$, and
$\Sigma=[\sigma_{ij}]$. The desired integral is
\[
P(a_1<X_1\le b_1,a_2<X_2\le b_2)=\int_{a_1}^{b_1}\int_{a_2}^{b_2}f(x_1,x_2)\,dx_1\,dx_2,
\]
where $f(x_1,x_2)$ is the density function of the bivariate normal distribution. The first step is
to rewrite the joint density in marginal-conditional form, $f(x_1)f(x_2|x_1)$ using the results (see
Section A.9) that $x_1\sim N(\mu_1,\sigma_{11})$ and that $x_2|x_1\sim N(\mu_2',\sigma_{22}')$, where
\[
\mu_2'=\mu_2+(\sigma_{12}/\sigma_{11})(x_1-\mu_1),\qquad \sigma_{22}'=\sigma_{22}-\sigma_{12}^2/\sigma_{11}.
\]
Now let
\begin{align*}
y_1&=(x_1-\mu_1)/\sqrt{\sigma_{11}}, & y_2&=(x_2-\mu_2')/\sqrt{\sigma_{22}'},\\
a_1'&=(a_1-\mu_1)/\sqrt{\sigma_{11}}, & a_2'&=(a_2-\mu_2')/\sqrt{\sigma_{22}'},\\
b_1'&=(b_1-\mu_1)/\sqrt{\sigma_{11}}, & b_2'&=(b_2-\mu_2')/\sqrt{\sigma_{22}'}.
\end{align*}
With this transformation, $y_1$ and $y_2|y_1$ have standard normal distributions, enabling us to write
\begin{align*}
P(a_1<x_1\le b_1,a_2<x_2\le b_2)&=\int_{a_1'}^{b_1'}\int_{a_2'}^{b_2'}\phi(y_1)\phi(y_2|y_1)\,dy_1\,dy_2\\
&=\int_{a_1'}^{b_1'}\phi(y_1)\int_{a_2'}^{b_2'}\phi(y_2|y_1)\,dy_2\,dy_1\\
&=\int_{a_1'}^{b_1'}\phi(y_1)[\Phi(b_2'|y_1)-\Phi(a_2'|y_1)]\,dy_1.
\end{align*}
The algorithm is now implemented by choosing the standard normal distribution truncated to
$(a_1',b_1')$ as the importance distribution; that is, draw $Y_1^{(g)}$ from the truncated standard
normal using (5.1). With this choice for the importance function, we have
\[
P(a_1<X_1\le b_1,a_2<X_2\le b_2)\approx\frac1G\sum\big[\Phi\big(b_2'|Y_1^{(g)}\big)-\Phi\big(a_2'|Y_1^{(g)}\big)\big]\big[\Phi(b_1')-\Phi(a_1')\big].\text{''}\tag{5.3}
\]

5.6 Verify Equation (5.3).
\programacao Use the numbers of G~5.7 (next question). (1) Implement the truncated-normal draw
(5.1) and compare the mean of $10{,}000$ draws (\texttt{np.random.default\_rng(21)}) with
\texttt{scipy.stats.truncnorm}. (2) Compute by quadrature the exact variance of one term of (5.3)
and of one crude indicator $\1\{a_1<X_1\le b_1,a_2<X_2\le b_2\}$, their ratio, and check the bound
you prove in the solution.
\resolespaco
\begin{sol}
\solhead
\begin{notacao}
\item $p=P(a_1<X_1\le b_1,a_2<X_2\le b_2)$: target; $\phi,\Phi$: standard normal density and c.d.f.
\item $h(y_1)=\Phi(b_2'|y_1)-\Phi(a_2'|y_1)\in[0,1]$, where $\Phi(c'|y_1)=\Phi\big((c-\mu_2'(y_1))/\sqrt{\sigma_{22}'}\big)$
  and $\mu_2'(y_1)=\mu_2+(\sigma_{12}/\sigma_{11})\sqrt{\sigma_{11}}\,y_1$.
\item $m_1=\Phi(b_1')-\Phi(a_1')$; $q(y_1)=\phi(y_1)\1\{a_1'<y_1\le b_1'\}/m_1$: truncated normal (importance) density.
\end{notacao}

\passo{Inner integral}{Given $x_1$, $x_2\sim N(\mu_2',\sigma_{22}')$, so standardizing,}
\begin{equation}\label{q20:inner}
\int_{a_2}^{b_2}f(x_2|x_1)\,dx_2=\Phi\Big(\frac{b_2-\mu_2'}{\sqrt{\sigma_{22}'}}\Big)-\Phi\Big(\frac{a_2-\mu_2'}{\sqrt{\sigma_{22}'}}\Big)=h(y_1).
\end{equation}

\passo{Outer integral}{With $x_1=\mu_1+\sqrt{\sigma_{11}}y_1$, $dx_1=\sqrt{\sigma_{11}}\,dy_1$ and
$f(x_1)dx_1=\phi(y_1)dy_1$; substituting \eqref{q20:inner} into the marginal--conditional form:}
\begin{equation}\label{q20:p}
p=\int_{a_1}^{b_1}f(x_1)\Big(\int_{a_2}^{b_2}f(x_2|x_1)dx_2\Big)dx_1=\int_{a_1'}^{b_1'}\phi(y_1)\,h(y_1)\,dy_1 .
\end{equation}

\passo{Importance-sampling form}{Multiplying and dividing the integrand of \eqref{q20:p} by $m_1$:}
\begin{equation}\label{q20:is}
p=\int\Big(\frac{\phi(y_1)\1\{a_1'<y_1\le b_1'\}}{m_1}\Big)\big(m_1\,h(y_1)\big)\,dy_1=E_q\big[m_1\,h(Y_1)\big].
\end{equation}
The importance weight $\phi/q=m_1$ is constant on the support. If $Y_1^{(g)}$ are i.i.d.\ from $q$ ---
which is what (5.1) delivers with $F=\Phi$, $c_1=a_1'$, $c_2=b_1'$ --- then by \eqref{q20:is} each term
$m_1h(Y_1^{(g)})$ has mean $p$, and by the strong law of large numbers
\[
\frac1G\sum_{g=1}^G\big[\Phi(b_2'|Y_1^{(g)})-\Phi(a_2'|Y_1^{(g)})\big]\big[\Phi(b_1')-\Phi(a_1')\big]\xrightarrow{a.s.}p,
\]
which is (5.3), unbiased for every $G$. $\blacksquare$

\passo{It beats crude Monte Carlo}{Since $0\le h\le1$, $h^2\le h$, so}
\begin{equation}\label{q20:var}
\Var_q(m_1h)=m_1^2E_q[h^2]-p^2\le m_1^2E_q[h]-p^2=m_1p-p^2=p(m_1-p)\le p(1-p),
\end{equation}
the last term being the variance of one crude indicator. The gain is large when $m_1$ is small
(the first coordinate's constraint is handled exactly) and when $h$ varies little.

\passo{Numbers (G~5.7)}{$\sigma_{22}'=3-1^2/2=2.5$; $a_1'=(0.2-1)/\sqrt2=-0.565685$,
$b_1'=(2-1)/\sqrt2=0.707107$, $m_1=0.474446$; $p=0.300105$ (quadrature of \eqref{q20:p}). One GHK
term has variance $0.000797$, below the bound \eqref{q20:var} $p(m_1-p)=0.052321$; a crude
indicator has $p(1-p)=0.210042$: a ratio of $263.6$, so one GHK draw is worth about 264 crude draws.}

\codehead
{\small
\begin{verbatim}
a1p, b1p = (0.2 - 1) / np.sqrt(2), (2 - 1) / np.sqrt(2)
rng = np.random.default_rng(21); u = rng.uniform(size=10_000)
y1 = norm.ppf(norm.cdf(a1p) + u * (norm.cdf(b1p) - norm.cdf(a1p)))   # (5.1)
y1.mean(), truncnorm(a1p, b1p).mean()       # 0.0655  0.0617 (se ~0.004)
h = lambda y: (norm.cdf((5 - mu2c(y)) / np.sqrt(2.5))
               - norm.cdf((-1 - mu2c(y)) / np.sqrt(2.5)))   # mu2c(y) = -0.5 + y/sqrt(2)
m1 = norm.cdf(b1p) - norm.cdf(a1p)                            # 0.474446
p = quad(lambda y: norm.pdf(y) * h(y), a1p, b1p)[0]           # 0.300105
v = quad(lambda y: norm.pdf(y) / m1 * (m1 * h(y))**2, a1p, b1p)[0] - p**2   # 0.000797
p * (1 - p) / v                                               # 263.6
\end{verbatim}
}

\intuicao{GHK integrates the second coordinate exactly, given the first, and samples the first only
where it matters. Every simulation draw lands in the region and contributes a smooth probability,
instead of a 0/1 indicator that wastes most draws.}
\refbib{Greenberg (2008), §5.1 (eq.~5.1), §5.4 (importance sampling, pp.~70--72) and §5.7 (GHK).}
\end{sol}

% -----------------------------------------------------------------------------
\questao{G 5.7 --- importance sampling for a bivariate normal rectangle.}
\fonte{Greenberg, E. (2008), \textit{Introduction to Bayesian Econometrics}, 1st ed., Cambridge
University Press, Chapter~5, §5.8 Exercises, Exercise~5.7, p.~75 [PDF p.~90]; the book gives no
answer --- reproduced verbatim. The method is the one of (5.3), reproduced in Question~20.}

5.7 Use importance sampling to compute $P(0.2<x_1\le2,\ -1<x_2\le5)$, where
$(x_1,x_2)\sim N_2(\mu,\Sigma)$, $\mu=(1,-0.5)'$, and
\[
\Sigma=\begin{pmatrix}2&1\\1&3\end{pmatrix}.
\]
\programacao With $G=10{,}000$ and \texttt{np.random.default\_rng(21)} (the same truncated-normal
draws as Question~20), report the estimate (5.3) and its numerical standard error; compare with
the exact value (one-dimensional quadrature of the last integral of the text quoted in
Question~20, cross-checked with
\texttt{multivariate\_normal.cdf} and inclusion--exclusion) and with crude Monte Carlo using
$10{,}000$ bivariate normal draws from the same generator.
\resolespaco
\begin{sol}
\solhead
\begin{notacao}
\item As in Question~20, with $\mu_1=1$, $\mu_2=-0.5$, $\sigma_{11}=2$, $\sigma_{12}=1$, $\sigma_{22}=3$; $a_1=0.2$, $b_1=2$, $a_2=-1$, $b_2=5$.
\item $\hat p_G$: estimate (5.3); n.s.e.\ $=s/\sqrt G$, $s$ the standard deviation of the $G$ terms.
\end{notacao}

\passo{Ingredients}{From Question~20: $\sigma_{22}'=2.5$, $a_1'=-0.565685$, $b_1'=0.707107$,
$m_1=0.474446$, and}
\[
\mu_2'(y_1)=-0.5+\tfrac12\big((1+\sqrt2\,y_1)-1\big)=-0.5+\frac{y_1}{\sqrt2}.
\]

\passo{Algorithm}{For $g=1,\dots,G$: draw $U^{(g)}\sim U(0,1)$, set
$Y_1^{(g)}=\Phi^{-1}\big(\Phi(a_1')+U^{(g)}m_1\big)$ by (5.1), compute
$t_g=m_1\big[\Phi\big((5-\mu_2'(Y_1^{(g)}))/\sqrt{2.5}\big)-\Phi\big((-1-\mu_2'(Y_1^{(g)}))/\sqrt{2.5}\big)\big]$; return
$\hat p_G=\bar t$ and n.s.e.\ $s_t/\sqrt G$.}

\passo{Result}{$\hat p_G=0.300411$ with n.s.e.\ $0.000282$; the exact value is $0.300105$, about
$1.1$ n.s.e.\ away. Crude Monte Carlo with the same $G$ gives $0.3019$ with standard error $0.0046$,
sixteen times larger, consistent with the variance ratio $263.6$ of Question~20
($\sqrt{263.6}\approx16.2$).}

\codehead
{\small
\begin{verbatim}
t = m1 * h(y1)                                   # y1 from (5.1), Question 20
t.mean(), t.std(ddof=1) / np.sqrt(10_000)        # 0.300411  0.000282
mvn = multivariate_normal([1, -.5], [[2, 1], [1, 3]])
mvn.cdf([2, 5]) - mvn.cdf([.2, 5]) - mvn.cdf([2, -1]) + mvn.cdf([.2, -1])  # 0.300105
x = mvn.rvs(size=10_000, random_state=rng)
ind = (.2 < x[:, 0]) & (x[:, 0] <= 2) & (-1 < x[:, 1]) & (x[:, 1] <= 5)
ind.mean(), ind.std(ddof=1) / 100                # 0.3019  0.0046
\end{verbatim}
}

\intuicao{The rectangle is wide in $x_2$ (from $-1$ to $5$) and narrow in $x_1$, so once $x_1$ is
drawn inside its interval the conditional probability for $x_2$ is high and stable; nearly all the
uncertainty is removed analytically.}
\refbib{Greenberg (2008), §5.4 (pp.~70--72).}
\end{sol}

% -----------------------------------------------------------------------------
\questao{MS B.15 --- generating Poisson values by the inverse transform.}
\fonte{Morettin, P.~A. and Singer, J.~M. (2021), \textit{Estatística e Ciência de Dados} (versão
preliminar, abril 2021), Apêndice~B (Noções de simulação), §B.8 Exercícios, Exercício~15, p.~478
[PDF p.~488]; the book gives no answer --- reproduced verbatim.}

\begin{otherlanguage}{brazil}
15) Gere 5 valores de uma variável aleatória com distribuição de Poisson de parâmetro $\lambda=2$.
\end{otherlanguage}
\programacao Write \texttt{rpois\_inv(u, lam)} returning the smallest $j$ with $F(j)>u$, built with
the recursion $p_{j+1}=p_j\lambda/(j+1)$, and the number of comparisons it made. With
\texttt{np.random.default\_rng(22)}: generate the 5 values; then $100{,}000$ more, and report their
mean, variance, the average number of comparisons, and a chi-square goodness-of-fit test on the
classes $0,1,\dots,6,\ge7$.
\resolespaco
\begin{sol}
\solhead
\begin{notacao}
\item $p_j=P(X=j)=e^{-\lambda}\lambda^j/j!$, $F(j)=\sum_{i\le j}p_i$, $F(-1)=0$; $U\sim U(0,1)$.
\item $N$: number of comparisons ``$U<F(j)$?'' made by the algorithm.
\end{notacao}

\passo{Algorithm}{Set $X=j$ when $F(j-1)\le U<F(j)$: start with $j=0$, $p=e^{-\lambda}$, $F=p$;
while $U\ge F$, update $p\leftarrow p\lambda/(j+1)$, $j\leftarrow j+1$, $F\leftarrow F+p$. The
recursion holds because}
\begin{equation}\label{q22:rec}
\frac{p_{j+1}}{p_j}=\frac{e^{-\lambda}\lambda^{j+1}/(j+1)!}{e^{-\lambda}\lambda^j/j!}=\frac{\lambda}{j+1}.
\end{equation}

\passo{Correctness}{The intervals $[F(j-1),F(j))$ partition $[0,1)$ and have lengths $p_j$, so}
\begin{equation}\label{q22:ok}
P(X=j)=P\big(F(j-1)\le U<F(j)\big)=F(j)-F(j-1)=p_j .
\end{equation}

\passo{Cost}{The algorithm stops after checking $j=0,1,\dots,X$, so $N=X+1$ and
$E(N)=\lambda+1=3$: the inverse transform is cheap for small $\lambda$ and grows linearly in it.}

\passo{Five values}{With $\lambda=2$: $F(0)=0.135335$, $F(1)=0.406006$, $F(2)=0.676676$,
$F(3)=0.857123$. The seed-22 uniforms $0.366347$, $0.199295$, $0.088558$, $0.653192$, $0.459337$ fall
in $[F(0),F(1))$, $[F(0),F(1))$, $[0,F(0))$, $[F(1),F(2))$, $[F(1),F(2))$, giving $1,1,0,2,2$.}

\codehead
{\small
\begin{verbatim}
def rpois_inv(u, lam):
    j, p = 0, np.exp(-lam); F = p
    while u >= F:
        p *= lam / (j + 1); j += 1; F += p      # recursion (22.1)
    return j, j + 1                             # value, comparisons
rng = np.random.default_rng(22)
[rpois_inv(u, 2.)[0] for u in rng.uniform(size=5)]          # [1, 1, 0, 2, 2]
big = np.array([rpois_inv(u, 2.) for u in rng.uniform(size=100_000)])
big[:, 0].mean(), big[:, 0].var(ddof=1), big[:, 1].mean()   # 1.99556 2.01146 2.99556
# chi-square GOF, classes 0..6, >=7: X2 = 12.44 on 7 df, p = 0.0869
\end{verbatim}
}
Mean and variance are both close to $\lambda=2$, the comparisons average $3=\lambda+1$, and the
goodness-of-fit test does not reject at 5\%.

\intuicao{Every discrete distribution can be simulated by walking up its c.d.f.\ until it passes a
uniform draw; the Poisson recursion makes each step a single multiplication.}
\refbib{Morettin \& Singer (2021), Apêndice B (geração de variáveis discretas, método da
transformação inversa); Greenberg (2008), §5.1.}
\end{sol}

% -----------------------------------------------------------------------------
\questao{MS B.18 --- half-normal by accept--reject, then a normal sample.}
\fonte{Morettin, P.~A. and Singer, J.~M. (2021), \textit{Estatística e Ciência de Dados} (versão
preliminar, abril 2021), Apêndice~B, §B.8 Exercícios, Exercício~18, p.~478 [PDF p.~488]; the book
gives no answer --- reproduced verbatim.}

\begin{otherlanguage}{brazil}
18) Suponha que $Z\sim N(0,1)$ e considere $X=|Z|$. Obtenha a densidade $\pi(x)$ de $X$. Use o
método da aceitação-rejeição para obter uma amostra de $\pi(x)$ com $g(x)=\exp(-x)$. A seguir gere
uma amostra de $Z$, levando em conta que $Z=X$ ou $Z=-X$.
\end{otherlanguage}
\programacao With \texttt{np.random.default\_rng(23)}, draw proposals until $10{,}000$ are
accepted; report the empirical acceptance rate against $1/c$, the mean of $X$ against
$\sqrt{2/\pi}$, the mean and variance of the signed sample $Z$, and Kolmogorov--Smirnov p-values
of $X$ against the half-normal and of $Z$ against $N(0,1)$.
\resolespaco
\begin{sol}
\solhead
\begin{notacao}
\item $\phi,\Phi$: standard normal density and c.d.f.; $g(x)=e^{-x}$, $x\ge0$: proposal (exponential).
\item $c=\sup_x\pi(x)/g(x)$: envelope constant; $U\sim U(0,1)$; $S=\pm1$ with probability $1/2$ each, independent of $X$.
\end{notacao}

\passo{Density of $X$}{For $x\ge0$, $P(X\le x)=P(-x\le Z\le x)=2\Phi(x)-1$; differentiating,}
\begin{equation}\label{q23:pi}
\pi(x)=2\phi(x)=\sqrt{\frac2\pi}\,e^{-x^2/2},\qquad x\ge0 .
\end{equation}

\passo{Envelope}{Dividing \eqref{q23:pi} by $g$:}
\begin{equation}\label{q23:ratio}
\frac{\pi(x)}{g(x)}=\sqrt{\frac2\pi}\,e^{x-x^2/2}.
\end{equation}
The exponent $x-x^2/2$ is maximized at $x=1$ (derivative $1-x$), where it equals $1/2$. Substituting
$x=1$ into \eqref{q23:ratio}:
\begin{equation}\label{q23:c}
c=\sqrt{\frac2\pi}\,e^{1/2}=\sqrt{\frac{2e}{\pi}}=1.315489,\qquad \frac1c=0.760173 .
\end{equation}
Dividing \eqref{q23:ratio} by \eqref{q23:c}, the acceptance probability of a proposal $x$ is
\begin{equation}\label{q23:acc}
\frac{\pi(x)}{c\,g(x)}=e^{x-x^2/2-1/2}=e^{-(x-1)^2/2}.
\end{equation}

\passo{Why accepted draws follow $\pi$}{Draw $x\sim g$ and accept with probability \eqref{q23:acc}.
Then}
\[
P(X\le x,\ \text{accept})=\int_0^xg(t)\frac{\pi(t)}{cg(t)}\,dt=\frac1c\int_0^x\pi(t)\,dt,
\]
and letting $x\to\infty$, $P(\text{accept})=1/c$. Dividing, $P(X\le x\mid\text{accept})=\int_0^x\pi$:
accepted values have density $\pi$, and the number of proposals per accepted value is geometric
with mean $c$.

\passo{Back to $Z$}{Let $Z=SX$. For $z\ge0$, $P(Z\le z)=\frac12+\frac12P(X\le z)=\frac12+\frac12(2\Phi(z)-1)=\Phi(z)$;
for $z<0$, $P(Z\le z)=\frac12P(X\ge-z)=\frac12\big(2-2\Phi(-z)\big)=\Phi(z)$. So $Z\sim N(0,1)$.}

\codehead
{\small
\begin{verbatim}
rng = np.random.default_rng(23); acc, tries = [], 0
while len(acc) < 10_000:
    x, u = rng.exponential(1.0), rng.uniform(); tries += 1
    if u <= np.exp(-(x - 1)**2 / 2): acc.append(x)       # (23.4)
X = np.array(acc); Z = np.where(rng.uniform(size=X.size) <= .5, X, -X)
10_000 / tries, X.mean()             # 0.762195 (1/c = 0.760173)  0.793801 (0.797885)
Z.mean(), Z.var(ddof=1)              # 0.019998  0.990055
kstest(X, halfnorm.cdf).pvalue, kstest(Z, norm.cdf).pvalue    # 0.4935  0.3449
\end{verbatim}
}

\intuicao{The exponential has heavier tails than the half-normal, so a scaled copy covers it
everywhere; the price is that about one proposal in four is thrown away. Symmetry then turns a
half-normal into a normal with one coin flip.}
\refbib{Morettin \& Singer (2021), Apêndice B (método da aceitação-rejeição); Greenberg (2008),
§5.3 (Algorithm 5.2 and the normal-from-Laplace example, pp.~66--69).}
\end{sol}

% -----------------------------------------------------------------------------
\questao{Original math question --- importance sampling for a rare Poisson tail.}
\fonte{Original question (not from a book).}

A hospital's ICU receives $X\sim$ Poisson$(\lambda)$ emergency admissions per night; the
manager needs $p=P(X\ge k)$ for a capacity $k$ well above $\lambda$. Monte Carlo is used with
$G$ draws.
\begin{enumerate}[label=\textbf{(\alph*)},leftmargin=2.4em]
\item For the crude estimator $\hat p=\frac1G\sum_g\1\{X^{(g)}\ge k\}$, $X^{(g)}\sim$ Poisson$(\lambda)$,
derive its mean, variance and relative standard error, and prove that for fixed $G$ the relative
error tends to infinity as $p\to0$.
\item Importance sampling with proposal Poisson$(\mu)$: derive the weight $w_\mu(x)$, prove
unbiasedness, and derive the exact variance of one term in closed form,
$\sigma^2(\mu)=e^{\mu-2\lambda+\lambda^2/\mu}\,P(Y_{\lambda^2/\mu}\ge k)-p^2$ with $Y_\nu\sim$ Poisson$(\nu)$.
\item Prove that $\mu=\lambda$ gives back the crude variance, and that the proposal
$f_\lambda(x)\1\{x\ge k\}/p$ has zero variance. Why can it not be used?
\item Parameter analysis: for $\mu>\lambda$, prove that $w_\mu$ decreases in $x$, deduce
$\sigma^2(\mu)\le p\,(w_\mu(k)-p)$, and show that this bound is minimized at $\mu=k$. Interpret.
\end{enumerate}
\programacao With $\lambda=2$, $k=10$, $G=10^5$ and \texttt{np.random.default\_rng(24)}: compute $p$
exactly; run the crude estimator and the IS estimator with $\mu=10$ (in that order, same generator)
and print estimates with standard errors; tabulate $\sigma^2(\mu)$ for $\mu\in\{2,5,10,15\}$ from
(b), check it by summing $f_\lambda(x)^2/f_\mu(x)$ over $x\ge k$, and find the $\mu$ that minimizes
$\sigma^2(\mu)$ on a grid.
\resolespaco
\begin{sol}
\solhead
\begin{notacao}
\item $f_\nu(x)=e^{-\nu}\nu^x/x!$: Poisson p.m.f.; $Q_k(\nu)=P(Y_\nu\ge k)$; $p=Q_k(\lambda)$.
\item $w_\mu(x)=f_\lambda(x)/f_\mu(x)$: importance weight; $\hat p_\mu=\frac1G\sum_g\1\{X^{(g)}\ge k\}w_\mu(X^{(g)})$, $X^{(g)}\sim f_\mu$.
\item $\sigma^2(\mu)$: variance of one IS term; relative standard error $=\sqrt{\Var(\hat p)}/p$.
\end{notacao}

\passo{(a) Crude}{The indicators are i.i.d.\ Bernoulli$(p)$, so $E\hat p=p$ and $\Var\hat p=p(1-p)/G$.
The relative standard error is}
\begin{equation}\label{q24:rse}
\frac{\sqrt{p(1-p)/G}}{p}=\sqrt{\frac{1-p}{p\,G}}\xrightarrow[p\to0]{}\infty .
\end{equation}
For a relative error of $10\%$ one needs $G=(1-p)/(0.01\,p)$ draws, i.e.\ $G\approx100/p$.

\passo{(b) Weight and unbiasedness}{Dividing the two p.m.f.s,}
\begin{equation}\label{q24:w}
w_\mu(x)=\frac{e^{-\lambda}\lambda^x/x!}{e^{-\mu}\mu^x/x!}=e^{\mu-\lambda}\Big(\frac\lambda\mu\Big)^x .
\end{equation}
Then $E_\mu[\1\{X\ge k\}w_\mu(X)]=\sum_{x\ge k}f_\mu(x)\frac{f_\lambda(x)}{f_\mu(x)}=\sum_{x\ge k}f_\lambda(x)=p$.

\passo{(b) Second moment}{$E_\mu[\1\{X\ge k\}w_\mu^2]=\sum_{x\ge k}f_\lambda(x)^2/f_\mu(x)$, and}
\[
\frac{f_\lambda(x)^2}{f_\mu(x)}=\frac{e^{-2\lambda}\lambda^{2x}/(x!)^2}{e^{-\mu}\mu^x/x!}
=e^{\mu-2\lambda}\frac{(\lambda^2/\mu)^x}{x!}=e^{\mu-2\lambda+\lambda^2/\mu}\,f_{\lambda^2/\mu}(x).
\]
Summing over $x\ge k$ and subtracting $p^2$:
\begin{equation}\label{q24:var}
\sigma^2(\mu)=e^{\mu-2\lambda+\lambda^2/\mu}\,Q_k(\lambda^2/\mu)-p^2 .
\end{equation}

\passo{(c) Two special proposals}{Substituting $\mu=\lambda$ into \eqref{q24:var}:
$e^{\lambda-2\lambda+\lambda}Q_k(\lambda)-p^2=p-p^2=p(1-p)$, the crude variance. For
$q^*(x)=f_\lambda(x)\1\{x\ge k\}/p$, every term is $\1\{x\ge k\}f_\lambda(x)/q^*(x)=p$, a constant: zero
variance. But $q^*$ requires $p$, the unknown being estimated; it only tells us to put the proposal
mass where $\{x\ge k\}$ and in proportion to $f_\lambda$ there.}

\passo{(d) A bound}{For $\mu>\lambda$, $\lambda/\mu<1$, so by \eqref{q24:w} $w_\mu(x)$ decreases in
$x$; on $\{x\ge k\}$, $w_\mu(x)\le w_\mu(k)$. Hence}
\begin{equation}\label{q24:bound}
\sigma^2(\mu)=E_\mu[\1\{X\ge k\}w_\mu^2]-p^2\le w_\mu(k)\,E_\mu[\1\{X\ge k\}w_\mu]-p^2=p\big(w_\mu(k)-p\big).
\end{equation}

\passo{(d) Best tilt}{Taking logs of \eqref{q24:w} at $x=k$, $\log w_\mu(k)=\mu-\lambda+k\log\lambda-k\log\mu$;
differentiating in $\mu$:}
\[
\frac{d\log w_\mu(k)}{d\mu}=1-\frac k\mu,
\]
negative for $\mu<k$, positive for $\mu>k$: the bound \eqref{q24:bound} is minimized at $\mu=k$.
Reading: centre the proposal at the threshold, so that about half the draws land in the rare region,
each weighted down by roughly $w_k(k)$. As $\mu\to\infty$ the weights explode
($e^{\mu}$ in \eqref{q24:var}) and IS becomes worse than crude.

\passo{Numbers}{$\lambda=2$, $k=10$: $p=4.649808\times10^{-5}$; crude relative error with $G=10^5$
is $\sqrt{(1-p)/(pG)}=0.4637$ by \eqref{q24:rse}, and $10\%$ needs $G=2{,}150{,}527$ draws. With $\mu=10$:
$w_{10}(10)=e^{8}(0.2)^{10}=3.0525\times10^{-4}$, bound \eqref{q24:bound} $1.2031\times10^{-8}$, exact
\eqref{q24:var} $\sigma^2(10)=9.9339\times10^{-9}$, relative error $0.006778$. The grid minimizer of
\eqref{q24:var} is $\mu=10.035$, practically $k$. The variance ratio crude/IS is $4680.5$.}

\codehead
{\small
\begin{verbatim}
lam, k, G = 2., 10, 10**5; p = poisson.sf(k - 1, lam)          # 4.649808e-05
var_is = lambda mu: (np.exp(mu - 2*lam + lam**2/mu) * poisson.sf(k - 1, lam**2/mu)
                     - p**2)                                    # (24.3)
# mu = 2, 5, 10, 15: 4.6496e-05  8.4535e-08  9.9339e-09  2.8585e-08
rng = np.random.default_rng(24)
xn = rng.poisson(lam, G); (xn >= k).mean()        # 4.0e-05 (4 hits), se 2.0e-05
xi = rng.poisson(10., G)
w = (xi >= k) * np.exp(10. - lam) * (lam / 10.)**xi
w.mean(), w.std(ddof=1) / np.sqrt(G)              # 4.6844e-05  3.165e-07
\end{verbatim}
}
The crude run saw only 4 events in $10^5$ nights; IS is within about one standard error of the exact
$p$ with a relative error below 1\%.

\intuicao{Rare-event probabilities are hard for plain simulation because almost no draw hits the
event. Importance sampling simulates a world where the event is common and then corrects each hit
by how much likelier it was made; choosing that world well is the whole art.}
\refbib{Greenberg (2008), §5.4 (importance sampling, variance of $\hat g$, choice of $h$);
Morettin \& Singer (2021), Apêndice B.}
\rubrica{
\item (a) mean, variance and \eqref{q24:rse} with the limit: 3 pts.
\item (b) \eqref{q24:w} and unbiasedness: 2 pts; \eqref{q24:var} derived: 3 pts.
\item (c) both special cases and why $q^*$ is unusable: 3 pts.
\item (d) monotone weight, \eqref{q24:bound}, minimization at $\mu=k$, reading: 4 pts.
\item Programming: exact $p$, both estimators with standard errors, variance table: 3 pts.
}
\end{sol}

% -----------------------------------------------------------------------------
\questao{Original conceptual question --- simulation methods (10 points).}
\fonte{Original question (not from a book).}

\textbf{a)} (3 points) In the accept--reject algorithm with $f\le c\,g$, the number of proposals
needed to obtain one accepted draw is:
\begin{alternativas}
\item geometric with mean $1/c$, the acceptance probability.
\item Poisson with mean $c$, the envelope constant.
\item geometric with mean $c$, since each proposal is accepted with probability $1/c$.
\item binomial with parameters $c$ and $1/c$, one trial per unit of envelope.
\end{alternativas}
\gabarito{C}
\begin{sol}
\passo{Why}{C: each proposal is accepted independently with probability $1/c$ (Question~23), so the
count until the first success is geometric with mean $c$. A swaps the probability and the mean;
B and D are not distributions of a waiting time for a first success.}
\end{sol}

\textbf{b)} (3 points) The Monte Carlo estimator of $E[h(X)]$ from $G$ i.i.d.\ draws has standard error:
\begin{alternativas}
\item proportional to $1/G$, because averaging divides by $G$.
\item $\mathrm{sd}(h(X))/\sqrt G$, whatever the dimension of $X$.
\item $\mathrm{sd}(h(X))/\sqrt G$ only when $X$ is scalar and $h$ is bounded.
\item growing exponentially with the dimension of $X$, like quadrature.
\end{alternativas}
\gabarito{B}
\begin{sol}
\passo{Why}{B: $\Var(\bar h)=\Var(h)/G$ for i.i.d.\ draws with finite variance, in any dimension:
this is why Monte Carlo beats grids in high dimension. A forgets the square root; C adds
conditions that are not needed (finite variance suffices); D describes product-rule quadrature.}
\end{sol}

\textbf{c)} (2 points, T or F) An importance-sampling estimator is unbiased, yet a proposal with
lighter tails than the target can give it infinite variance.
\gabarito{T}
\begin{sol}
\passo{Why}{Unbiasedness only needs $h>0$ where $f>0$; the second moment is
$\int h^2f^2/g$, which diverges when $f/g$ grows too fast in the tails (Greenberg §5.4).}
\end{sol}

\textbf{d)} (2 points, T or F) Accept--reject cannot be used unless the normalizing constant of the
target density is known.
\gabarito{F}
\begin{sol}
\passo{Why}{If $f=\kappa r$ with $r$ known, choose $c$ with $r\le cg$ and accept with probability
$r/(cg)$: accepted values still follow $f$ (Greenberg p.~67). The constant only affects the
acceptance rate. MH shares this property.}
\end{sol}

\programacao Check (a) on the half-normal sampler of Question~23: with
\texttt{np.random.default\_rng(25)} draw $2\times10^6$ exponential proposals and uniforms, mark the
accepted ones, take the first $100{,}000$ acceptances, and compare the mean and variance of the
number of proposals per acceptance with $c$ and $c(c-1)$.
\begin{sol}
\codehead
{\small
\begin{verbatim}
rng = np.random.default_rng(25); x = rng.exponential(1, 2_000_000)
acc = rng.uniform(size=x.size) <= np.exp(-(x - 1)**2 / 2)
waits = np.diff(np.r_[-1, np.flatnonzero(acc)[:100_000]])
waits.mean(), waits.var(ddof=1)     # 1.31879 (c = 1.315489)  0.42047 (c(c-1) = 0.41502)
\end{verbatim}
}
\end{sol}

% =============================================================================
\bloco{6}{The bootstrap}{I 5.2, 5.4; original}

\questao{I 5.2 --- probability that an observation is in a bootstrap sample.}
\fonte{James, G., Witten, D., Hastie, T., Tibshirani, R. and Taylor, J. (2023), \textit{An
Introduction to Statistical Learning with Applications in Python}, Springer, Chapter~5, §5.4
Exercises (Conceptual), Exercise~2, pp.~224--225 [PDF pp.~232--233]; the book gives no answer ---
reproduced verbatim (the margin index entry ``\texttt{np.empty()}'' is omitted).}

2. We will now derive the probability that a given observation is part of a bootstrap sample.
Suppose that we obtain a bootstrap sample from a set of $n$ observations.
\begin{enumerate}[label=(\alph*),leftmargin=2.4em]
\item What is the probability that the first bootstrap observation is \emph{not} the $j$th
observation from the original sample? Justify your answer.
\item What is the probability that the second bootstrap observation is \emph{not} the $j$th
observation from the original sample?
\item Argue that the probability that the $j$th observation is \emph{not} in the bootstrap sample
is $(1-1/n)^n$.
\item When $n=5$, what is the probability that the $j$th observation is in the bootstrap sample?
\item When $n=100$, what is the probability that the $j$th observation is in the bootstrap sample?
\item When $n=10{,}000$, what is the probability that the $j$th observation is in the bootstrap
sample?
\item Create a plot that displays, for each integer value of $n$ from 1 to 100,000, the probability
that the $j$th observation is in the bootstrap sample. Comment on what you observe.
\item We will now investigate numerically the probability that a bootstrap sample of size $n=100$
contains the $j$th observation. Here $j=4$. We first create an array \texttt{store} with values that
will subsequently be overwritten using the function \texttt{np.empty()}. We then repeatedly create
bootstrap samples, and each time we record whether or not the fifth observation is contained in the
bootstrap sample.
{\small
\begin{verbatim}
rng = np.random.default_rng(10)
store = np.empty(10000)
for i in range(10000):
    store[i] = np.sum(rng.choice(100, replace=True) == 4)
          > 0
np.mean(store)
\end{verbatim}
}
Comment on the results obtained.
\end{enumerate}
\programacao Let $D$ be the number of \emph{distinct} original observations in a bootstrap sample of
size $n=100$. Derive $E(D)$ and $\Var(D)$ from indicator variables, then check them with $10{,}000$
bootstrap samples drawn with \texttt{np.random.default\_rng(26)}.
\resolespaco
\begin{sol}
\solhead
\begin{notacao}
\item $n$: sample size; $j\in\{1,\dots,n\}$: a fixed original observation.
\item $Z_1,\dots,Z_n$: indices of the bootstrap draws, i.i.d.\ uniform on $\{1,\dots,n\}$ (sampling
  with replacement).
\item $p_n=P(j\text{ is in the bootstrap sample})$; $A_j=\1\{j\text{ is absent}\}$; $D=n-\sum_jA_j$.
\end{notacao}

\passo{(a)}{Each draw picks one of the $n$ original observations with equal probability $1/n$, so}
\begin{equation}\label{q26:a}
P(Z_1\neq j)=1-\frac1n .
\end{equation}

\passo{(b)}{Sampling is with replacement, so $Z_2$ has the same distribution as $Z_1$ and is
independent of it:}
\begin{equation}\label{q26:b}
P(Z_2\neq j)=1-\frac1n .
\end{equation}

\passo{(c)}{$j$ is absent only if every draw misses it; by independence and \eqref{q26:a}--\eqref{q26:b}
applied to all $n$ draws:}
\begin{equation}\label{q26:c}
P(j\notin\text{sample})=\prod_{m=1}^nP(Z_m\neq j)=\Big(1-\frac1n\Big)^n .
\end{equation}

\passo{(d)--(f)}{Taking the complement of \eqref{q26:c}:}
\begin{equation}\label{q26:p}
p_n=1-\Big(1-\frac1n\Big)^n .
\end{equation}
Substituting into \eqref{q26:p}:
\begin{align*}
n=5:&\quad p_5=1-(4/5)^5=1-0.32768=0.67232,\\
n=100:&\quad p_{100}=1-(0.99)^{100}=1-0.366032=0.633968,\\
n=10{,}000:&\quad p_{10000}=1-(0.9999)^{10000}=1-0.367861=0.632139.
\end{align*}

\passo{(g) Plot and limit}{Since $n\log(1-1/n)=n\big(-\tfrac1n-\tfrac1{2n^2}-\dots\big)\to-1$,}
\begin{equation}\label{q26:lim}
\lim_{n\to\infty}p_n=1-e^{-1}=0.632121 .
\end{equation}
\begin{center}
\begin{tikzpicture}
\begin{semilogxaxis}[width=0.72\linewidth,height=5cm,xlabel={$n$ (log scale)},
  ylabel={$p_n$},ymin=0.6,ymax=1.02,xmin=1,xmax=100000,grid=major,
  every axis plot/.append style={thick}]
\addplot[blue!60!black,mark=*,mark size=1.3pt] coordinates {(1,1.0000) (2,0.7500) (3,0.7037)
  (5,0.6723) (10,0.6513) (20,0.6415) (50,0.6358) (100,0.6340) (1000,0.6323) (10000,0.6321)
  (100000,0.6321)};
\addplot[red!70!black,dashed,domain=1:100000] {0.632121};
\end{semilogxaxis}
\end{tikzpicture}
\end{center}
$p_n$ falls from 1 at $n=1$ and flattens almost immediately at $1-e^{-1}\approx0.632$ (dashed): for
$n\ge100$ the curve is indistinguishable from the limit.

\passo{(h) The printed code}{\texttt{rng.choice(100, replace=True)} without \texttt{size} returns a
\emph{single} integer, so each iteration checks whether one draw equals 4, an event of probability
$1/100$. Run as printed, \texttt{np.mean(store)} returns $0.0089$, not a bootstrap probability. The
intended code draws a whole bootstrap sample, \texttt{rng.choice(100, size=100, replace=True)},
and then returns $0.6362$, close to $p_{100}=0.633968$ from (e) (Monte Carlo standard error
$\sqrt{0.634\cdot0.366/10000}\approx0.005$). Index 4 is the fifth observation because Python
counts from 0, which is why the text says ``fifth''.}

\passo{Added: distinct observations}{$E(A_j)=a=(1-1/n)^n$ by \eqref{q26:c}; for $i\ne j$ both are
absent iff every draw avoids two indices, $E(A_iA_j)=b=(1-2/n)^n$. Hence}
\begin{equation}\label{q26:D}
E(D)=n(1-a),\qquad \Var(D)=\Var\Big(\sum_jA_j\Big)=na(1-a)+n(n-1)(b-a^2).
\end{equation}
For $n=100$: $E(D)=63.396766$ and $\mathrm{sd}(D)=3.120912$.

\codehead
{\small
\begin{verbatim}
rng = np.random.default_rng(10); store = np.empty(10000)
for i in range(10000):
    store[i] = np.sum(rng.choice(100, size=100, replace=True) == 4) > 0
store.mean()                                   # 0.6362 (printed code: 0.0089)
rng = np.random.default_rng(26)
D = [np.unique(rng.choice(100, size=100, replace=True)).size for _ in range(10000)]
np.mean(D), np.std(D, ddof=1)                  # 63.3994  3.1583  (exact 63.3968, 3.1209)
\end{verbatim}
}

\intuicao{A bootstrap sample of size $n$ contains only about 63.2\% of the distinct original
observations; the other 36.8\% are ``out of bag''. This is why bootstrap estimates of test error
that reuse the training points are optimistic, and why out-of-bag points can serve as a free
validation set.}
\refbib{ISLP (2023), §5.2 (the bootstrap, pp.~212--215).}
\end{sol}

% -----------------------------------------------------------------------------
\questao{I 5.4 --- bootstrap standard deviation of a prediction.}
\fonte{James, G., Witten, D., Hastie, T., Tibshirani, R. and Taylor, J. (2023), \textit{An
Introduction to Statistical Learning with Applications in Python}, Springer, Chapter~5, §5.4
Exercises (Conceptual), Exercise~4, p.~225 [PDF p.~233]; the book gives no answer --- reproduced
verbatim.}

4. Suppose that we use some statistical learning method to make a prediction for the response $Y$
for a particular value of the predictor $X$. Carefully describe how we might estimate the standard
deviation of our prediction.
\programacao Fit the Poisson regression $\log\mu_t=\beta_0+\beta_1x_t$, $x_t=(t-56.5)/10$, to the coal
series ($t=1,\dots,112$ for 1851--1962) with \texttt{statsmodels} GLM. Predict $\mu$ for 1900
($t=50$). Compute the delta-method standard error of $\hat\mu(1900)$ and the bootstrap standard
deviation from $B=2{,}000$ pairs-bootstrap refits (\texttt{np.random.default\_rng(27)}), with the
95\% percentile interval. Compare the two using the Pearson dispersion statistic.
\resolespaco
\begin{sol}
\solhead
\begin{notacao}
\item $(x_i,y_i)$, $i=1,\dots,n$: training data; $\hat f$: fitted method; $x_0$: the predictor value; $\hat y_0=\hat f(x_0)$.
\item $B$: number of bootstrap samples; $\hat y_0^{*b}$: prediction refitted on bootstrap sample $b$.
\item Coal model: $\hat\beta=(\hat\beta_0,\hat\beta_1)$, $\hat V$: its estimated covariance; $\hat\mu_0=\exp(x_0'\hat\beta)$.
\end{notacao}

\passo{The procedure}{(1) Draw $n$ indices with replacement from $\{1,\dots,n\}$ and form the
bootstrap sample $\{(x_{i^*},y_{i^*})\}$. (2) Refit the \emph{whole} method on it (including any
tuning or variable selection) and compute $\hat y_0^{*b}=\hat f^{*b}(x_0)$. (3) Repeat for
$b=1,\dots,B$. (4) Estimate the standard deviation of the prediction by}
\begin{equation}\label{q27:sd}
\widehat{\mathrm{SD}}(\hat y_0)=\sqrt{\frac1{B-1}\sum_{b=1}^B\big(\hat y_0^{*b}-\bar y_0^{*}\big)^2},\qquad
\bar y_0^*=\frac1B\sum_b\hat y_0^{*b}.
\end{equation}

\passo{What it measures}{\eqref{q27:sd} is the sampling variability of the \emph{fitted value}
$\hat f(x_0)$, the analogue of $\mathrm{SE}(\hat\alpha)$ in ISLP §5.2. For the error of predicting
a new response $Y_0=f(x_0)+\varepsilon$ one must add the irreducible noise:
$\mathrm{SD}(Y_0-\hat y_0)\approx\sqrt{\widehat{\mathrm{SD}}(\hat y_0)^2+\widehat{\Var}(\varepsilon\mid x_0)}$.
Resampling pairs keeps the joint behaviour of $(X,Y)$ and does not assume the model is correct;
resampling residuals instead assumes a correct mean and homoskedastic errors.}

\passo{Delta method (benchmark)}{$\hat\mu_0=g(\hat\beta)$ with $g(\beta)=\exp(x_0'\beta)$ and
gradient $\nabla g=\hat\mu_0x_0$, so}
\begin{equation}\label{q27:delta}
\mathrm{SE}(\hat\mu_0)\approx\hat\mu_0\sqrt{x_0'\hat Vx_0}.
\end{equation}

\passo{Coal numbers}{$\hat\beta_0=0.363109$ (s.e.\ $0.084922$), $\hat\beta_1=-0.183776$ (s.e.\
$0.024728$) per decade; $\hat\mu(1900)=1.620223$; \eqref{q27:delta} gives $0.125930$; the bootstrap
\eqref{q27:sd} gives $0.134530$ with percentile interval $[1.3633,\ 1.8748]$. The Pearson
dispersion is $\hat\phi=1.133236$: the counts are slightly more variable than a Poisson with a
linear trend allows, and $0.125930\sqrt{1.133236}=0.134057$, almost exactly the bootstrap value.
For a new count in 1900 the predictive SD is about $\sqrt{1.620223+0.134530^2}=1.2800$, dominated by
Poisson noise.}

\codehead
{\small
\begin{verbatim}
t = np.arange(1, 113); X = sm.add_constant((t - 56.5) / 10)
fit = sm.GLM(y, X, family=sm.families.Poisson()).fit()
x0 = np.array([1., (50 - 56.5) / 10])
mu0 = np.exp(x0 @ fit.params)                               # 1.620223
se = mu0 * np.sqrt(x0 @ fit.cov_params() @ x0)              # 0.125930
rng = np.random.default_rng(27); pred = np.empty(2000)
for b in range(2000):
    i = rng.integers(0, 112, 112)
    pred[b] = np.exp(x0 @ sm.GLM(y[i], X[i], family=sm.families.Poisson()).fit().params)
pred.std(ddof=1), np.quantile(pred, [.025, .975])           # 0.134530 [1.3633 1.8748]
\end{verbatim}
}

\intuicao{The bootstrap answers ``how much would my prediction move if I had drawn another sample
of the same size?'' by drawing those samples from the data itself. It agrees with the model-based
formula when the model is right, and corrects it (here for mild overdispersion) when it is not.}
\refbib{ISLP (2023), §5.2 (pp.~212--215) and the §5.3.3 lab (\texttt{boot\_SE}); Wooldridge (2016),
§17.3 (Poisson regression and overdispersion).}
\end{sol}

% -----------------------------------------------------------------------------
\questao{Original math question --- nonparametric versus parametric bootstrap for a Poisson mean.}
\fonte{Original question (not from a book).}

Let $x_1,\dots,x_n$ be counts, $\bar x$ their mean and $\hat\sigma^2=\frac1n\sum_i(x_i-\bar x)^2$.
\begin{enumerate}[label=\textbf{(\alph*)},leftmargin=2.4em]
\item Nonparametric bootstrap ($X_i^*$ i.i.d.\ from the empirical distribution): find
$E^*(\overline X^*)$ and $\Var^*(\overline X^*)$ exactly, without simulation.
\item Parametric bootstrap ($X_i^*$ i.i.d.\ Poisson$(\bar x)$): find $\Var^*(\overline X^*)$. If the
$X_i$ are i.i.d.\ with mean $\theta$ and variance $\sigma^2$, prove $E(\hat\sigma^2)=\frac{n-1}n\sigma^2$
and compare the expectations of the two bootstrap variances with the true $\Var(\overline X)$, both
when the Poisson model holds and when $\sigma^2=\phi\theta$ with $\phi\ne1$.
\item For $\psi=e^{-\theta}$ and $\hat\psi=e^{-\bar x}$, find the parametric-bootstrap bias estimate
$\hat b=E^*(e^{-\overline X^*})-e^{-\bar x}$ in closed form and the bias-corrected estimator.
\item Parameter analysis: prove $\hat b>0$, give its leading term in $1/n$, and say how it moves with
$\bar x$. Interpret.
\end{enumerate}
\programacao For the coal counts 1851--1891 ($n=41$): print $\bar x$, $\hat\sigma^2$, the two exact
bootstrap variances, $\hat\sigma^2/\bar x$, $\hat\psi$, $\hat b$, its leading-term approximation and
the corrected estimate. Check them with $B=20{,}000$ resamples of each kind
(\texttt{np.random.default\_rng(28)}).
\resolespaco
\begin{sol}
\solhead
\begin{notacao}
\item $E^*,\Var^*$: moments under the bootstrap distribution, data held fixed.
\item $\theta=E(X)$, $\sigma^2=\Var(X)$, $\phi=\sigma^2/\theta$: dispersion ($\phi=1$ for Poisson).
\item $\psi=e^{-\theta}$, $\hat\psi=e^{-\bar x}$, $\hat b$: bootstrap bias estimate, $\tilde\psi=\hat\psi-\hat b$.
\end{notacao}

\passo{(a) One draw}{$X^*$ takes each value $x_i$ with probability $1/n$, so}
\begin{equation}\label{q28:one}
E^*(X^*)=\frac1n\sum_ix_i=\bar x,\qquad \Var^*(X^*)=\frac1n\sum_i(x_i-\bar x)^2=\hat\sigma^2 .
\end{equation}
The $n$ draws are i.i.d., so by \eqref{q28:one}
\begin{equation}\label{q28:np}
E^*(\overline X^*)=\bar x,\qquad \Var^*(\overline X^*)=\frac{\hat\sigma^2}{n}.
\end{equation}

\passo{(b) Parametric}{$\Var^*(X^*)=\bar x$ for a Poisson$(\bar x)$, so}
\begin{equation}\label{q28:par}
\Var^*(\overline X^*)=\frac{\bar x}{n}.
\end{equation}

\passo{(b) Expectation of $\hat\sigma^2$}{Using $\sum_i(X_i-\overline X)^2=\sum_iX_i^2-n\overline X^2$
and $E(X^2)=\sigma^2+\theta^2$, $E(\overline X^2)=\sigma^2/n+\theta^2$:}
\[
E\Big[\sum_i(X_i-\overline X)^2\Big]=n(\sigma^2+\theta^2)-n\Big(\frac{\sigma^2}n+\theta^2\Big)=(n-1)\sigma^2 .
\]
Dividing by $n$:
\begin{equation}\label{q28:Es2}
E(\hat\sigma^2)=\frac{n-1}{n}\sigma^2 .
\end{equation}

\passo{(b) Comparison}{The target is $\Var(\overline X)=\sigma^2/n=\phi\theta/n$. By \eqref{q28:np} and
\eqref{q28:Es2}, $E[\hat\sigma^2/n]=\frac{n-1}n\cdot\frac{\sigma^2}n$: biased down by the factor
$(n-1)/n$, whatever $\phi$. By \eqref{q28:par}, $E[\bar x/n]=\theta/n=\Var(\overline X)/\phi$: exact when
the Poisson model holds ($\phi=1$), too small by $1/\phi$ under overdispersion ($\phi>1$) and too
large under underdispersion ($\phi<1$). The nonparametric bootstrap is robust to the variance model;
the parametric one is more precise when the model is right.}

\passo{(c) Bias estimate}{Under the parametric bootstrap $n\overline X^*\sim$ Poisson$(n\bar x)$, and
$e^{-\overline X^*}=(e^{-1/n})^{n\overline X^*}$. Substituting $s=e^{-1/n}$ into the generating
function \eqref{q4:pgf} with $n\theta$ replaced by $n\bar x$:}
\begin{equation}\label{q28:Estar}
E^*(e^{-\overline X^*})=\exp\{n\bar x(e^{-1/n}-1)\},\qquad
\hat b=\exp\{n\bar x(e^{-1/n}-1)\}-e^{-\bar x},
\end{equation}
and the corrected estimator is
\begin{equation}\label{q28:bc}
\tilde\psi=\hat\psi-\hat b=2e^{-\bar x}-\exp\{n\bar x(e^{-1/n}-1)\}.
\end{equation}

\passo{(d) Sign and order}{By \eqref{q4:tan}, $n\bar x(e^{-1/n}-1)>-\bar x$ (for $\bar x>0$), so
$\hat b>0$: the bootstrap detects the upward Jensen bias of Question~4. Expanding,
$n(e^{-1/n}-1)=-1+\frac1{2n}-\frac1{6n^2}+\dots$, so the exponent in \eqref{q28:Estar} is
$-\bar x+\frac{\bar x}{2n}+O(n^{-2})$ and}
\begin{equation}\label{q28:lead}
\hat b=e^{-\bar x}\Big(e^{\bar x/(2n)+O(n^{-2})}-1\Big)=\frac{\bar x\,e^{-\bar x}}{2n}+O(n^{-2}).
\end{equation}
The leading term vanishes at rate $1/n$; as a function of $\bar x$ it is proportional to
$\bar xe^{-\bar x}$, whose derivative $(1-\bar x)e^{-\bar x}$ is positive for $\bar x<1$ and negative
for $\bar x>1$: the bias is largest, relative to $n$, when the mean is about one event.

\passo{(d) Reading}{The bootstrap estimates the bias of a nonlinear plug-in estimator by repeating
the plug-in inside the fitted model; here it reproduces the exact Jensen gap to first order. The
correction is worth making only when $\hat b$ is large relative to the standard error.}

\passo{Coal numbers}{$n=41$, $\bar x=3.097561$, $\hat\sigma^2=2.478287$ ($s^2=2.540244$):
\eqref{q28:np} $=2.478287/41=0.060446$; \eqref{q28:par} $=3.097561/41=0.075550$; $\hat\sigma^2/\bar x=0.800077$:
these years are \emph{under}dispersed, so the Poisson bootstrap is the larger one.
$\hat\psi=e^{-3.097561}=0.045159$; \eqref{q28:Estar}: $\exp\{127(e^{-1/41}-1)\}=0.046883$, so
$\hat b=0.046883-0.045159=0.001724$; \eqref{q28:lead}: $3.097561\cdot0.045159/82=0.001706$;
\eqref{q28:bc}: $\tilde\psi=0.043435$.}

\codehead
{\small
\begin{verbatim}
x = y[:41]; n, xb, s2h = 41, x.mean(), x.var(ddof=0)
s2h / n, xb / n, s2h / xb              # 0.060446  0.075550  0.800077
Estar = np.exp(n * xb * (np.exp(-1 / n) - 1))
Estar - np.exp(-xb), 2 * np.exp(-xb) - Estar        # 0.001724  0.043435
rng = np.random.default_rng(28)
xnp = rng.choice(x, size=(20_000, n), replace=True).mean(1)
xpar = rng.poisson(xb, size=(20_000, n)).mean(1)
xnp.var(ddof=1), xpar.var(ddof=1)                   # 0.060686  0.076385
np.exp(-xpar).mean() - np.exp(-xb)                  # 0.001655
\end{verbatim}
}

\intuicao{The nonparametric bootstrap uses the data's own spread, the parametric one the spread
the model implies. When they disagree, the data are telling you the variance model is off.}
\refbib{ISLP (2023), §5.2; Greenberg (2008), §5.6 (using simulated output); Question~4 for the
exact bias of $\hat\psi$.}
\rubrica{
\item (a) \eqref{q28:one}--\eqref{q28:np}: 3 pts.
\item (b) \eqref{q28:par}: 1 pt; \eqref{q28:Es2} proved: 2 pts; comparison with $\phi$: 3 pts.
\item (c) \eqref{q28:Estar}--\eqref{q28:bc}: 3 pts.
\item (d) sign, \eqref{q28:lead}, dependence on $\bar x$, reading: 3 pts.
\item Programming: exact values and simulation check: 3 pts.
}
\end{sol}

% -----------------------------------------------------------------------------
\questao{Original conceptual question --- the bootstrap (10 points).}
\fonte{Original question (not from a book).}

\textbf{a)} (3 points) For large $n$, the fraction of the original observations that appear at least
once in a bootstrap sample of size $n$ is about:
\begin{alternativas}
\item $0.5$, since each observation is drawn or not with equal chance.
\item $e^{-1}\approx0.368$, the probability that an observation is never drawn.
\item $1$, because the bootstrap sample has the same size as the data.
\item $1-e^{-1}\approx0.632$, the complement of never being drawn.
\end{alternativas}
\gabarito{D}
\begin{sol}
\passo{Why}{D: \eqref{q26:lim}. B is the out-of-bag fraction; A has no derivation; C ignores
repetitions (Question~26 shows $E(D)=63.4$ for $n=100$).}
\end{sol}

\textbf{b)} (3 points) As $B\to\infty$, the nonparametric bootstrap standard error of $\bar x$ equals:
\begin{alternativas}
\item $\hat\sigma/\sqrt n$, with $\hat\sigma^2$ computed with divisor $n$.
\item $s/\sqrt n$, with $s^2$ computed with divisor $n-1$.
\item $\sqrt{\bar x/n}$, the Poisson standard error.
\item $0$, because every bootstrap sample has mean $\bar x$.
\end{alternativas}
\gabarito{A}
\begin{sol}
\passo{Why}{A: \eqref{q28:np}. B is the usual unbiased-variance formula, not what resampling gives;
C is the parametric bootstrap \eqref{q28:par}; D confuses $E^*(\overline X^*)=\bar x$ with each
resample's mean.}
\end{sol}

\textbf{c)} (2 points, T or F) For a fixed sample, the bootstrap standard error converges to the true
standard error of the estimator as the number of bootstrap replications $B$ grows.
\gabarito{F}
\begin{sol}
\passo{Why}{As $B\to\infty$ it converges to the \emph{ideal} bootstrap value computed from the
empirical distribution (e.g.\ $\hat\sigma/\sqrt n$), which still differs from the true value; only
$n\to\infty$ removes that gap.}
\end{sol}

\textbf{d)} (2 points, T or F) With overdispersed counts, a parametric Poisson bootstrap understates
the standard error of $\bar x$.
\gabarito{T}
\begin{sol}
\passo{Why}{By (b) of Question~28, $E[\bar x/n]=\Var(\overline X)/\phi<\Var(\overline X)$ when $\phi>1$.}
\end{sol}

\programacao Illustrate (b)--(c) on the coal counts 1851--1891: with \texttt{np.random.default\_rng(29)}
compute the bootstrap standard error of $\bar x$ with $B=100$, $1{,}000$ and $10{,}000$ resamples and
compare them with the $B\to\infty$ value $\hat\sigma/\sqrt n$.
\begin{sol}
\codehead
{\small
\begin{verbatim}
rng = np.random.default_rng(29); x = y[:41]
for B in (100, 1000, 10000):
    print(B, rng.choice(x, size=(B, 41), replace=True).mean(1).std(ddof=1))
# 100 0.259285   1000 0.242725   10000 0.244059   ; sigma_hat/sqrt(n) = 0.245858
\end{verbatim}
}
Increasing $B$ only removes the resampling noise around $0.245858$; no $B$ can tell how far that
value is from the true standard error.
\end{sol}

% =============================================================================
\bloco{7}{MCMC: invariance, Metropolis--Hastings, accept--reject versus MH}{G 6.8, 7.4, 7.6; original}

\questao{G 6.8 --- an invariant density for a continuous-state kernel.}
\fonte{Greenberg, E. (2008), \textit{Introduction to Bayesian Econometrics}, 1st ed., Cambridge
University Press, Chapter~6, §6.5 Exercises, Exercise~6.8, p.~89 [PDF p.~104]; the book gives no
answer --- reproduced verbatim.}

6.8 (Bhattacharya and Waymire, 1990, p.~216) Show that
\[
\pi(y)=2(1-y),\qquad 0\le y\le1
\]
is an invariant density for the transition kernel
\[
p(x,y)=\begin{cases}e^{1-x}, & \text{if } y<x,\\ e^{1-x}-e^{y-x}, & \text{if } y>x.\end{cases}
\]
\programacao (1) Check by quadrature that $\int_0^1p(x,y)\,dy=1$ at $x=0.3,0.7$ and that
$\int_0^1\pi(x)p(x,y)\,dx=\pi(y)$ at $y=0.1,0.5,0.9$. (2) Derive the c.d.f.\ of $p(x,\cdot)$, simulate the
chain by inversion (\texttt{brentq} on the branch $y>x$) from $x_0=0.5$ with
\texttt{np.random.default\_rng(30)}, discard 1{,}000 draws, keep 20{,}000, and compare the sample mean
and $P(Y<0.5)$ with their values under $\pi$; report the batch-means n.s.e.\ and the lag-1
autocorrelation.
\resolespaco
\begin{sol}
\solhead
\begin{notacao}
\item State space $[0,1]$; $p(x,\cdot)$: density of the next state given the current state $x$.
\item $\pi$ is invariant if $\int_0^1\pi(x)p(x,y)\,dx=\pi(y)$ for all $y$ (Greenberg §6.3).
\item $F_x(y)=\int_0^yp(x,t)\,dt$: c.d.f.\ of the next state.
\end{notacao}

\passo{$p(x,\cdot)$ is a density}{For $y>x$, $e^{1-x}-e^{y-x}=e^{-x}(e-e^y)>0$ since $y<1$. Integrating
over the two branches:}
\begin{equation}\label{q30:norm}
\int_0^1p(x,y)\,dy=\int_0^xe^{1-x}dy+\int_x^1\big(e^{1-x}-e^{y-x}\big)dy=xe^{1-x}+(1-x)e^{1-x}-e^{-x}(e-e^{x})=1 .
\end{equation}

\passo{Split the invariance integral}{For fixed $y$, the points $x<y$ use the branch ``$y>x$'' and
the points $x>y$ use ``$y<x$'':}
\begin{equation}\label{q30:split}
\int_0^1\pi(x)p(x,y)\,dx=\int_0^y2(1-x)\big(e^{1-x}-e^{y-x}\big)dx+\int_y^12(1-x)e^{1-x}dx .
\end{equation}
Collecting the two $e^{1-x}$ terms into one integral over $[0,1]$:
\begin{equation}\label{q30:two}
\int_0^1\pi(x)p(x,y)\,dx=\int_0^12(1-x)e^{1-x}dx-e^{y}\int_0^y2(1-x)e^{-x}dx .
\end{equation}

\passo{First integral}{With $u=1-x$ and $\frac{d}{du}[(u-1)e^u]=ue^u$:}
\[
\int_0^12(1-x)e^{1-x}dx=\int_0^12ue^u\,du=2\big[(u-1)e^u\big]_0^1=2\,(0-(-1))=2 .
\]

\passo{Second integral}{Since $\frac{d}{dx}[xe^{-x}]=(1-x)e^{-x}$,}
\[
\int_0^y2(1-x)e^{-x}dx=2\big[xe^{-x}\big]_0^y=2ye^{-y}.
\]

\passo{Conclusion}{Substituting both results into \eqref{q30:two}:}
\[
\int_0^1\pi(x)p(x,y)\,dx=(2)-e^y\big(2ye^{-y}\big)=2-2y=\pi(y). \qquad\blacksquare
\]

\passo{Simulating the kernel}{From \eqref{q30:norm}'s pieces, the c.d.f.\ of the next state is}
\begin{equation}\label{q30:cdf}
F_x(y)=\begin{cases}ye^{1-x}, & y\le x,\\ ye^{1-x}-e^{y-x}+1, & y>x,\end{cases}
\end{equation}
continuous at $y=x$ and equal to 1 at $y=1$. Given $u\sim U(0,1)$: if $u\le xe^{1-x}$, $y=ue^{x-1}$;
otherwise solve $F_x(y)=u$ on $(x,1)$. Under $\pi$: $E(Y)=\int_0^12y(1-y)dy=1/3$ and
$P(Y<0.5)=2(0.5)-0.5^2=0.75$.

\codehead
{\small
\begin{verbatim}
def draw(x, u):
    if u <= x * np.exp(1 - x): return u * np.exp(x - 1)
    return brentq(lambda y: y*np.exp(1 - x) - np.exp(y - x) + 1 - u, x, 1.)  # (30.4)
rng = np.random.default_rng(30); x = 0.5; ch = []
for g in range(21_000):
    x = draw(x, rng.uniform())
    if g >= 1000: ch.append(x)
ch = np.array(ch); ch.mean(), np.mean(ch < .5)      # 0.334981 (1/3)  0.74745 (0.75)
ch.reshape(50, 400).mean(1).std(ddof=1) / np.sqrt(50)   # n.s.e. 0.001868
np.corrcoef(ch[:-1], ch[1:])[0, 1]                  # lag-1 autocorrelation 0.0768
\end{verbatim}
}
Quadrature confirms \eqref{q30:norm} and the invariance identity to $10^{-9}$; the simulated mean is
within one n.s.e.\ of $1/3$.

\intuicao{Invariance means that if today's state is drawn from $\pi$, tomorrow's is too: the kernel
moves mass around but leaves its profile unchanged. MCMC is the art of building kernels with a
prescribed invariant law; this exercise checks one by direct integration.}
\refbib{Greenberg (2008), §6.3 (continuous state spaces, invariant distributions, pp.~84--87).}
\end{sol}

% -----------------------------------------------------------------------------
\questao{G 7.4 --- the acceptance probability of Algorithm 7.5.}
\fonte{Greenberg, E. (2008), \textit{Introduction to Bayesian Econometrics}, 1st ed., Cambridge
University Press, Chapter~7, §7.5 Exercises, Exercise~7.4, p.~106 [PDF p.~121]; the book gives no
answer --- reproduced verbatim. Algorithm 7.5 and Equation (7.6), to which it refers, are reproduced
from §7.2, pp.~98--99 [PDF pp.~113--114].}

\emph{Referred to (Greenberg 2008, p.~98):}
\[
\alpha(x,y)=\begin{cases}\min\Big\{\dfrac{f(y)q(y,x)}{f(x)q(x,y)},1\Big\}, & \text{if } f(x)q(x,y)\neq0,\\[8pt] 0, & \text{otherwise.}\end{cases}\tag{7.6}
\]
\emph{Referred to (Greenberg 2008, p.~99):} ``\textbf{Algorithm 7.5: MH for Beta(3, 4) with U(0, 1) proposal}
\begin{itemize}[leftmargin=2.6em]
\item[7.1] Set $x^{(0)}$ equal to a number between zero and one.
\item[7.2] At the $g$th iteration, generate $U_1$ and $U_2$ from $U(0,1)$.
\item[7.3] If
\[
U_1\le\alpha(x^{(g-1)},U_2)=\frac{U_2^2(1-U_2)^3}{(x^{(g-1)})^2(1-x^{(g-1)})^3},
\]
set $x^{(g)}=U_2$. Otherwise set $x^{(g)}=x^{(g-1)}$.
\item[7.4] Go to 2 and continue until the desired number of iterations is obtained.''
\end{itemize}

7.4 Derive $\alpha(x,y)$ for Algorithm 7.5.
\programacao Run Algorithm 7.5 from $x^{(0)}=0.5$ with \texttt{np.random.default\_rng(31)}, discard
500 iterations and keep 5{,}000 (Greenberg reports an acceptance rate of 0.57 and a mean of 0.4296).
Report your acceptance rate and mean against the true mean $3/7$, and compute the exact long-run
acceptance probability $\iint f(x)\min\{f(y)/f(x),1\}\,dx\,dy$ by two-dimensional quadrature.
\resolespaco
\begin{sol}
\solhead
\begin{notacao}
\item Target $f(x)=x^2(1-x)^3/B(3,4)$ on $(0,1)$ (Beta$(3,4)$); proposal $q(x,y)=1$ for $y\in(0,1)$ (independence, uniform).
\item $\alpha(x,y)$: probability of moving from $x$ to the proposal $y$.
\end{notacao}

\passo{Apply (7.6)}{Both $x,y\in(0,1)$, so $f(x)q(x,y)\neq0$. Substituting $f$ and $q\equiv1$ into (7.6):}
\[
\alpha(x,y)=\min\Big\{\frac{\big(y^2(1-y)^3/B(3,4)\big)\cdot(1)}{\big(x^2(1-x)^3/B(3,4)\big)\cdot(1)},1\Big\}.
\]
The constant $B(3,4)$ cancels:
\begin{equation}\label{q31:alpha}
\alpha(x,y)=\min\Big\{\frac{y^2(1-y)^3}{x^2(1-x)^3},\,1\Big\}.
\end{equation}

\passo{Why the algorithm can omit the min}{For $U_1\sim U(0,1)$ and any $r\ge0$,
$P(U_1\le\min(r,1))=\min(r,1)=P(U_1\le r)$. So step 7.3, which compares $U_1$ with the raw ratio,
accepts with exactly the probability \eqref{q31:alpha}.}

\passo{Detailed balance}{For $x\ne y$, by \eqref{q31:alpha},
$f(x)q(x,y)\alpha(x,y)=\min\{f(y),f(x)\}=f(y)q(y,x)\alpha(y,x)$; this symmetry is the reversibility
condition from which (7.6) was derived, so Beta$(3,4)$ is invariant (Question~33(b) proves the
general step).}

\codehead
{\small
\begin{verbatim}
fb = lambda v: v**2 * (1 - v)**3
rng = np.random.default_rng(31); x, acc, out = 0.5, 0, []
for g in range(5500):
    U1, U2 = rng.uniform(), rng.uniform()
    if U1 <= fb(U2) / fb(x):                  # (31.1), min not needed
        x = U2; acc += g >= 500
    if g >= 500: out.append(x)
acc / 5000, np.mean(out)                     # 0.5682  0.432394  (true 3/7 = 0.428571)
dblquad(lambda y, x: beta(3, 4).pdf(x) * min(1, fb(y) / fb(x)),
        0, 1, 0, 1)[0]                       # 0.572589
\end{verbatim}
}
The exact long-run acceptance probability $0.5726$ matches Greenberg's reported $0.57$ and our
$0.5682$.

\intuicao{With an independence proposal the chain compares how much more likely the target makes
the new point than the current one; it always moves uphill and sometimes downhill, and the
normalizing constant never matters.}
\refbib{Greenberg (2008), §7.2 (pp.~97--100).}
\end{sol}

% -----------------------------------------------------------------------------
\questao{G 7.6 --- accept--reject versus Metropolis--Hastings for a beta mean.}
\fonte{Greenberg, E. (2008), \textit{Introduction to Bayesian Econometrics}, 1st ed., Cambridge
University Press, Chapter~7, §7.5 Exercises, Exercise~7.6, p.~106 [PDF p.~121]; the book gives no
answer --- reproduced verbatim. Algorithm 5.2 and the condition on $c$, to which it refers, are
reproduced from §5.3, p.~66 [PDF p.~81].}

\emph{Referred to (Greenberg 2008, p.~66):} ``The accept--reject (AR) algorithm can be used to
simulate values from a density function $f(\cdot)$ if it is possible to simulate values from a
density $g(\cdot)$ and if a number $c$ can be found such that $f(Y)\le cg(Y)$, $c\ge1$, for all $Y$ in
the support of $f(\cdot)$. [\dots] \textbf{Algorithm 5.2: Accept--reject method}
\begin{itemize}[leftmargin=2.6em]
\item[5.1] Generate a value $y$ from $g(\cdot)$.
\item[5.2] Draw a value $u$ from $U(0,1)$.
\item[5.3] Return $y$ as a draw from $f(\cdot)$ if $u\le f(y)/cg(y)$. If not, reject it and return to step 1.
(The effect of this step is to accept $y$ with probability $f(y)/cg(y)$.)''
\end{itemize}

7.6 Estimate the mean of a Beta(3.7, 4.8) distribution with (1) an AR algorithm and a Beta(4, 5)
proposal density (you will need to determine the value of $c$ needed in Algorithm 5.2); (2) an MH
algorithm with a Beta(4, 5) proposal density. After the break-in sample, graph the values of the
mean against the iteration number to monitor convergence. Compare your answers to the true value.
\programacao With \texttt{np.random.default\_rng(32)}: (1) evaluate $f/g$ near $0$ and at $0.5$ for
the Beta$(4,5)$ proposal; (2) run AR with a proposal for which a finite $c$ exists (justify it) until
$10{,}000$ draws are accepted; (3) run the independence MH with Beta$(4,5)$ (1{,}000 burn-in,
10{,}000 kept), report the acceptance rate, the running mean at 1{,}000, 2{,}500, 5{,}000 and 10{,}000
kept iterations, and the batch-means n.s.e.; plot the running mean.
\resolespaco
\begin{sol}
\solhead
\begin{notacao}
\item $f$: Beta$(3.7,4.8)$ density, mean $3.7/(3.7+4.8)=0.435294$; $g_{a,b}$: Beta$(a,b)$ density;
  $B(\cdot,\cdot)$: beta function.
\item $c$: AR envelope constant; $w(x)=f(x)/g_{4,5}(x)$: MH importance ratio.
\end{notacao}

\passo{(1) No finite $c$ for Beta(4,5)}{Dividing the densities,}
\begin{equation}\label{q32:ratio45}
\frac{f(x)}{g_{4,5}(x)}=\frac{x^{2.7}(1-x)^{3.8}/B(3.7,4.8)}{x^{3}(1-x)^{4}/B(4,5)}=\frac{B(4,5)}{B(3.7,4.8)}\,x^{-0.3}(1-x)^{-0.2},
\end{equation}
which tends to $\infty$ as $x\to0$ or $x\to1$. No $c$ satisfies $f\le cg_{4,5}$ on $(0,1)$: \emph{the AR
algorithm the exercise asks for does not exist.} The proposal has thinner tails at both ends than
the target ($x^3$ versus $x^{2.7}$ near 0). Numerically the ratio is $0.9528$ at $x=0.5$ but $169.2$
at $x=10^{-8}$.

\passo{(1) A valid AR proposal (added)}{A Beta$(a,b)$ proposal works iff $a\le3.7$ and $b\le4.8$. With
Beta$(3,4)$, $f/g_{3,4}\propto x^{0.7}(1-x)^{0.8}$; its log-derivative $0.7/x-0.8/(1-x)$ vanishes at}
\begin{equation}\label{q32:xstar}
x^*=\frac{0.7}{0.7+0.8}=0.466667,\qquad c=\frac{f(x^*)}{g_{3,4}(x^*)}=1.115321,\qquad \frac1c=0.896603 .
\end{equation}

\passo{(2) MH with Beta(4,5)}{An independence kernel only needs $g_{4,5}>0$ wherever $f>0$, which
holds. By (7.6) (Question~31) with $q(x,y)=g_{4,5}(y)$,}
\begin{equation}\label{q32:mh}
\alpha(x,y)=\min\Big\{\frac{f(y)\,g_{4,5}(x)}{f(x)\,g_{4,5}(y)},1\Big\}=\min\Big\{\frac{w(y)}{w(x)},1\Big\}.
\end{equation}
Because $w$ is unbounded the chain is not uniformly ergodic: after landing near $0$ or $1$ (large
$w$) it can stay put for a while. It still converges to $f$, but convergence should be monitored,
which is the point of the running-mean graph.

\passo{Results (programming)}{AR with Beta$(3,4)$: acceptance $0.894775$ (theory $0.896603$), mean
$0.434564$ (s.e.\ $0.001621$). MH with Beta$(4,5)$: acceptance $0.9645$; running mean after
1{,}000, 2{,}500, 5{,}000 and 10{,}000 kept iterations: $0.4284$, $0.4333$, $0.4375$, $0.437056$;
n.s.e.\ $0.001744$. Both are within about one standard error of $0.435294$.}

\codehead
{\small
\begin{verbatim}
f, g34, g45 = beta(3.7, 4.8), beta(3, 4), beta(4, 5)
f.pdf(1e-8) / g45.pdf(1e-8), f.pdf(.5) / g45.pdf(.5)     # 169.2  0.9528
c = f.pdf(0.7 / 1.5) / g34.pdf(0.7 / 1.5)                # 1.115321
rng = np.random.default_rng(32); dr, tries = [], 0
while len(dr) < 10_000:
    y, u = rng.beta(3, 4), rng.uniform(); tries += 1
    if u <= f.pdf(y) / (c * g34.pdf(y)): dr.append(y)
np.mean(dr), 10_000 / tries                              # 0.434564  0.894775
w = lambda v: f.pdf(v) / g45.pdf(v); x, acc, out = .5, 0, []
for g in range(11_000):
    y = rng.beta(4, 5)
    if rng.uniform() <= min(1, w(y) / w(x)): x = y; acc += g >= 1000
    if g >= 1000: out.append(x)
run = np.cumsum(out) / np.arange(1, 10_001)              # plot run against iteration
acc / 10_000, run[[999, 2499, 4999, 9999]]   # 0.9645 [0.4284 0.4333 0.4375 0.4371]
\end{verbatim}
}

\intuicao{Accept--reject needs the proposal to dominate the target everywhere, so it must have
heavier tails; MH only needs the proposal to reach everywhere, and pays for poor tails with sticky
stretches instead of impossibility. Beta$(4,5)$ is a near-perfect body match (MH acceptance 96\%)
with tails that are slightly too thin.}
\refbib{Greenberg (2008), §5.3 (AR, Algorithm 5.2), §7.2 (independence kernels), §7.3 (convergence
monitoring).}
\end{sol}

% -----------------------------------------------------------------------------
\questao{Original math question --- random-walk MH on the log scale for a Poisson rate.}
\fonte{Original question (not from a book).}

The coal posterior of Question~14 is $\pi(\theta\mid y)\propto\theta^{a-1}e^{-b\theta}$ with $a=127.5$,
$b=42$. Pretend it is not recognized as a gamma and sample it by Metropolis--Hastings with the
multiplicative random walk $\theta'=\theta e^{u}$, $u\sim N(0,s^2)$.
\begin{enumerate}[label=\textbf{(\alph*)},leftmargin=2.4em]
\item Derive the proposal density $q(\theta,\theta')$ and the acceptance probability. Show that the
same probability arises from a symmetric random walk on $\varphi=\log\theta$ with the target written
in $\varphi$.
\item Prove that the MH kernel (with its rejection mass) leaves $\pi$ invariant, via detailed balance.
\item Suppose the factor $\theta'/\theta$ is forgotten. Find the distribution the chain then samples and
the bias in the posterior mean.
\item Parameter analysis: prove the limits of the acceptance probability as $s\to0$ and $s\to\infty$
and explain why the numerical standard error is large at both extremes.
\end{enumerate}
\programacao Implement the sampler (start $\theta=1$, 2{,}000 burn-in, 50{,}000 kept,
\texttt{np.random.default\_rng(33)} for each run) for $s\in\{0.02,0.2,1\}$, reporting acceptance
rate, posterior mean and batch-means n.s.e.\ (50 batches) against the exact mean $a/b$; repeat
$s=0.2$ without the Jacobian factor.
\resolespaco
\begin{sol}
\solhead
\begin{notacao}
\item $\pi(\theta)\equiv\pi(\theta\mid y)\propto\theta^{a-1}e^{-b\theta}$, $\theta>0$; $q(\theta,\theta')$: proposal density.
\item $K(\theta,A)$: MH transition probability into a set $A$; $r(\theta)$: probability of staying put.
\item $\varphi=\log\theta$; $\tilde\pi(\varphi)$: density of $\varphi$; $s>0$: step size.
\end{notacao}

\passo{(a) Proposal density}{$\log\theta'=\log\theta+u\sim N(\log\theta,s^2)$, so $\theta'$ is log-normal:}
\begin{equation}\label{q33:q}
q(\theta,\theta')=\frac{1}{s\,\theta'}\,\phi\Big(\frac{\log(\theta'/\theta)}{s}\Big),\qquad \theta'>0 .
\end{equation}
Since $\phi$ is even, $\phi(\log(\theta/\theta')/s)=\phi(\log(\theta'/\theta)/s)$, and dividing
\eqref{q33:q} with the arguments swapped by \eqref{q33:q}:
\begin{equation}\label{q33:qr}
\frac{q(\theta',\theta)}{q(\theta,\theta')}=\frac{1/(s\theta)}{1/(s\theta')}=\frac{\theta'}{\theta}.
\end{equation}

\passo{(a) Acceptance}{Substituting \eqref{q33:qr} into (7.6) of Question~31:}
\begin{equation}\label{q33:alpha}
\alpha(\theta,\theta')=\min\Big\{\frac{\pi(\theta')}{\pi(\theta)}\cdot\frac{\theta'}{\theta},1\Big\}
=\min\Big\{\Big(\frac{\theta'}{\theta}\Big)^{a}e^{-b(\theta'-\theta)},1\Big\}.
\end{equation}

\passo{(a) Log scale}{With $\theta=e^\varphi$, $d\theta=e^\varphi d\varphi$, the density of $\varphi$ is
$\tilde\pi(\varphi)=\pi(e^\varphi)e^{\varphi}$. A symmetric walk $\varphi'=\varphi+u$ accepts with
$\min\{\tilde\pi(\varphi')/\tilde\pi(\varphi),1\}=\min\{\pi(\theta')\theta'/(\pi(\theta)\theta),1\}$, the
same as \eqref{q33:alpha}: the factor $\theta'/\theta$ is the Jacobian.}

\passo{(b) Detailed balance}{For $\theta\ne\theta'$, by \eqref{q33:alpha} in general form,}
\begin{equation}\label{q33:db}
\pi(\theta)q(\theta,\theta')\alpha(\theta,\theta')=\min\{\pi(\theta)q(\theta,\theta'),\ \pi(\theta')q(\theta',\theta)\}
=\pi(\theta')q(\theta',\theta)\alpha(\theta',\theta).
\end{equation}
The kernel is $K(\theta,A)=\int_Aq(\theta,\theta')\alpha(\theta,\theta')d\theta'+r(\theta)\1\{\theta\in A\}$,
with $r(\theta)=1-\int q(\theta,\theta')\alpha(\theta,\theta')d\theta'$. Integrating against $\pi$ and
using \eqref{q33:db} to swap the roles of $\theta$ and $\theta'$ in the first term:
\[
\int\pi(\theta)K(\theta,A)d\theta=\int_A\pi(\theta')\Big[\int q(\theta',\theta)\alpha(\theta',\theta)d\theta\Big]d\theta'+\int_A\pi(\theta)r(\theta)d\theta .
\]
The bracket is $1-r(\theta')$, so
\[
\int\pi(\theta)K(\theta,A)d\theta=\int_A\pi(\theta')\big(1-r(\theta')\big)d\theta'+\int_A\pi(\theta)r(\theta)d\theta=\int_A\pi=\pi(A). \qquad\blacksquare
\]

\passo{(c) Forgetting the Jacobian}{Without $\theta'/\theta$ the acceptance is
$\min\{\pi(\theta')/\pi(\theta),1\}$ while $q$ is not symmetric. Writing it as a correct MH acceptance
for some target $\pi^\dagger$ requires $\pi^\dagger(\theta')\theta'/(\pi^\dagger(\theta)\theta)=\pi(\theta')/\pi(\theta)$,
i.e.\ $\pi^\dagger(\theta)\propto\pi(\theta)/\theta$. By (b) the chain is then invariant for}
\begin{equation}\label{q33:bug}
\pi^\dagger(\theta)\propto\theta^{a-2}e^{-b\theta}=G(a-1,b),\qquad \E^\dagger(\theta)-\E(\theta)=\frac{a-1}{b}-\frac ab=-\frac1b .
\end{equation}
For the coal posterior: $126.5/42=3.011905$ instead of $127.5/42=3.035714$.

\passo{(d) Limits}{Write $R=(\theta'/\theta)^ae^{-b(\theta'-\theta)}=e^{au-b\theta(e^u-1)}$ with $u=s Z$,
$Z\sim N(0,1)$. As $s\to0$, $u\to0$ and $R\to1$ for every $Z$; since $\min\{R,1\}\le1$, dominated
convergence gives $E[\alpha]\to1$. As $s\to\infty$, for $Z>0$, $u\to\infty$ and
$au-b\theta(e^u-1)\to-\infty$; for $Z<0$, $u\to-\infty$ and $au\to-\infty$ while $b\theta(1-e^{u})\le b\theta$;
so $R\to0$ for every $Z\ne0$ and $E[\alpha]\to0$.}

\passo{(d) Reading}{Tiny steps are almost always accepted but barely move, so successive draws are
highly correlated; huge steps are almost always rejected and the chain sits still. Both inflate the
n.s.e.; the best $s$ is of the order of the posterior standard deviation of $\log\theta$, here
$1/\sqrt{a}\approx0.089$.}

\codehead
{\small
\begin{verbatim}
a, b = 127.5, 42.; lp = lambda t: (a - 1) * np.log(t) - b * t
def rw_log(s, jac, seed, G=50_000, burn=2_000):
    rng = np.random.default_rng(seed); th, acc, out = 1., 0, np.empty(G)
    for g in range(G + burn):
        pr = th * np.exp(s * rng.standard_normal())
        la = lp(pr) - lp(th) + (np.log(pr / th) if jac else 0.)   # (33.3)
        if np.log(rng.uniform()) <= la: th = pr; acc += g >= burn
        if g >= burn: out[g - burn] = th
    return acc / G, out.mean(), out.reshape(50, -1).mean(1).std(ddof=1) / np.sqrt(50)
# s     accept   mean      n.s.e.        (exact mean 3.035714)
# 0.02  0.9284   3.037804  0.011093
# 0.2   0.4636   3.037605  0.001913
# 1     0.1127   3.032593  0.005092
# 0.2, no Jacobian: 0.4664  3.015205  0.002112   (target of (33.5): 3.011905)
\end{verbatim}
}
The middle step size has the smallest n.s.e., about six times smaller than $s=0.02$; the
Jacobian-free chain settles at $3.0152$, $1.6$ n.s.e.\ from its target $3.0119$ in \eqref{q33:bug}
and $9.7$ n.s.e.\ below the true mean.

\intuicao{A proposal that moves multiplicatively is not symmetric in $\theta$, and MH must be told:
the Jacobian is the price of proposing on one scale and targeting on another. Forgetting it does not
make the sampler crash; it silently samples the wrong distribution.}
\refbib{Greenberg (2008), §7.2 (MH algorithm, random-walk kernels, pp.~97--100); Tsay (2010),
§12.3 (Metropolis--Hastings).}
\rubrica{
\item (a) \eqref{q33:q}--\eqref{q33:alpha}: 4 pts; log-scale equivalence: 1 pt.
\item (b) \eqref{q33:db} and the integration with the rejection mass: 4 pts (detailed balance stated without the rejection term: max.\ 2).
\item (c) \eqref{q33:bug}: 3 pts.
\item (d) both limits proved and the reading: 3 pts.
\item Programming: three runs and the Jacobian-free run with correct comparisons: 3 pts.
}
\end{sol}

% -----------------------------------------------------------------------------
\questao{Original conceptual question --- MCMC (10 points).}
\fonte{Original question (not from a book).}

\textbf{a)} (3 points) In MH with a symmetric random-walk proposal, $q(x,y)=q(y,x)$, the probability
of accepting a move from $x$ to $y$ is:
\begin{alternativas}
\item $\min\{f(y)/f(x),1\}$, the ratio of target densities capped at one.
\item $\min\{f(x)/f(y),1\}$, the ratio of current to proposed density.
\item $f(y)/\{f(x)+f(y)\}$, the share of the proposed point in the pair.
\item $1$, because a symmetric proposal needs no correction.
\end{alternativas}
\gabarito{A}
\begin{sol}
\passo{Why}{A: (7.6) with the $q$'s cancelling (Greenberg p.~99). B inverts the ratio, so it would
favour moves downhill; C is Barker's rule, a different (valid but less efficient) acceptance; D
ignores the target altogether.}
\end{sol}

\textbf{b)} (3 points) The Gibbs sampler is a special case of MH in which:
\begin{alternativas}
\item the proposal is a symmetric random walk on each block in turn.
\item every proposal is accepted with probability $1/2$.
\item the draws are independent, as in accept--reject.
\item each block is proposed from its full conditional, and the acceptance probability is $1$.
\end{alternativas}
\gabarito{D}
\begin{sol}
\passo{Why}{D: with $q$ equal to the full conditional, the ratio in (7.6) is exactly 1 (Greenberg
§7.2). A and B describe other kernels; C is false: Gibbs draws are dependent (Question~11 found an
inefficiency factor of 1.17).}
\end{sol}

\textbf{c)} (2 points, T or F) If a kernel satisfies detailed balance with respect to $\pi$, then
$\pi$ is an invariant distribution of the kernel.
\gabarito{T}
\begin{sol}
\passo{Why}{Question~33(b): integrate detailed balance over the starting point.}
\end{sol}

\textbf{d)} (2 points, T or F) Detailed balance is necessary for $\pi$ to be invariant: every
MCMC kernel with invariant distribution $\pi$ is reversible.
\gabarito{F}
\begin{sol}
\passo{Why}{The converse of (c) fails. The systematic-scan Gibbs sampler of Question~11 (update
$\theta_1$, then $\theta_2$, then $k$) leaves the posterior invariant, since each step does, but the
composed kernel is not reversible in general. G~6.8's kernel was checked for invariance directly,
not through detailed balance.}
\end{sol}

\programacao Check (d) exactly on a toy target: for $p(x,y)$ on $\{0,1\}^2$ with
$p(0,0)=0.3$, $p(0,1)=0.2$, $p(1,0)=0.1$, $p(1,1)=0.4$, build the $4\times4$ transition matrices
$P_x$ (update $x$ from $p(x\mid y)$) and $P_y$ (update $y$ from $p(y\mid x)$), the systematic scan
$P=P_xP_y$, and report $\max|\pi P-\pi|$ and $\max_{i,j}|\pi_iP_{ij}-\pi_jP_{ji}|$.
\begin{sol}
\codehead
{\small
\begin{verbatim}
p = np.array([[.3, .2], [.1, .4]]); st = [(0, 0), (0, 1), (1, 0), (1, 1)]
pi = np.array([p[s] for s in st])
Px = np.array([[p[a, d] / p[:, d].sum() if b == d else 0 for (a, b) in st]
               for (c, d) in st])                 # row: current (c, d)
Py = np.array([[p[c, b] / p[c, :].sum() if a == c else 0 for (a, b) in st]
               for (c, d) in st])
P = Px @ Py; F = pi[:, None] * P
np.abs(pi @ P - pi).max(), np.abs(F - F.T).max()      # 5.6e-17   0.053333
\end{verbatim}
}
$\pi$ is invariant to machine precision, yet the flow $\pi_iP_{ij}$ is not symmetric (largest gap
$0.053333$): invariance without reversibility. Each single update $P_x$, $P_y$ is reversible; their
product is not.
\end{sol}

\end{document}
